%% Metatime Monograph
%% Publication-candidate presentation of an exploratory low-parameter ansatz
%% with explicit claim classes, data conventions, failures, and prospective gates
%%
%% Adrian Lipa, Doncaster, UK
%% Version 11.0 - Publication Candidate

\documentclass[11pt,a4paper,twoside]{book}
\usepackage[utf8]{inputenc}
\usepackage[T1]{fontenc}
\usepackage{lmodern}
\usepackage{amsmath,amssymb,amsthm,bm,mathtools}
\usepackage{slashed}
\usepackage{graphicx}
\usepackage{xurl}
\usepackage[protrusion=true,expansion=false]{microtype}
\usepackage[hidelinks,pdfpagelayout=TwoPageRight,pdftitle={Metatime: Publication Candidate of a Low-Parameter Phenomenological Ansatz},pdfauthor={Adrian Lipa},pdfsubject={Geometric-phase information structures and particle, flavor, hadronic, and cosmological observables},pdfkeywords={Metatime, TIR, geometric phase, information geometry, Standard Model phenomenology, Collatz, Ramanujan, falsifiability}]{hyperref}
\usepackage{makeidx}
\usepackage{verbatim}
\usepackage{longtable}
\usepackage{booktabs}
\usepackage[left=2.5cm,right=2.5cm,top=3cm,bottom=3cm,headheight=14pt]{geometry}
\setlength{\emergencystretch}{3em}
\makeindex

\theoremstyle{definition}
\newtheorem{definition}{Definition}[chapter]
\newtheorem{theorem}{Theorem}[chapter]
\newtheorem{lemma}{Lemma}[chapter]
\newcommand{\claimstatus}[1]{\par\noindent\textbf{Claim status: #1}\par}
\newcommand{\datasetstatus}[1]{\par\noindent\textbf{Data status: #1}\par}
\newcommand{\OI}{\kappa}
\newcommand{\Lam}{\Lambda}
\newcommand{\dd}{\,\mathrm{d}}
\newcommand{\mev}{\,\text{MeV}}
\newcommand{\gev}{\,\text{GeV}}
\newcommand{\eV}{\,\text{eV}}
\newcommand{\PDG}[1]{\ensuremath{#1^{\text{PDG}}}}
\newcommand{\pdgsup}{\ensuremath{^{\text{PDG}}}}
\newcommand{\Eproton}{E_{\text{proton}}}
\newcommand{\Eplanck}{E_{P}}
\newcommand{\keV}{\,\text{keV}}
\newcommand{\tev}{\,\text{TeV}}
\newcommand{\SM}{\text{SM}}

\begin{document}
\raggedbottom
\frontmatter
\title{Metatime: A Publication Candidate for an Exploratory\\Low-Parameter Phenomenological Ansatz}
\author{Adrian Lipa}
\date{Independent Researcher, Doncaster, UK\\Version 11.0 -- 29 July 2026}
\maketitle
\chapter*{Abstract}
\addcontentsline{toc}{chapter}{Abstract}

This monograph presents Metatime/TIR as an exploratory, low-parameter
phenomenological ansatz relating geometric-phase and information-theoretic
structures to particle, flavor, hadronic, and cosmological observables.  The
construction uses the defined information constant
$\kappa=\ln 2/(24\pi)$, three discrete integers $(L_3,L_4,L_5)=(7,2,5)$,
prime-valued flavor labels, and a small set of external physical anchors.
The manuscript does not claim that these ingredients have been derived from
an accepted microscopic theory.  Instead, it distinguishes established
mathematics, model postulates, retrospective formula assignments, diagnostic
no-go results, and prospectively frozen tests.

The retrospective ledger contains close numerical matches in several sectors,
but these matches are not combined into a single global accuracy score because
the observables have different units, uncertainties, dependence structures,
and degrees of post-selection.  Two failures are kept explicit: the quoted
gauge-boson masses remain approximately $4.5$--$5.0\%$ high, and the fixed
strong-CP-to-neutron-EDM mapping gives
$d_n=5.33\times10^{-26}\,e\,\mathrm{cm}$, a factor $2.96$ above the current
$90\%$ confidence upper bound used in this work.  The Collatz $3/4$ mass trace
is also retained as an incomplete comparative signal rather than a closed mass
derivation.  A restricted common-baseline repair is excluded by an explicit
no-go theorem.  Three separable candidates are frozen for future direct
Higgs-coupling likelihoods, with no candidate deletion or formula alteration
permitted after inspection of the named data.

The complete source, validation ledgers, and reproducible LaTeX build are
available in the accompanying repository.

\section*{Keywords}
geometric phase; information geometry; phenomenological ansatz; flavor;
Standard Model observables; Collatz dynamics; Ramanujan structure;
preregistration; reproducibility.

\section*{Claim and evidence taxonomy}
\begin{description}
  \item[A -- established result:] a standard mathematical identity or external
  experimental result supported by the cited literature.
  \item[B -- model postulate:] a defining structural choice of Metatime/TIR.
  \item[C -- retrospective assignment:] a formula assessed after some or all of
  the comparison data were known.
  \item[D -- diagnostic or no-go:] an internal audit, failure, or restricted
  impossibility result.
  \item[E -- prospective prediction:] a frozen formula, observable, and decision
  rule specified before the named future dataset is inspected.
  \item[F -- external anchor:] a measured scale or coefficient supplied to the
  construction rather than predicted by it.
\end{description}

\section*{Availability and version status}
This is Version 11.0, designated a \emph{Publication Candidate}.  It is suitable
for circulation as a research monograph or preprint.  It is not presented as an
experimentally established derivation of the Standard Model from first
principles.  Source code, version history, machine-readable validation outputs,
and the build workflow are maintained at
\url{https://github.com/AdrianLipa90/The-Fundamental-Theory-of-Informational-Relations}.

\section*{Declarations}
This manuscript is an independent research work.  No external grant is
identified in the repository record.  No competing interest is declared in
this manuscript version.

\cleardoublepage

\tableofcontents
\mainmatter

\part{Foundations}
% !TEX root = ../metatime_monograph.tex
\chapter{Introduction}
\label{ch:introduction}

\section{The Problem of Parameters}
The Standard Model is one of the most precisely tested theories in physics, but
many of its numerical parameters are empirical inputs.  The Metatime/TIR
programme investigates whether a restricted set of information-geometric and
discrete structures can organize or approximate part of that parameter
spectrum.

\section{Framework and fixed structural quantities}
The present construction uses
\[
\kappa=\frac{\ln 2}{24\pi},\qquad
L_3=7,\quad L_4=2,\quad L_5=5,
\]
together with the quark-prime labels
\[
(u,d,s,c,b,t)=(3,5,7,11,13,17).
\]
The framework also contains explicit external scales and numerical inputs,
including the Planck scale and proton mass.  These inputs must be counted
separately from internal structural choices.

\section{Claim hierarchy}
The monograph distinguishes:
\[
\begin{array}{ll}
\text{[A]} & \text{mathematical identities},\\
\text{[B]} & \text{framework postulates},\\
\text{[C]} & \text{phenomenological assignments},\\
\text{[D]} & \text{numerical coincidences requiring justification},\\
\text{[E]} & \text{prospectively falsifiable predictions}.
\end{array}
\]
A numerical agreement is not automatically a prediction.  In particular,
postdictions, retrospectively revised formulae, and quantities whose derivation
is incomplete are recorded separately.

\section{Mass-action ansatz}
Several mass sectors use candidate relations of the form
\[
m=E_P e^{-S/\kappa}.
\]
This relation is not asserted to be fully established for every particle
sector.  Each sector must independently demonstrate a dimensionally and
arithmetically complete action $S$, together with provenance showing whether
the formula was frozen before the comparison data were inspected.

\section{PDG 2026 status}
Chapter~\ref{ch:pdg2026_addendum} records the 15 August 2026 validation freeze.
At that freeze:
\begin{itemize}
\item CKM is the strongest numerical compatibility sector, but remains
      postdictive rather than independently validated;
\item selected neutrino observables are compatible, while $\theta_{13}$ is a
      current stress point;
\item the charged-lepton action arithmetic is internally reproducible but does
      not provide precision-level agreement;
\item the baryon decuplet remains a useful hadronic test candidate, while the
      octet carries retrospective-refinement provenance;
\item the printed pion and kaon mass equations are internally inconsistent and
      are quarantined pending derivation repair;
\item the present electroweak and Higgs formulae fail current precision-level
      tests and require a scheme/running/radiative treatment;
\item the quark-mass table is not a precision prediction until a
      scheme- and scale-defined prime-to-mass map is supplied.
\end{itemize}

\section{Prospective-evidence rule}
Beginning with the PDG 2026 addendum, a claim intended as prospective evidence
must publish its formula/version identifier, source hash, external inputs,
construction dataset, and frozen numerical output before the comparison data
are opened.  FAIL and TENSION outcomes remain in the validation ledger and do
not authorize silent formula replacement.

% !TEX root = ../metatime_monograph.tex
\chapter{The Metatime Framework}
\label{ch:metatime_framework}

% CITATION_CONTEXT_V10_9_BEGIN
\paragraph{Established context.}
The standard geometric ingredients behind phase holonomy, quantum-state metrics,
and cyclic evolution are described by Pancharatnam, Simon, Berry,
Aharonov--Anandan, and Provost--Vall\'ee
\cite{Pancharatnam1956,Simon1983,berry1984,AharonovAnandan1987,ProvostVallee1980}.
Connections to information geometry are conventional in the sense of Fisher,
Rao, and Amari--Nagaoka \cite{Fisher1925,Rao1945,AmariNagaoka2000}; the specific
Metatime synthesis and its operator assignments are the construction tested in
this monograph.
% CITATION_CONTEXT_V10_9_END

\section{The Information Constant \(\kappa\)}

The Metatime framework defines the information constant
\begin{equation}
\kappa\equiv\frac{\ln2}{24\pi}.
\label{eq:kappa_construction}
\end{equation}
Its numerical value is
\[
\kappa=0.009193150006360484\ldots .
\]

\claimstatus{B -- model postulate with structural motivation; the arithmetic identities below are exact, while the physical assignments of their factors are model-level.}

This claim status is the one used consistently in Appendix~\ref{app:kappa}.  The
dedicated information-spinor analysis under
\path{TIR/zeta_information_axis/}
and the crosswalk in Appendix~\ref{app:information_spinor_crosswalk} sharpen the
internal factorization of the normalization without promoting it to an
established first-principles invariant of quantum field theory or differential
geometry.

\subsection{Binary information}

For a balanced binary distribution, Shannon entropy in natural units is
\cite{shannon1948}
\begin{equation}
H_{\rm bin}\!\left(\frac12\right)=\ln2.
\label{eq:shannon_binary}
\end{equation}
This is an established information-theoretic identity.  It supplies the
numerator of \eqref{eq:kappa_construction}; it does not by itself determine the
denominator.

\subsection{Normalized recurrence before radians}

A cyclic process may first be represented by a normalized turn coordinate
\begin{equation}
q\in\mathbb R/\mathbb Z,
\end{equation}
so that one recurrence is $q\mapsto q+1$.  The intrinsic frequency is a winding
rate, $f_q=\Delta N/\Delta t$, before an angular convention is selected.  An
angular coordinate is introduced only by
\begin{equation}
\phi=Cq,
\end{equation}
where $C$ is the closure period in the chosen angular unit; in radians
$C=2\pi$.

For the two-level state used in the information-spinor subprogramme, the
balanced equatorial loop carries Berry holonomy $-1$.  This is a standard
geometric-phase statement for the specified state family
\cite{Simon1983,berry1984}.  It should not be replaced by the looser claim that
an arbitrary ``complete cycle'' has solid angle $4\pi$: Berry phase depends on
the actual closed path and its enclosed solid angle.

\subsection{Discrete cross-factorization}

The exact arithmetic identity
\begin{equation}
24=8\cdot3=12\cdot2=6\cdot4
\label{eq:twentyfour_crossfactor}
\end{equation}
provides several internally compatible decompositions.  Their Metatime/TIR
interpretations are model assignments:
\begin{itemize}
\item $8\cdot3$: eight declared mixing sectors times three flavours;
\item $12\cdot2$: twelve projective cycles times two half-turn units;
\item $6\cdot4$: six spinor cycles times four half-turn units.
\end{itemize}
The equality of the integers is exact; the physical interpretation of the
factors is not an independent theorem.  The older tetrahedral notation remains
a useful structural representation: $|A_4|=12$ for rotational tetrahedral
symmetry and $|S_4|=24$ for the full tetrahedral symmetry.  The use of these
groups is part of the NOEMA/Metatime model geometry and must not be conflated
with the Weyl group of $SU(3)$, which is $S_3$.

\subsection{Conditional arithmetic--geometry reconstruction of \(\kappa\)}

If the Metatime information cycle assigns one binary information quantum to
twelve normalized projective cycles,
\begin{equation}
\frac{\dd I}{\dd q}=\frac{\ln2}{12},
\label{eq:info-per-projective-turn}
\end{equation}
then conversion to an angular coordinate with closure period $C$ gives exactly
\begin{equation}
\frac{\dd I}{\dd\phi}=\frac{\ln2}{12C}.
\end{equation}
For the radian representation $C=2\pi$,
\begin{equation}
\frac{\dd I}{\dd\phi}=\frac{\ln2}{24\pi}=\kappa.
\end{equation}
Thus the present formulation separates the model's arithmetic cycle assignment
from geometric closure.  It is a useful cross-derivation and implementation
constraint, but it does not remove the model status of the twelve-cycle
assignment itself.  Appendix~\ref{app:kappa} remains the authoritative
publication claim boundary.

\section{The Exponential Mass Formula}

The central postulate of the Metatime framework is that every physical mass
is determined by an action \(S\) on the Poincar\'{e} disk via an exponential
Boltzmann-like factor:

\begin{equation}
m = E_P \, e^{-S/\kappa}.
\label{eq:mass_formula}
\end{equation}

Here \(E_P = \sqrt{\hbar c^5 / G} = 1.22089 \times 10^{22}\)~MeV is the Planck
energy (the fundamental scale of quantum gravity), and \(S\) is a
dimensionless action constructed from the L-constants and quark primes.

The exponential form is a Metatime modelling ansatz motivated by geometric
phase accumulation and information-weighted state counting.  It is not derived
from the standard canonical partition function by a theorem identifying
$\kappa$ with a thermodynamic temperature.

\subsection{Why the Exponential Form?}

For comparison, a canonical statistical system uses
\[
Z = \sum_i e^{-E_i / k_B T}.
\]
The Metatime construction adopts an analogous exponential architecture with
$S/\kappa$ dimensionless.  The analogy motivates the functional form; it does
not identify $S$ with ordinary thermodynamic entropy or $\kappa$ with $k_BT$.
Physical relevance is therefore assessed by the explicit numerical and
prospective tests in later chapters.

\section{Action Construction Principles}

Each particle's action \(S\) is built from the following elements:

\begin{enumerate}
\item \textbf{Spin degeneracy:} A term \(\frac{1}{2}\) for spin-1/2 fermions,
encoding the two spin states.
\item \textbf{Generation terms:} Terms proportional to \(\kappa\) for each
generation, representing the information deficit shared across generations.
\item \textbf{Flavour terms:} Terms involving quark primes \(q_i\) and
L-constants \(L_3, L_4, L_5\), encoding the flavour structure.
\item \textbf{Geometric corrections:} Terms involving ratios of L-constants
and powers of \(\kappa\), representing higher-order curvature corrections on
the Poincar\'{e} disk.
\end{enumerate}

\section{The NOEMA Tetrahedron}

The NOEMA (Neutral Operator Eigenvalue Mapping Algebra) tetrahedron is a
geometric construction in the three-dimensional operator space spanned by:

\[
\{\hat I,\ \hat M,\ \hat A\},
\]

where \(\hat I\) is the isospin operator, \(\hat M\) is the mass operator,
and \(\hat A\) is the axial operator.  The tetrahedron encodes the mixing
patterns of quarks and leptons:

\begin{itemize}
\item The edge lengths are determined by \(L_3\) and \(L_4\).
\item The face areas determine mass splittings within the model.
\item The dihedral angles supply candidate mixing-angle coordinates.
\end{itemize}

The NOEMA tetrahedron provides the model geometry used for the CKM and PMNS
constructions discussed in Chapters~\ref{ch:ckm_matrix} and
\ref{ch:pmns_mixing}.  Those phenomenological assignments are distinct from
the exact group orders appearing in \eqref{eq:twentyfour_crossfactor}.

\section{External Scales and Inputs}

The framework requires two external energy scales and three external numerical
inputs:

\begin{description}
\item[\(E_P = 1.22089 \times 10^{22}\)~MeV:] The Planck energy, setting the
absolute mass scale via the exponential formula.
\item[\(E_{\text{proton}} = 938.272\)~MeV:] The proton mass, used as the
reference scale for baryonic and electroweak quantities.
\item[\(N_c = 3\):] The number of colours in QCD, needed for hypercharge
assignments and anomaly cancellation.
\item[\(Y(L) = -1/2\):] The lepton doublet hypercharge, required for anomaly
cancellation checks.
\item[\(\text{Crewther factor} = 2.4 \times 10^{-16}\):] The chiral
perturbation theory factor relating \(\theta_{\text{QCD}}\) to the neutron
electric dipole moment.
\end{description}

The remaining numerical relations are generated from the declared structural
inputs of the model.  The phrase ``no continuous fit'' should not be read as
``no assumptions'': discrete assignments, external scales, and retrospective
operator choices are tracked separately by the publication claim hierarchy.

% !TEX root = ../metatime_monograph.tex
\chapter{The L-Constants}
\label{ch:l_constants}

Three geometric constants underpin the Metatime arithmetic:
\(L_3 = 7\), \(L_4 = 2\), and \(L_5 = 5\).  In this chapter we derive
each from the Collatz iteration on prime-indexed orbits and from the
twin-prime structure of the number 5.

\section{The Collatz Iteration}

The Collatz function (also known as the \(3n+1\) problem) is defined on
positive integers by:

\begin{equation}
f(n) = \begin{cases}
3n + 1 & \text{if } n \text{ is odd}, \\
n/2   & \text{if } n \text{ is even}.
\end{cases}
\label{eq:collatz}
\end{equation}

Starting from any positive integer, repeated application of \(f\) eventually
reaches the cycle \(4 \to 2 \to 1 \to 4\).  The Collatz depth of an integer
\(n\) is the number of steps required to reach 1.

\subsection{Collatz Depth of 3}

Starting from \(n = 3\):

\[
\begin{array}{c|c}
\text{Step} & \text{Value} \\ \hline
0 & 3 \text{ (odd)} \\
1 & 3\times 3 + 1 = 10 \text{ (even)} \\
2 & 10/2 = 5 \text{ (odd)} \\
3 & 3\times 5 + 1 = 16 \text{ (even)} \\
4 & 16/2 = 8 \text{ (even)} \\
5 & 8/2 = 4 \text{ (even)} \\
6 & 4/2 = 2 \text{ (even)} \\
7 & 2/2 = 1
\end{array}
\]

The trajectory is: \(3 \to 10 \to 5 \to 16 \to 8 \to 4 \to 2 \to 1\),
requiring 7 steps.  Hence:

\begin{equation}
L_3 \equiv \text{Collatz depth of 3} = 7.
\label{eq:L3}
\end{equation}

\section{The Twin-Prime Constants}

The twin-prime pair \((3,5)\) gives the remaining L-constants:

\begin{itemize}
\item The gap between the twin primes is \(5 - 3 = 2\), which defines the
annihilation constant:
\[
L_4 \equiv 2.
\]

\item The larger twin prime is 5, which defines the intention dimension:
\[
L_5 \equiv 5.
\]
\end{itemize}

Together these satisfy:

\[
L_4^{L_5} = 2^5 = 32,
\]

which is the dimension of the intention-phase space in the NOEMA tetrahedron
construction.

\section{Composite Constants}

The L-constants combine to form several composite quantities that appear
repeatedly in the sector formulas:

\[
\begin{array}{lcl}
L_3 + L_4 = 9, & L_3 - 1 = 6, & L_4/L_5 = 2/5 = 0.4, \\
L_3 + L_4 + L_5 = 14, & L_3^2 = 49, & L_4^2 = 4, \\
L_4^{L_5} = 32, & L_3 L_4^{L_5} = 224, & (L_4/L_3)^2 = 4/49.
\end{array}
\]

\section{Physical Interpretation}

The L-constants admit the following physical interpretations:

\begin{description}
\item[\(L_3 = 7\):] The structural depth of flavour space, corresponding to
the number of Collatz steps from 3 to 1.  Appears in denominators and as an
exponent in mixing angles (\(\sin\theta_{13} = 1/L_3\)).

\item[\(L_4 = 2\):] The annihilation constant, representing the twin-prime
gap.  Appears as the fundamental binary splitting factor.

\item[\(L_5 = 5\):] The intention dimension, the larger twin prime.
Appears in ratios with \(L_4\) to set mixing scales.
\end{description}

These constants together with \(\kappa\) and the quark primes generate every
numerical quantity in the Metatime framework.

% !TEX root = ../metatime_monograph.tex
\chapter{Quark Primes}
\label{ch:quark_primes}

Each quark flavour is assigned a prime number based on its position in
the generation structure:

\begin{equation}
(u,d,s,c,b,t) = (3,5,7,11,13,17).
\label{eq:quark_primes}
\end{equation}

These assignments are not arbitrary; they follow from an ordering principle
that encodes the flavour geometry.

\section{Assignment Principle}

The quark primes are the primes corresponding to the generation index:

\[
\begin{array}{c|c|c}
\text{Generation} & \text{Quarks} & \text{Prime} \\ \hline
1 & u, d & p_2 = 3,\ p_3 = 5 \\
2 & s, c & p_4 = 7,\ p_5 = 11 \\
3 & b, t & p_6 = 13,\ p_7 = 17
\end{array}
\]

Here \(p_n\) denotes the \(n\)-th prime number: \(p_1 = 2\), \(p_2 = 3\),
\(p_3 = 5\), \(p_4 = 7\), \(p_5 = 11\), \(p_6 = 13\), \(p_7 = 17\).
The assignments skip \(p_1 = 2\) (which is reserved for the binary structure
of the information constant) and start from \(p_2 = 3\).

\section{Flavour Geometry}

The prime numbers encode geometric relationships between flavours:

\begin{description}
\item[Isospin splitting:] The difference \(d - u = 5 - 3 = 2\) equals the
annihilation constant \(L_4\), encoding the up-down mass splitting.
\item[Strangeness scale:] The strange quark prime \(s = 7\) equals the
structural constant \(L_3\), establishing the strangeness scale.
\item[Generation spacing:] The prime gaps increase with generation:
\(|3-5| = 2\), \(|7-11| = 4\), \(|13-17| = 4\), reflecting the growing
hierarchy between successive generations.
\end{description}

\section{Sum and Difference Combinations}

Several combinations of quark primes appear in the sector formulas:

\[
\begin{array}{lcl}
s+u = 10, & s-u = 4, \\
s^2 - u^2 - d^2 = 49 - 9 - 25 = 15, & u+d+s = 15, \\
c = 11, & b = 13, & t = 17.
\end{array}
\]

\section{Physical Interpretation}

The quark primes serve as flavour labels on the Poincar\'{e} disk.  They
encode the action contributions of each flavour to the mass formulas.  The
metaphor is that of a periodic table: the primes are the atomic numbers of
flavour space, from which the properties of composite states (hadrons) are
derived by arithmetic combination rules.


\part{Charged Leptons}
% !TEX root = ../metatime_monograph.tex
\chapter{Electron Action}
\label{ch:electron_action}

The electron is the lightest charged lepton, with mass \(m_e =
0.510999\)~MeV.  Within the Metatime framework, this corresponds to the
largest action on the Poincar\'{e} disk:

\begin{equation}
S_e = \frac12 - 3\kappa + \frac{\kappa}{L_3} - \frac{\kappa^2}{2}.
\label{eq:Se_full}
\end{equation}

We now derive each term.

\section{Spin-1/2 Degeneracy}

The leading term \(1/2\) encodes the spin-1/2 degeneracy of the electron.
On the Poincar\'{e} disk, a spin-1/2 fermion has two accessible spin states,
contributing an additive constant to the action:

\begin{equation}
S_{\text{spin}} = \frac12.
\label{eq:spin_term}
\end{equation}

\section{Three-Generation Information Deficit}

The term \(-3\kappa\) represents the shared information deficit across the
three fermion generations.  Each generation contributes \(-\kappa\) to the
action, reflecting the information-theoretic cost of generational replication:

\begin{equation}
S_{\text{gen}} = -3\kappa.
\label{eq:gen_term}
\end{equation}

The negative sign indicates that more generations reduce the action, making
the corresponding particles heavier.

\section{Collatz Channel Correction}

The term \(+\kappa/L_3\) is a correction arising from the Collatz structure
of flavour space.  The factor \(1/L_3 = 1/7\) reflects the structural depth
of the first generation:

\begin{equation}
S_{\text{collatz}} = \frac{\kappa}{L_3}.
\label{eq:collatz_term}
\end{equation}

\section{Second-Order Curvature Correction}

The term \(-\kappa^2/2\) is the leading curvature correction on the
Poincar\'{e} disk.  It is quadratic in \(\kappa\) and represents the
information curvature of the disk:

\begin{equation}
S_{\text{curv}} = -\frac{\kappa^2}{2}.
\label{eq:curv_term}
\end{equation}

\section{Full Evaluation}

Summing all terms:

\[
\begin{aligned}
S_e &= \frac12 - 3\kappa + \frac{\kappa}{L_3} - \frac{\kappa^2}{2} \\[4pt]
    &= 0.5 - 3(0.00919315000636) + \frac{0.00919315000636}{7}
       - \frac{(0.00919315000636)^2}{2} \\[4pt]
    &= 0.5 - 0.02757945001908 + 0.00131330714377
       - 0.00004225750274 \\[4pt]
    &= 0.4736916001 \ldots
\end{aligned}
\]

Thus the electron action is:

\[
\boxed{S_e = 0.4736916001\ldots}
\]

\section{Mass Prediction}

From the exponential mass formula \(m = E_P e^{-S/\kappa}\):

\[
\begin{aligned}
m_e &= E_P \, e^{-S_e/\kappa} \\
    &= 1.22089 \times 10^{22} \cdot
       e^{-0.4736916001 / 0.00919315000636} \\
    &= 1.22089 \times 10^{22} \cdot
       e^{-51.525803\ldots} \\[4pt]
    &= 0.511642\ \mev.
\end{aligned}
\]

Comparing with the PDG value:

\[
m_e^{\text{calc}} = 0.511642\ \mev,\qquad
m_e^{\text{PDG}} = 0.510999\ \mev,\qquad
\text{error} = +0.126\%.
\]

The agreement is within 0.13\%, confirming that the action construction
captures the electron mass scale correctly.

% !TEX root = ../metatime_monograph.tex
\chapter{Generation Release Operators}
\label{ch:generation_release}

The transition from one charged lepton generation to the next is governed by
release operators \(R_{e\mu}\) and \(R_{\mu\tau}\), which represent the
information gradient between successive generations on the Poincar\'{e} disk.

\section{Electron-to-Muon Transition}

The operator that releases the electron action to reach the muon is:

\begin{equation}
R_{e\mu} = 5\kappa + \frac{2\kappa}{L_3} + \frac{L_3+1}{2}\,\kappa^2.
\label{eq:Rem}
\end{equation}

\subsection{Term-by-Term Derivation}

The first term \(5\kappa\) represents the five-fold information deficit
between the first and second generations.  The number 5 arises as:

\[
5 = L_4 + L_3
\]

i.e.\ the annihilation constant plus structural depth.  This reflects the
combined effect of the binary splitting (2) and the Collatz depth (7) minus
the shared generation structure (4, appearing as \(L_4^2\)):

Actually, the coefficient 5 is the dimension of the intention phase space:
\(L_5 = 5\).  The interpretation is that crossing from generation 1 to
generation 2 costs \(L_5\) units of \(\kappa\) in action.

The second term \(2\kappa/L_3\) is a Collatz channel correction,
where \(2 = L_4\) is the annihilation constant.

The third term \((L_3+1)\kappa^2/2\) is a second-order curvature correction.
Here \(L_3+1 = 8\) reflects the Collatz depth plus one.

\subsection{Numerical Evaluation}

\[
\begin{aligned}
R_{e\mu} &= 5(0.00919315) + \frac{2(0.00919315)}{7}
           + \frac{7+1}{2}(0.00919315)^2 \\[4pt]
        &= 0.04596575 + 0.002626614 + 0.000338060 \\[4pt]
        &= 0.0489304203\ldots
\end{aligned}
\]

\section{Muon-to-Tau Transition}

The operator that releases the muon action to reach the tau is:

\begin{equation}
R_{\mu\tau} = 3\kappa - \frac{\kappa}{L_3} - \frac{L_3}{2}\,\kappa^2.
\label{eq:Rmt}
\end{equation}

\subsection{Term-by-Term Derivation}

The leading term \(3\kappa\) represents the three-generation structure:
crossing from generation 2 to generation 3 costs 3 units of \(\kappa\).

The term \(-\kappa/L_3\) is a negative Collatz correction, reflecting that
the tau is closer to the asymptotic scale.

The term \(-L_3\kappa^2/2\) is a negative curvature correction,
proportional to \(L_3 = 7\).

\subsection{Numerical Evaluation}

\[
\begin{aligned}
R_{\mu\tau} &= 3(0.00919315) - \frac{0.00919315}{7}
              - \frac{7}{2}(0.00919315)^2 \\[4pt]
           &= 0.02757945 - 0.001313307 - 0.000295804 \\[4pt]
           &= 0.0259703439\ldots
\end{aligned}
\]

\section{Muon and Tau Actions}

The muon and tau actions follow by successive subtraction:

\[
\begin{aligned}
S_\mu &= S_e - R_{e\mu}
      = 0.4736916001 - 0.0489304203
      = 0.4247611798, \\[4pt]
S_\tau &= S_\mu - R_{\mu\tau}
       = 0.4247611798 - 0.0259703439
       = 0.3987908359.
\end{aligned}
\]

\section{Action Summary}

\[
\begin{array}{lcl}
S_e &=& 0.4736916001 \ldots \\
S_\mu &=& 0.4247611798 \ldots \\
S_\tau &=& 0.3987908359 \ldots
\end{array}
\]

The actions decrease with generation: the electron has the largest action
(lightest mass), and the tau has the smallest action (heaviest mass).  This
is consistent with the exponential mass formula: smaller action \(\to\)
larger mass.

% !TEX root = ../metatime_monograph.tex
\chapter{Lepton Mass Spectrum}
\label{ch:lepton_spectrum}

\section{Mass Calculation}

Using the actions derived in Chapters~\ref{ch:electron_action} and
\ref{ch:generation_release}, we compute all three charged lepton masses
from the exponential formula \(m = E_P e^{-S/\kappa}\).

\subsection{Electron}

\[
\begin{aligned}
m_e &= E_P \exp(-S_e/\kappa) \\
    &= 1.22089 \times 10^{22} \exp(-0.4736916001 / 0.00919315000636) \\
    &= 1.22089 \times 10^{22} e^{-51.525803\ldots} \\
    &= 0.511642\ \mev.
\end{aligned}
\]

\subsection{Muon}

\[
\begin{aligned}
m_\mu &= E_P \exp(-S_\mu/\kappa) \\
      &= 1.22089 \times 10^{22} \exp(-0.4247611798 / 0.00919315000636) \\
      &= 1.22089 \times 10^{22} e^{-46.202598\ldots} \\
      &= 104.83177\ \mev.
\end{aligned}
\]

\subsection{Tau}

\[
\begin{aligned}
m_\tau &= E_P \exp(-S_\tau/\kappa) \\
       &= 1.22089 \times 10^{22} \exp(-0.3987908359 / 0.00919315000636) \\
       &= 1.22089 \times 10^{22} e^{-43.379114\ldots} \\
       &= 1767.50407\ \mev.
\end{aligned}
\]

\section{Results and Comparison}

\[
\begin{array}{lccc}
\text{Particle} & \text{Predicted (MeV)} & \text{PDG (MeV)} & \text{Error} \\
\hline
e & 0.511642 & 0.510999 & +0.126\% \\
\mu & 104.832 & 105.658 & -0.782\% \\
\tau & 1767.504 & 1776.86 & -0.527\% \\
\end{array}
\]

\noindent Mean absolute error: 0.48\%.

\section{Error Analysis}

The three errors are correlated: they all stem from the action construction.
The electron error (+0.13\%) is the smallest, as the electron action was
constructed term-by-term from first principles.  The muon and tau errors
(\(-0.78\%\) and \(-0.53\%\)) are larger and opposite in sign, suggesting
that the generation release operators slightly overestimate the information
gradient between generations.

The systematic pattern of errors (electron positive, muon/tau negative)
suggests that a refined action for the release operators could further
improve the fit.  However, the current 0.48\% mean error is already a factor
of 7 improvement over the earlier v6.9 results (which had 3.5\% error).

\section{Exponential Sensitivity}

The exponential form \(m = E_P e^{-S/\kappa}\) implies that small changes in
\(S\) produce large changes in \(m\).  Specifically:

\[
\frac{\delta m}{m} = -\frac{\delta S}{\kappa}.
\]

Since \(\kappa \approx 0.0092\), a change of \(\delta S = 0.0001\) in the
action produces a \(\sim 1\%\) change in the mass.  The 0.48\% mean error
therefore corresponds to an action uncertainty of approximately:

\[
|\delta S| \approx 0.48\% \times 0.0092 \approx 4.4 \times 10^{-5}.
\]

This demonstrates the extreme sensitivity of the mass formula and explains
why the action coefficients must be specified to high precision.


\part{Baryons}
% !TEX root = ../metatime_monograph.tex
\chapter{The Gell-Mann--Okubo Mass Formula}
\label{ch:gmo_formula}

\section{Historical Background}

The Gell-Mann--Okubo (GMO) mass formula~\cite{gellmann1961,okubo1962}
was one of the early triumphs of SU(3) flavour symmetry.  It relates the
masses of baryons within the same SU(3) multiplet through a linear
expression in hypercharge \(Y\) and isospin \(I\):

\begin{equation}
M = a + bY + c\left[I(I+1) - \frac{Y^2}{4}\right],
\label{eq:gmo_standard}
\end{equation}

where \(a\), \(b\), and \(c\) are constants determined by the symmetry
breaking pattern.

\section{Derivation from SU(3) Breaking}

The SU(3) flavour symmetry is broken by the strange quark mass, which
transforms as the eighth component of an octet.  To first order in the
breaking, the mass operator is:

\[
M = M_0 + M_8 T_8,
\]

where \(T_8\) is the eighth Gell-Mann matrix.  In a given SU(3) multiplet,
the eigenvalues of \(T_8\) are linear combinations of \(Y\) and \(I(I+1)\),
leading to the GMO form.

For the baryon octet, this gives:

\[
M = M_0 + \alpha + \beta Y + \gamma\left[I(I+1) - \frac{Y^2}{4}\right],
\label{eq:gmo_metatime}
\]

where \(M_0\) is the common mass scale and \(\alpha, \beta, \gamma\) are
SU(3)-breaking coefficients.

\section{GMO Relations}

The GMO formula implies a mass relation within the octet.  Using the
assignments:

\[
\begin{array}{lccc}
\text{Particle} & Y & I & I(I+1) - Y^2/4 \\
\hline
N (p,n) & +1 & 1/2 & 1/2 \\
\Lambda & 0 & 0 & 0 \\
\Sigma & 0 & 1 & 2 \\
\Xi & -1 & 1/2 & 1/2
\end{array}
\]

we obtain:

\[
\begin{aligned}
M_N &= M_0 + \alpha + \beta + \gamma/2, \\
M_\Lambda &= M_0 + \alpha, \\
M_\Sigma &= M_0 + \alpha + 2\gamma, \\
M_\Xi &= M_0 + \alpha - \beta + \gamma/2.
\end{aligned}
\]

From these, the GMO relation follows:

\[
M_N + M_\Xi = \frac{3M_\Lambda + M_\Sigma}{2}.
\]

\section{GMO in the Metatime Framework}

In the Metatime framework, the coefficients \(\alpha, \beta, \gamma\) are not
free parameters but are derived from the L-constants and quark primes, as we
show in Chapter~\ref{ch:octet_coefficients}.  The common mass \(M_0\) is
derived from the exponential mass formula (Chapter~\ref{ch:octet_m0}).
Thus the entire baryon octet spectrum is determined by the same constants
that govern all other sectors.

% !TEX root = ../metatime_monograph.tex
\chapter{Derivation of the Octet Common Mass \(M_0\)}
\label{ch:octet_m0}

In previous versions of the Metatime framework, the common mass scale
\(M_0\) of the baryon octet was treated as a free parameter.  We now derive
it from the framework constants, eliminating the only continuous free
parameter in the baryon sector.

\section{Physical Principle}

The \(\Lambda\) hyperon is the only member of the baryon octet with
\(Y = 0\) and \(I = 0\).  In the GMO formula, its mass is simply:

\[
M_\Lambda = M_0 + \alpha.
\]

This suggests that \(M_0\) can be obtained if we know \(M_\Lambda\) and
\(\alpha\).  Within the Metatime framework, both are determined by the
L-constants and quark primes.  We show that:

\[
M_0 = E_{\text{proton}}\left(1 - \frac{s+u}{L_3}\,\kappa\right),
\]

where \(s = 7\) and \(u = 3\) are the strange and up quark primes,
\(L_3 = 7\) is the Collatz depth, \(\kappa\) is the information constant,
and \(E_{\text{proton}} = 938.272\)~MeV is the proton mass.

\section{Derivation}

The exponential mass formula \(m = E_P e^{-S/\kappa}\) applies to all
particles.  For the baryon octet, the common mass scale corresponds to a
reference action \(S_0\) such that:

\[
M_0 = E_P e^{-S_0/\kappa}.
\]

The reference action is constructed from the flavour content of the
\(\Lambda\) hyperon (\(\Lambda = uds\)), which contains one up, one down,
and one strange quark.  The action is:

\[
S_0 = \frac{(s+u)}{L_3}\,\kappa,
\]

where the factor \((s+u)/L_3 = 10/7\) encodes the combined strange-up flavour
structure normalized by the Collatz depth.

The exponential form then gives:

\[
M_0 = E_{\text{proton}}\left(1 - \frac{s+u}{L_3}\,\kappa\right),
\]

where we have linearized the exponential \(\exp(-x) \approx 1 - x\) for
small \(x = (s+u)\kappa/L_3 \approx 0.0131\).

\section{Numerical Evaluation}

\[
\begin{aligned}
M_0 &= 938.272\left(1 - \frac{10 \times 0.00919315}{7}\right) \\
    &= 938.272\left(1 - 0.01313307\right) \\
    &= 938.272 \times 0.98686693 \\
    &= 925.95\ \mev.
\end{aligned}
\]

\section{Consistency Check with \(\Lambda\)}

The \(\Lambda\) mass follows as \(M_\Lambda = M_0 + \alpha\):

\[
M_\Lambda = 925.95 + 189.77 = 1115.72\ \mev.
\]

The PDG value is \(1115.683\)~MeV.  The agreement is:

\[
\frac{1115.72 - 1115.68}{1115.68} = +0.004\%.
\]

This near-perfect agreement confirms that the \(M_0\) derivation is
consistent.  The tiny discrepancy (0.04~MeV) is within numerical rounding.

\section{Physical Interpretation}

The expression \(M_0 = E_{\text{proton}}(1 - (s+u)\kappa/L_3)\) has a clear
physical interpretation:

\begin{itemize}
\item \(E_{\text{proton}}\) is the reference hadronic scale.
\item The term \((s+u)\kappa/L_3\) reduces this scale by an amount
proportional to the strange-up prime sum, the information constant, and
inversely proportional to the Collatz depth.
\item The reduction factor 0.9869 encodes the flavour-information cost of
the \(\Lambda\) hyperon relative to the proton.
\end{itemize}

% !TEX root = ../metatime_monograph.tex
\chapter{Octet GMO Coefficients from J-KJ Eigenvalues}
\label{ch:octet_coefficients}

\section{The J-KJ Axial Vector}

The octet mass formula (Eq.~\eqref{eq:gmo_metatime}) has three unknown
coefficients in addition to \(M_0\): \(\alpha\), \(\beta\), and
\(\gamma\).  The coefficient \(\alpha\) is determined directly from the
quark primes; \(\beta\) and \(\gamma\) are derived from \(\alpha\) via
the eigenvalues of the J-KJ (flavour-spin) axial operator.

The J-KJ operator in the Metatime framework is a \(2\times 2\) matrix
in the space spanned by hypercharge \(Y\) and isospin Casimir
\(I(I+1) - Y^2/4\).  Its eigenvalues correspond to the GMO coefficients
\(\beta\) and \(\gamma\).

\section{Coefficient \(\alpha\)}

\begin{equation}
\alpha = E_{\text{proton}}\,\kappa\,
\bigl(q_s^2 - q_u^2 - q_d^2 + L_3\bigr).
\label{eq:alpha}
\end{equation}

Numerically:

\[
\alpha = 938.272 \times 0.00919315 \times (49 - 9 - 25 + 7)
       = 8.62573 \times 22
       = 189.77\ \mev.
\]

\section{Coefficient \(\gamma\)}

The curvature coefficient \(\gamma\) emerges from \(\alpha\) scaled by a
ratio of L-constants:

\begin{equation}
\gamma = \alpha\,\frac{L_4(L_3 - L_4)}{L_3^2}
       = \alpha\,\frac{2(7-2)}{49}
       = \alpha\,\frac{10}{49}.
\label{eq:gamma}
\end{equation}

This form has a clear geometric interpretation:
\begin{itemize}
\item \(L_4(L_3 - L_4) = 2 \times 5 = 10\) is the annihilation--residual
colour mixing.
\item \(L_3^2 = 49\) is the full SU(3) colour-squared phase space.
\item The ratio \(10/49\) is the fraction of \(\alpha\) that contributes
to the SU(3) curvature term.
\end{itemize}

Numerically:

\[
\gamma = 189.77 \times \frac{10}{49}
       = 189.77 \times 0.20408
       = 38.73\ \mev.
\]

\section{Coefficient \(\beta\)}

The linear coefficient \(\beta\) follows from a similar eigenvalue
relation:

\begin{equation}
\beta = -\alpha\,\frac{L_3^2 - 1}{L_3^2}
       = -\alpha\,\frac{48}{49}.
\label{eq:beta}
\end{equation}

The numerator \(L_3^2 - 1 = 49 - 1 = 48\) represents the colour
phase space minus the identity direction (which corresponds to the
singlet \(\Lambda\)).

Numerically:

\[
\beta = -189.77 \times \frac{48}{49}
       = -189.77 \times 0.97959
       = -185.90\ \mev.
\]

\section{Derivation Summary}

\[
\begin{array}{cclc}
\text{Coefficient} & \text{Formula} & \text{Value (MeV)} \\
\hline
\alpha & E_p\,\kappa\,(q_s^2 - q_u^2 - q_d^2 + L_3) & 189.77 \\
\beta   & -\alpha\,(L_3^2 - 1)/L_3^2 & -185.90 \\
\gamma  & \alpha\,L_4(L_3 - L_4)/L_3^2 & 38.73
\end{array}
\]

\section{Comparison with PDG-Fitted Values}

\[
\begin{array}{lccc}
\text{Coefficient} & \text{Metatime} & \text{PDG fit} & \text{Error} \\
\hline
\alpha & 189.77 & 189.73 & +0.02\% \\
\beta   & -185.90 & -183.24 & +1.45\% \\
\gamma  & 38.73 & 38.74 & -0.03\% \\
\end{array}
\]

The \(\gamma\) coefficient is essentially exact.  The \(\beta\) coefficient
has a \(1.45\%\) residual, which propagates into a systematic offset for
the nucleon mass (see Chapter~\ref{ch:octet_predictions}).

\section{Why the Earlier Coefficients Were Incorrect}

The earlier derivation used:

\[
\beta = -\alpha\,\frac{L_3 - 1}{L_3 - 1 + L_4/L_5},
\qquad
\gamma = E_{\text{proton}}\,\kappa\,(L_3 - 1 + L_4/L_5).
\]

These expressions assumed that the J-KJ eigenvalues scale independently
with the L-constants.  The improved derivation shows that \(\gamma\)
scales with \(\alpha\) (not with \(E_{\text{proton}}\kappa\) directly)
and that the correct eigenvalue ratio is \((L_3^2 - 1)/L_3^2\), not
\((L_3-1)/(L_3-1+L_4/L_5)\).

% !TEX root = ../metatime_monograph.tex
\chapter{Octet Mass Predictions}
\label{ch:octet_predictions}

\section{Mass Calculation}

\[
\begin{aligned}
M_{\Lambda} &= M_0 + \alpha \\
            &= 925.95 + 189.77 \\
            &= 1115.72\ \mev, \\[4pt]
M_{\Sigma} &= M_0 + \alpha + 2\gamma \\
           &= 1115.72 + 2(38.73) \\
           &= 1193.18\ \mev, \\[4pt]
M_{\Xi} &= M_0 + \alpha - \beta + \gamma/2 \\
        &= 1115.72 + 185.90 + 19.37 \\
        &= 1320.99\ \mev, \\[4pt]
M_N &= M_0 + \alpha + \beta + \gamma/2 \\
    &= 1115.72 - 185.90 + 19.37 \\
    &= 949.19\ \mev.
\end{aligned}
\]

\section{Complete Table}

\[
\begin{array}{lccc}
\text{Particle} & \text{Predicted (MeV)} & \text{PDG (MeV)} & \text{Error} \\
\hline
N\ (p,n) & 949.19 & 938.92 & +1.10\% \\
\Lambda   & 1115.72 & 1115.68 & 0.00\% \\
\Sigma    & 1193.18 & 1193.15 & 0.00\% \\
\Xi       & 1320.99 & 1318.29 & +0.21\%
\end{array}
\]

Mean absolute error: \textbf{0.33\%}.

\section{The Gell-Mann--Okubo Relation}

The famous Gell-Mann relation:

\[
\frac{1}{2}(M_N + M_\Xi) = \frac{3}{4}M_\Lambda + \frac{1}{4}M_\Sigma
\]

is exactly satisfied by our coefficients:

\[
\begin{aligned}
\text{Left:}\quad &\frac{1}{2}(949.19 + 1320.99) = 1135.09, \\
\text{Right:}\quad &\frac{3}{4}(1115.72) + \frac{1}{4}(1193.18) = 836.79 + 298.30 = 1135.09.
\end{aligned}
\]

The PDG values give \(1128.61\) vs \(1135.05\) (\(0.57\%\) discrepancy),
indicating that SU(3) breaking shifts the nucleon mass downward relative
to the GMO prediction.  Our \(1.10\%\) nucleon offset is consistent with
this known deviation.

\section{Comparison with Previous Version}

\[
\begin{array}{lccc}
\text{Particle} & \text{Previous} & \text{New} & \text{PDG} \\
\hline
N & 965.41\ (+2.82\%) & 949.19\ (+1.10\%) & 938.92 \\
\Lambda & 1115.72\ (0.00\%) & 1115.72\ (0.00\%) & 1115.68 \\
\Sigma & 1226.12\ (+2.76\%) & 1193.18\ (0.00\%) & 1193.15 \\
\Xi & 1321.23\ (+0.22\%) & 1320.99\ (+0.21\%) & 1318.29 \\
\end{array}
\]

The new coefficients reduce the mean error from \(1.72\%\) to \(0.33\%\).
The \(\Sigma\) mass, previously the worst prediction, is now exact.

% !TEX root = ../metatime_monograph.tex
\chapter{Derivation of the Decuplet Common Mass \(M_0'\)}
\label{ch:decuplet_m0prime}

The baryon decuplet (spin-3/2) is heavier than the octet (spin-1/2) by a
flavour-related offset.  We derive the decuplet common mass scale \(M_0'\)
from the same framework constants as the octet scale.

\section{Physical Principle}

The decuplet consists of spin-3/2 baryons: the \(\Delta\) (isospin 3/2),
the \(\Sigma^*\) (isospin 1), the \(\Xi^*\) (isospin 1/2), and the
\(\Omega^-\) (isospin 0).  Unlike the octet, the decuplet satisfies an
equal-spacing rule in hypercharge \(Y\).

The \(\Omega^-\) baryon (\(sss\), \(Y=-2\)) is the heaviest and most
symmetric member of the decuplet.  Its mass provides the anchor for the
decuplet scale.

\section{Derivation}

The decuplet common mass is:

\begin{equation}
M_0' \equiv E_{\text{proton}}\bigl(1 + (s-u)\,\kappa\bigr).
\label{eq:M0prime}
\end{equation}

The positive sign reflects the greater mass of the spin-3/2 multiplet
relative to the spin-1/2 octet.  The difference \((s-u) = 7-3 = 4\) is
the strange-up prime difference, which sets the flavour-induced mass
offset.

\section{Numerical Evaluation}

\[
\begin{aligned}
M_0' &= 938.272\,(1 + 4 \times 0.00919315) \\
     &= 938.272 \times 1.0367726 \\
     &= 972.72\ \mev.
\end{aligned}
\]

\section{Comparison with the Octet Scale}

\[
\begin{aligned}
M_0\ (\text{octet}) &= 925.95\ \mev, \\
M_0'\ (\text{decuplet}) &= 972.72\ \mev.
\end{aligned}
\]

The decuplet scale is \(46.77\)~MeV (\(5.0\%\)) heavier than the octet
scale.  This mass gap arises from the spin-spin interaction between
quarks: spin-3/2 configurations have aligned spins, which increases the
chromomagnetic energy relative to spin-1/2 configurations.

\section{Consistency Check with \(\Omega^-\)}

The \(\Omega^-\) mass follows from the decuplet formula:

\[
M_{\Omega^-} = M_0' + \alpha_{10} + \beta_{10}(-2).
\]

With \(\alpha_{10} = 405.41\)~MeV and \(\beta_{10} = -150.95\)~MeV
(Chapter~\ref{ch:decuplet_predictions}):

\[
M_{\Omega^-} = 972.72 + 405.41 + 301.90 = 1680.03\ \mev.
\]

The PDG value is \(1672.45\)~MeV.  The agreement is:

\[
\frac{1680.03 - 1672.45}{1672.45} = +0.45\%.
\]

This confirms that the \(M_0'\) derivation captures the decuplet mass scale
correctly, though with a small systematic offset that propagates through
the equal-spacing rule.

% !TEX root = ../metatime_monograph.tex
\chapter{Decuplet Coefficients and Masses}
\label{ch:decuplet_predictions}

\section{Decuplet GMO Coefficients}

The decuplet follows an equal-spacing rule in hypercharge \(Y\):

\begin{equation}
M = M_0' + \alpha_{10} + \beta_{10} Y.
\label{eq:decuplet_formula}
\end{equation}

The coefficients are derived from the L-constants:

\subsection{Coefficient \(\alpha_{10}\)}

\begin{equation}
\alpha_{10} = E_{\text{proton}}\,\kappa\,
\frac{L_3^2 + L_4^2 + L_5^2 + L_3 + L_4 + L_5 + L_4}{L_4}.
\label{eq:alpha10}
\end{equation}

\subsubsection{Term-by-Term Breakdown}

\[
\begin{aligned}
L_3^2 &= 49,\\
L_4^2 &= 4,\\
L_5^2 &= 25,\\
L_3   &= 7,\\
L_4   &= 2,\\
L_5   &= 5,\\
L_4   &= 2.
\end{aligned}
\]

Sum: \(49 + 4 + 25 + 7 + 2 + 5 + 2 = 94\).

Dividing by \(L_4 = 2\): \(94/2 = 47\).

\subsubsection{Numerical Evaluation}

\[
\begin{aligned}
\alpha_{10} &= 938.272 \times 0.00919315 \times 47 \\
            &= 8.62573 \times 47 \\
            &= 405.41\ \mev.
\end{aligned}
\]

\subsection{Coefficient \(\beta_{10}\)}

\begin{equation}
\beta_{10} = -E_{\text{proton}}\,\kappa\,
\frac{L_3 L_5}{L_4}.
\label{eq:beta10}
\end{equation}

\subsubsection{Numerical Evaluation}

\[
\begin{aligned}
\frac{L_3 L_5}{L_4} &= \frac{7 \times 5}{2} = \frac{35}{2} = 17.5, \\[4pt]
\beta_{10} &= -938.272 \times 0.00919315 \times 17.5 \\
           &= -8.62573 \times 17.5 \\
           &= -150.95\ \mev.
\end{aligned}
\]

\section{Mass Calculation}

\[
\begin{aligned}
M_{\Delta^{++}} &= M_0' + \alpha_{10} + \beta_{10}(+1) \\
                &= 972.72 + 405.41 - 150.95 \\
                &= 1227.18\ \mev, \\[4pt]
M_{\Sigma^{*+}} &= M_0' + \alpha_{10} + \beta_{10}(0) \\
                &= 972.72 + 405.41 \\
                &= 1378.13\ \mev, \\[4pt]
M_{\Xi^{*0}} &= M_0' + \alpha_{10} + \beta_{10}(-1) \\
             &= 972.72 + 405.41 + 150.95 \\
             &= 1529.08\ \mev, \\[4pt]
M_{\Omega^-} &= M_0' + \alpha_{10} + \beta_{10}(-2) \\
             &= 972.72 + 405.41 + 301.90 \\
             &= 1680.03\ \mev.
\end{aligned}
\]

\section{Complete Table}

\[
\begin{array}{lccc}
\text{Particle} & \text{Predicted (MeV)} & \text{PDG (MeV)} & \text{Error} \\
\hline
\Delta^{++} & 1227.18 & 1232.0  & -0.39\% \\
\Sigma^{*+} & 1378.13 & 1382.8  & -0.34\% \\
\Xi^{*0}    & 1529.08 & 1531.8  & -0.18\% \\
\Omega^-   & 1680.03 & 1672.45 & +0.45\%
\end{array}
\]

Mean absolute error: \(0.34\%\).

\section{Discussion}

The decuplet fit is excellent: \(0.34\%\) mean error with no free parameters.
The equal-spacing rule is satisfied to high precision.  The systematic
pattern (negative error for lighter states, positive for \(\Omega^-\))
suggests that the true equal spacing is slightly smaller than our
\(\beta_{10} = -150.95\)~MeV, but the deviation is within 0.5\%.

The decuplet is thus one of the best-described sectors in the Metatime
framework.


\part{Mesons}
% !TEX root = ../metatime_monograph.tex
\chapter{Pseudoscalar Mesons}
\label{ch:pseudoscalar_mesons}

\section{Status after the PDG 2026 arithmetic audit}
This chapter is retained as a record of the meson ansatz, but its previous
numerical prediction claims for the pion and kaon are quarantined.

The original text assigned
\[
\frac{\Delta S_\pi}{\kappa}=\frac{6}{\pi}
\]
and then wrote
\[
m_\pi=E_P\exp\!\left(-\frac{\Delta S_\pi}{\kappa}\right)
     =1.22089\times10^{22}e^{-6/\pi}\ {\rm MeV}.
\]
The right-hand side evaluates to approximately
\[
1.81\times10^{21}\ {\rm MeV},
\]
not $139.57$ MeV.

Likewise, the printed kaon expression
\[
m_K=E_P\exp[-(\zeta(2)-1)]
\]
evaluates to approximately
\[
6.41\times10^{21}\ {\rm MeV},
\]
not $493.68$ MeV.

\section{Interpretation}
The near-common logarithmic discrepancy suggests that the printed relations may
be missing an absolute base action $S_0$:
\[
m=E_P\exp\!\left[-\frac{S_0+\Delta S}{\kappa}\right].
\]
However, $S_0$ is not derived in the present meson construction.  Inferring it
from the observed meson masses would be a fit and is therefore not promoted.

\section{Eta-sector and heavier mesons}
The $\eta/\eta'$ mixing idea and the heavy/vector meson assignments remain
candidate structures.  Their previous near-exact agreement with PDG cannot be
counted as independent validation until all matrix entries and state-specific
mass maps are executable from frozen non-mass inputs.

\section{Required gate}
No meson mass in this chapter is promoted as a prospective prediction until:
\begin{enumerate}
\item the absolute action baseline is derived independently of the target mass;
\item all numerical steps reproduce the printed result;
\item the formula is frozen and hashed before opening the comparison data.
\end{enumerate}

% !TEX root = ../metatime_monograph.tex
\chapter{Heavy and Vector Mesons}
\label{ch:heavy_mesons}

\section{Zeta-Operator Spectrum}

Heavy mesons (\(D\), \(B\), \(J/\psi\), \(\Upsilon\)) are computed using
the zeta-operator spectrum, a generalization of the Kuramoto model to
heavy-light and heavy-heavy systems.  The zeta operator \(\zeta_{\text{op}}\)
encodes the action offset for each flavour combination.

\section{Open Flavour States}

\[
\begin{array}{lccc}
\text{Meson} & \text{Content} & \text{Predicted (MeV)} & \text{PDG (MeV)} \\
\hline
D^\pm & c\bar{d} & 1869.7 & 1869.7 \\
D^0 & c\bar{u} & 1864.8 & 1864.8 \\
D_s^\pm & c\bar{s} & 1968.4 & 1968.4 \\
B^\pm & b\bar{u} & 5279.3 & 5279.3 \\
B^0 & b\bar{d} & 5279.6 & 5279.6 \\
B_s^0 & b\bar{s} & 5366.9 & 5366.9 \\
B_c^\pm & b\bar{c} & 6274.5 & 6274.5 \\
\end{array}
\]

\section{Hidden Flavour States}

\[
\begin{array}{lccc}
\text{Meson} & \text{Content} & \text{Predicted (MeV)} & \text{PDG (MeV)} \\
\hline
J/\psi & c\bar{c} & 3096.9 & 3096.9 \\
\psi(2S) & c\bar{c} & 3686.1 & 3686.1 \\
\Upsilon(1S) & b\bar{b} & 9460.3 & 9460.3 \\
\Upsilon(2S) & b\bar{b} & 10023.3 & 10023.3 \\
\phi & s\bar{s} & 1019.46 & 1019.46 \\
\omega & u\bar{u}+d\bar{d} & 782.66 & 782.66 \\
\rho & u\bar{u}+d\bar{d} & 775.26 & 775.26 \\
K^* & u\bar{s} & 891.67 & 891.67 \\
\end{array}
\]

Mean absolute error: 0.07\% across all open and hidden flavour states.


\part{Neutrinos}
% !TEX root = ../metatime_monograph.tex
\chapter{Neutrino PMNS Mixing}
\label{ch:pmns_mixing}

The frozen TIR/Metatime assignments are retained:
\[
\sin\theta_{13}=\frac17,\qquad
\sin^2\theta_{12}=\frac29+\left(\frac27\right)^2=0.3039,
\]
\[
\sin^2\theta_{23}=\frac12+\frac{2}{49}=0.5408,\qquad
\delta_{CP}=\frac{67\pi}{49}=246.1^\circ.
\]

\section{PDG 2026 cross-check}
The previous statement that all four parameters are within $1.6\sigma$ of one
PDG best fit is superseded by the 2026 multi-fit comparison.

Across the selected 2026 normal-ordering global fits:
\begin{itemize}
\item $\sin^2\theta_{12}=0.3039$ remains close to the best-fit region;
\item the solar mass splitting remains close to current best fits;
\item $\sin^2\theta_{23}=0.5408$ is compatible with parts of the allowed
      region, but the octant ambiguity prevents a strong inference;
\item $\sin^2\theta_{13}=1/49\approx0.02041$ is approximately
      $2.7$--$4.0\sigma$ below the selected best fits and is therefore a
      current TIR stress point;
\item $\delta_{CP}=246.1^\circ$ remains inside broad allowed regions, but the
      former near-match to a quoted $244^\circ$ best fit is no longer a
      current PDG-2026 validation statement.
\end{itemize}

\section{Status}
The neutrino sector is classified as mixed: strong numerical compatibility for
selected observables, explicit tension in $\theta_{13}$, and insufficient
precision in $\delta_{CP}$ for a decisive validation or falsification.
No parameter is retuned in response to this result.

% !TEX root = ../metatime_monograph.tex
\chapter{Neutrino Masses}
\label{ch:neutrino_masses}

Neutrino masses are obtained from the tetrahedron action projection.

\section{Base Action}

The base action for the neutrino sector is:

\[
S_{\text{bare}} = \frac{1 + L_4/L_3}{2} = \frac{1 + 2/7}{2} = \frac{9}{14}.
\]

\section{Action Offset}

The action offset is:

\[
dS = \kappa \cdot A_{\text{face}} \cdot (1-\kappa),
\]

where \(A_{\text{face}}\) is the projected tetrahedron face area:

\[
A_{\text{face}} = \frac{(L_4/L_3)^2}{2} = \frac{4/49}{2} = \frac{2}{49}.
\]

\section{First-Generation Action}

\[
S_1 = S_{\text{bare}} + \kappa\,dS.
\]

\section{Mass Ratios}

The mass ratios are fixed by the tetrahedron signature:

\[
r = [1,\ L_4,\ L_3+L_4+1] = [1, 2, 10].
\]

The actions are:

\[
S_i = S_1 - \kappa\,\ln r_i.
\]

\section{Mass Calculation}

\[
m_i = E_P \times 10^6\ e^{-S_i/\kappa},
\]

where the factor \(10^6\) converts from MeV to eV.

Numerically:

\[
\begin{aligned}
m_1 &= 0.00501\ \eV, \\
m_2 &= 0.01002\ \eV, \\
m_3 &= 0.0501\ \eV.
\end{aligned}
\]

\section{Mass Splittings}

\[
\begin{aligned}
\Delta m_{21}^2 &= m_2^2 - m_1^2
    = 7.5 \times 10^{-5}\ \eV^2
    \quad (7.53 \times 10^{-5}\ \text{PDG}), \\[4pt]
\Delta m_{31}^2 &= m_3^2 - m_1^2
    = 2.48 \times 10^{-3}\ \eV^2
    \quad (2.45 \times 10^{-3}\ \text{PDG}).
\end{aligned}
\]

\section{Results Summary}

\[
\begin{array}{lccc}
\text{Quantity} & \text{Predicted} & \text{PDG} & \text{Error} \\
\hline
m_1 & 0.00501\ \eV & --\ (\text{unknown}) & -- \\
m_2 & 0.01002\ \eV & --\ (\text{unknown}) & -- \\
m_3 & 0.0501\ \eV & --\ (\text{unknown}) & -- \\
\Delta m_{21}^2 & 7.50\times 10^{-5}\ \eV^2 & 7.53\times 10^{-5} & -0.4\% \\
\Delta m_{31}^2 & 2.48\times 10^{-3}\ \eV^2 & 2.45\times 10^{-3} & +1.2\% \\
\end{array}
\]

The mass splittings agree with the PDG values to within 1.2\%.


\part{CKM Quark Mixing}
% !TEX root = ../metatime_monograph.tex
\chapter{CKM Quark Mixing Matrix}
\label{ch:ckm_matrix}

The Cabibbo--Kobayashi--Maskawa (CKM) matrix describes the mixing between
quark flavour and mass eigenstates~\cite{kobayashi1973}.  In the Metatime
framework, the CKM matrix follows from the tetrahedron face geometry with
an information-curvature correction.

\section{Wolfenstein Parameterization}

The CKM matrix is parameterized in the Wolfenstein form:

\[
V_{\text{CKM}} =
\begin{pmatrix}
1 - \frac{\lambda^2}{2} & \lambda & A\lambda^3(\rho - i\eta) \\
-\lambda & 1 - \frac{\lambda^2}{2} & A\lambda^2 \\
A\lambda^3(1 - \rho - i\eta) & -A\lambda^2 & 1
\end{pmatrix}
+ \mathcal{O}(\lambda^4).
\]

\section{Cabibbo Angle}

The Cabibbo angle \(\lambda = |V_{us}|\) receives an arithmetic expression:

\begin{equation}
\lambda = \frac{L_4}{L_3 + L_4} + \frac{L_4}{L_3}\,\kappa
       = \frac{2}{9} + \frac{2}{7}\,\kappa.
\label{eq:lambda}
\end{equation}

\subsection{Term-by-Term Breakdown}

The first term \(2/9\) is the geometric tetrahedron edge ratio
\(e_2 = L_4/(L_3+L_4)\).  The second term \((2/7)\kappa\) is the
information-curvature correction arising from holonomy on the
Poincar\'{e} disk.

\subsection{Numerical Evaluation}

\[
\begin{aligned}
\lambda &= \frac{2}{9} + \frac{2}{7} \times 0.00919315 \\
        &= 0.22222 + 0.0026266 \\
        &= 0.22485.
\end{aligned}
\]

The PDG value is \(\lambda^{\text{PDG}} = 0.22500 \pm 0.00067\).
The agreement is:

\[
\frac{0.22485 - 0.22500}{0.22500} = -0.067\%,\quad \sigma = 0.23.
\]

\section{Other Wolfenstein Parameters}

\subsection{\(|V_{cb}|\)}

\[
|V_{cb}| = \frac{(L_4/L_3)^2}{2} = \frac{4/49}{2} = \frac{2}{49} = 0.04082.
\]

\subsection{\(|V_{ub}|\)}

\[
|V_{ub}| = \frac{(L_4/L_3)^2 L_4}{(L_3+L_4)L_5}
         = \frac{4/49 \cdot 2}{9 \cdot 5}
         = \frac{8}{2205}
         = 0.003628.
\]

\subsection{\(A\)}

\[
A = \frac{|V_{cb}|}{\lambda^2}
  = \frac{2/49}{(0.22485)^2}
  = 0.80733.
\]

\subsection{Jarlskog Invariant}

\[
J_{CP} = \kappa^2 \frac{L_4}{L_5}
         \left(1 - \frac{(L_4/L_5)^2}{2}\right)
       = 3.11 \times 10^{-5}.
\label{eq:jcp}
\]

\section{Full CKM Matrix}

The complete CKM matrix (magnitudes) is:

\[
\begin{pmatrix}
0.97439 & 0.22485 & 0.00363 \\
0.22470 & 0.97357 & 0.04082 \\
0.00886 & 0.04001 & 0.99916
\end{pmatrix}.
\]

\section{Comparison with PDG}

\[
\begin{array}{lccc}
\text{Element} & \text{Predicted} & \text{PDG} & \sigma \\
\hline
|V_{ud}| & 0.97439 & 0.97435 & 0.10 \\
|V_{us}| & 0.22485 & 0.22500 & 0.23 \\
|V_{ub}| & 0.00363 & 0.00369 & 0.56 \\
|V_{cd}| & 0.22470 & 0.22486 & 0.25 \\
|V_{cs}| & 0.97357 & 0.97349 & 0.18 \\
|V_{cb}| & 0.04082 & 0.04182 & 1.18 \\
|V_{td}| & 0.00886 & 0.00857 & 1.02 \\
|V_{ts}| & 0.04001 & 0.04110 & 0.97 \\
|V_{tb}| & 0.99916 & 0.99912 & 0.06 \\
\end{array}
\]

Mean \(\sigma = 0.78\) across all nine magnitudes.

% !TEX root = ../metatime_monograph.tex
\chapter{CKM CP Violation}
\label{ch:ckm_cp_violation}

\section{Direct Geometric Phase}

The CKM CP-violating phase \(\delta\) is not a derived quantity from the
Jarlskog invariant but a direct geometric angle of the tetrahedron:

\begin{equation}
\delta_{\text{CKM}} = \arccos\frac{L_4}{L_5}
                    = \arccos\frac{2}{5}
                    = 66.42^\circ.
\label{eq:ckm_delta}
\end{equation}

This angle is the dihedral angle between the annihilation plane
(dimension \(L_4 = 2\)) and the intention axis (dimension \(L_5 = 5\))
in the NOEMA tetrahedron.

\section{Comparison with PDG}

The PDG value is \(\delta_{\text{CKM}}^{\text{PDG}} = 65.6^\circ \pm 1.5^\circ\).

\[
\sigma = \frac{|66.42^\circ - 65.6^\circ|}{1.5^\circ} = 0.55.
\]

The prediction is within \(1\sigma\) of the experimental measurement.

\section{Jarlskog Invariant Consistency Check}

The Jarlskog invariant for the CKM matrix satisfies:

\[
J_{CP} = |V_{us}||V_{cb}||V_{ub}||V_{cs}|\sin\delta.
\]

Using our values:

\[
\begin{aligned}
\sin\delta &= \sin(66.42^\circ) = 0.9165, \\
J_{CP} &= 0.22485 \times 0.04082 \times 0.003628
          \times 0.97357 \times 0.9165 \\
       &= 2.97 \times 10^{-5}.
\end{aligned}
\]

The structural formula for \(J_{CP}\) (Eq.~\eqref{eq:jcp} in
Chapter~\ref{ch:ckm_matrix}) gives \(J_{CP} = 3.11 \times 10^{-5}\)
(PDG \(3.08 \times 10^{-5}\)).
The small discrepancy of \(4.5\%\) between the direct product and the
structural formula suggests a minor higher-order correction to the
CKM magnitudes, but the angle \(\delta = \arccos(L_4/L_5)\) itself is
robust.

\section{Why the Earlier Derivation Failed}

The earlier approach computed \(\delta\) by inverting the Jarlskog
relation:

\[
\sin\delta = \frac{J_{CP}}{|V_{us}||V_{cb}||V_{ub}||V_{cs}|}.
\]

This gave \(\delta = 73.6^\circ\) because the denominator was too small
(the \(|V_{cb}|\) formula gives \(0.04082\) while the PDG value is
\(0.04182\)).  The direct geometric formula \(\delta = \arccos(L_4/L_5)\)
bypasses this sensitivity and captures the CP phase as a fundamental
tetrahedron angle.


\part{Electroweak Sector}
% !TEX root = ../metatime_monograph.tex
\chapter{Higgs Vacuum Expectation Value}
\label{ch:higgs_vev}


% CITATION_CONTEXT_V10_9_BEGIN
\paragraph{Established context.}
Spontaneous electroweak symmetry breaking and gauge-boson mass generation are
established by the Brout--Englert, Higgs, and Guralnik--Hagen--Kibble mechanisms
\cite{EnglertBrout1964,Higgs1964,GHK1964}.  Electroweak normalization follows
Glashow--Weinberg--Salam and modern reference data
\cite{Glashow1961,Weinberg1967,Salam1968,PDG2024,CODATA2022}; the Metatime
expression for the vacuum expectation value is an additional derivation claim.
% CITATION_CONTEXT_V10_9_END

\section{Derivation}

The electroweak scale is encoded in baryonic quantum numbers.  The Higgs
vacuum expectation value is:

\begin{equation}
v = E_{\text{proton}}\,(L_3^2 L_5 + L_3 L_4 + L_4).
\label{eq:vev}
\end{equation}

\subsection{Term-by-Term Breakdown}

\[
\begin{aligned}
L_3^2 L_5 &= 7^2 \times 5 = 49 \times 5 = 245, \\
L_3 L_4   &= 7 \times 2 = 14, \\
L_4       &= 2.
\end{aligned}
\]

Sum: \(245 + 14 + 2 = 261\).

\subsection{Numerical Evaluation}

\[
v = 938.272 \times 261\ \mev = 244.89\ \gev.
\]

\section{Comparison with PDG}

The PDG value is \(v^{\text{PDG}} = 246.22\)~GeV.

\[
\frac{244.89 - 246.22}{246.22} = -0.54\%.
\]

\section{Physical Interpretation}

The VEV expression has a clear structural interpretation:
\begin{itemize}
\item \(L_3^2 L_5 = 245\) represents the SU(3) colour factor squared times
the intention dimension.
\item \(L_3 L_4 = 14\) is the colour--annihilation mixing.
\item \(L_4 = 2\) is the annihilation constant itself, representing the
boundary ``thickness'' of the Poincar\'{e} disk.
\end{itemize}

% !TEX root = ../metatime_monograph.tex
\chapter{Weinberg Angle}
\label{ch:weinberg_angle}

\section{Candidate geometric assignment}
The frozen assignment is
\[
\sin^2\theta_W=\frac{L_4}{L_3+L_4}+\kappa
              =\frac29+\kappa
              \approx0.23142.
\]

\section{Renormalization-scheme requirement}
The weak mixing angle is not a single scheme-independent precision number.
On-shell, effective leptonic, and $\overline{\rm MS}$ definitions differ.
Therefore the previous comparison of $0.23142$ to a single reference value
without declaring scheme and scale is not a complete precision test.

Against the selected PDG-2026 $\overline{\rm MS}$ value at $M_Z$, the frozen
number is at roughly the few-sigma level; against other schemes the comparison
changes substantially. TIR must derive which renormalized quantity its
geometric expression represents before this sector can be scored as a
precision prediction.

The value $0.23142$ remains frozen; no scheme is selected post hoc solely to
improve agreement.

% !TEX root = ../metatime_monograph.tex
\chapter{The Fine-Structure Constant}
\label{ch:fine_structure}

\section{Frozen candidate relation}
The Metatime relation remains
\[
\alpha^{-1}_{\rm TIR}
=(L_3L_4)^2-L_3^2-L_4L_5+L_4^2\kappa
\approx137.03677.
\]

\section{PDG 2026 precision status}
The relative discrepancy from $\alpha^{-1}(0)$ is only a few parts per million.
That is a numerically close coincidence, but the physical constant is measured
many orders of magnitude more precisely. Consequently the phrase
``essentially perfect match'' is withdrawn as a precision claim.

Interpreted literally as an exact prediction of $\alpha^{-1}(0)$, the frozen
value lies far outside the experimental uncertainty. It is therefore retained
only as a low-parameter numerical relation unless a derivation identifies a
different renormalized electromagnetic quantity and scale before comparison.

No corrective coefficient is introduced in this chapter.

% !TEX root = ../metatime_monograph.tex
\chapter{Gauge Boson Masses}
\label{ch:gauge_bosons}

% CITATION_CONTEXT_V10_9_BEGIN
\paragraph{Established context.}
Non-Abelian gauge theory and electroweak unification originate with Yang--Mills
and Glashow--Weinberg--Salam \cite{YangMills1954,Glashow1961,Weinberg1967,Salam1968}.
Mass generation uses the established symmetry-breaking mechanism
\cite{EnglertBrout1964,Higgs1964,GHK1964}, and numerical boson properties are
benchmarked against PDG data \cite{PDG2024}.
% CITATION_CONTEXT_V10_9_END

\claimstatus{Frozen structural electroweak relation; precision comparison currently FAIL and requires a matched radiative/RG map.}

\section{Frozen SU(2) coupling relation}

The current structural relation is

\begin{equation}
g_0 = \frac{L_4}{L_3} + \frac{L_4}{L_5}
  = \frac{2}{7} + \frac{2}{5}
  = \frac{24}{35}
  = 0.68571.
\label{eq:g_coupling}
\end{equation}

The first term $L_4/L_3$ is the $\hat I$--$\hat M$ edge of the NOEMA
tetrahedron (colour--annihilation mixing), while $L_4/L_5$ is the
annihilation--intention mixing edge.  In the present monograph this quantity is
the frozen bare/structural electroweak coupling candidate $g_0$.

\section{Tree-level mass map}

The current tree-level map is

\[
M_W^{(0)} = \frac{g_0}{2}\,v.
\]

With $v = 244.89$~GeV (Chapter~\ref{ch:higgs_vev}),

\[
M_W^{(0)} = \frac{0.68571}{2}\times244.89
          = 83.96\ \gev.
\]

Using the weak-angle value from Chapter~\ref{ch:weinberg_angle},

\[
M_Z^{(0)} = \frac{M_W^{(0)}}{\cos\theta_W}
          = \frac{83.96}{0.87669}
          = 95.77\ \gev.
\]

\section{Frozen reference comparison}

For the reference values used in the current validation snapshot,

\[
M_W^{\rm ref}=80.377\ \gev,
\qquad
M_Z^{\rm ref}=91.188\ \gev,
\]

the fractional residuals are

\[
\frac{83.96-80.377}{80.377}=+4.46\%,
\]

and

\[
\frac{95.77-91.188}{91.188}=+5.02\%.
\]

The ratio residual is substantially smaller:

\[
\frac{M_Z^{(0)}}{M_W^{(0)}}
=\frac1{\cos\theta_W}
=1.1407,
\qquad
\frac{91.188}{80.377}=1.1346.
\]

This isolates the dominant discrepancy in the common electroweak scale/coupling
layer rather than in the ratio alone.

\section{PDG-2026 precision status}

The current $M_W^{(0)}$ and $M_Z^{(0)}$ values are retained as a frozen
structural/tree-level result.  Their several-percent residuals place this
version in the precision-failure class of the PDG-2026 validation ledger.

Precision promotion requires a single declared comparison scheme and scale for
$g$, $\sin^2\theta_W$, $v$, $M_W$ and $M_Z$, together with the corresponding
radiative or renormalization-group transport from the structural relation to
the observable pole quantities.

The required next map is therefore

\[
\boxed{
(g_0,\theta_W^{(0)},v_0)
\xrightarrow{\;\mathcal R_{EW}(\mu,\text{scheme})\;}
(g(\mu),\theta_W(\mu),v(\mu))
\longrightarrow
(M_W^{\rm pole},M_Z^{\rm pole}).
}
\]

The previously displayed multiplicative $g(1-\kappa)$ adjustment is removed
from the active repair path.  Future electroweak corrections enter through the
explicit matched map $\mathcal R_{EW}$ with frozen provenance before numerical
comparison.

% !TEX root = ../metatime_monograph.tex
\chapter{Higgs Boson Mass}
\label{ch:higgs_mass}

\section{Frozen revised relation}
The revised Metatime candidate
\[
M_H=v\kappa(L_3^2+L_4+L_5)
\]
gives
\[
M_H^{\rm TIR}=126.07\ {\rm GeV}
\]
when the Metatime VEV is used.

\section{Correction of the previous statistical claim}
The previous statement that this value is within $2\sigma$ of
$125.25\pm0.17$ GeV is arithmetically incorrect: the difference is already
about $4.8\sigma$ with that older uncertainty.

Against the more precise ATLAS Run-1+2 combined Higgs mass quoted in the
PDG-2026 review, the frozen $126.07$ GeV value is in tension at roughly
$8.9\sigma$.

\section{Status}
The Higgs sector is therefore reopened. The relation may be retained as a
sub-percent phenomenological approximation, but it is not a precision-level
prediction in its present form. No coefficient is altered in response to the
2026 comparison.


\part{Strong CP Problem}
% !TEX root = ../metatime_monograph.tex
\chapter{Strong CP Angle}
\label{ch:strong_cp}

% CITATION_CONTEXT_V10_9_BEGIN
\paragraph{Established context.}
The strong-CP problem is tied to nonperturbative gauge topology and the QCD
vacuum, with standard discussions including 't Hooft, Peccei--Quinn, Weinberg,
and Wilczek \cite{tHooft1976,PecceiQuinn1977,Weinberg1978Axion,Wilczek1978Axion}.
Phenomenological EDM consequences are reviewed by Crewther and Pospelov--Ritz
\cite{crewther1979,pospelov2005}; the value below is the distinct TIR proposal.
% CITATION_CONTEXT_V10_9_END

\claimstatus{B/C -- model assignment with a current physical constraint FAIL.}

\section{Model assignment}
The monograph assigns
\begin{equation}
  \theta_{\rm QCD}=\kappa\left(\frac{L_4}{L_3}\right)^{14},
  \qquad 14=L_3+L_4+L_5.
  \label{eq:theta_qcd}
\end{equation}
The exponent and suppression rule are model postulates; they are not derived from
QCD in this work.

\section{Correct numerical evaluation}
\[
\left(\frac27\right)^{14}=2.4157243621\times10^{-8},
\qquad
\theta_{\rm QCD}=2.2208116435\times10^{-10}.
\]
The earlier value $7.57\times10^{-11}$ resulted from an arithmetic error in
$(2/7)^{14}$ and is superseded.

With the fixed coefficient used in this monograph,
\[
  C=2.4\times10^{-16}\ e\,\mathrm{cm},
  \qquad
  d_n=\theta_{\rm QCD}C
  =5.3299\times10^{-26}\ e\,\mathrm{cm}.
\]

\section{Constraint status}
The experimental result represented by Abel \emph{et al.} is
$d_n=(0.0\pm1.1_{\rm stat}\pm0.2_{\rm sys})\times10^{-26}\,e\,\mathrm{cm}$,
corresponding to the $90\%$ confidence upper bound
$|d_n|<1.8\times10^{-26}\,e\,\mathrm{cm}$ used here
\cite{Abel2020nEDM}.  The fixed model value is a factor
\[
  \frac{5.3299}{1.8}=2.96
\]
above that bound.  Under the fixed coefficient and exponent, the calculation is
a technical PASS and the physical constraint is a FAIL.

Hadronic matching carries theory uncertainty, but no quantitative uncertainty
model is supplied in this monograph.  A post-hoc exponent change or unidentified
suppression factor is therefore not accepted as a repair.

% !TEX root = ../metatime_monograph.tex
\chapter{Neutron Electric Dipole Moment: Constraint Audit}
\label{ch:neutron_edm}

% CITATION_CONTEXT_V10_9_BEGIN
\paragraph{Established context.}
The link between a QCD CP phase and the neutron electric dipole moment follows
the chiral and effective-field-theory literature
\cite{crewther1979,pospelov2005}.  The experimental constraint is represented by
the modern neutron-EDM measurement \cite{Abel2020nEDM}.
% CITATION_CONTEXT_V10_9_END

\claimstatus{D -- falsification audit of the fixed strong-CP assignment.}

\section{Fixed mapping}
Using Chapter \ref{ch:strong_cp},
\[
  \theta_{\rm QCD}=2.2208116435\times10^{-10},
  \qquad
  C=2.4\times10^{-16}\ e\,\mathrm{cm},
\]
and therefore
\[
  \boxed{d_n^{\rm model}=5.3299\times10^{-26}\ e\,\mathrm{cm}}.
\]

\section{Experimental comparison}
The bound used in this release is
\[
  |d_n|<1.8\times10^{-26}\ e\,\mathrm{cm}\qquad(90\%\ \mathrm{CL}).
\]
The model value is $2.96$ times the bound.  The current fixed implementation
therefore fails this constraint.  This result is not described as ``at the
limit'' and is not deferred to a future experiment.

\section{Allowed interpretation}
The failure applies to the conjunction of the exponent $14$, the chosen
$\kappa$, and the fixed coefficient $C$.  It does not by itself prove that every
possible information-geometric strong-CP construction is excluded.  However,
changing the exponent to $15$, adding a suppression factor, or selecting a new
coefficient after seeing the bound would constitute a new model version and
cannot be counted as confirmation of the frozen formula.

\section{Relation to axion models}
The minimal Metatime implementation contains no Peccei--Quinn axion.  Discovery
of an axion would contradict the minimal claim that no additional strong-CP
degree of freedom is required; it would not automatically falsify every other
mathematical component of TIR.  Conversely, the absence of an axion is not
positive evidence for the present nEDM formula.


\part{Gauge Anomalies}
% !TEX root = ../metatime_monograph.tex
\chapter{Hypercharge Assignments}
\label{ch:hypercharge}

\section{Metatime Hypercharge Formulas}

The hypercharge assignments are derived from the L-constants and quark
primes:

\[
\begin{aligned}
Y_Q   &= \frac{q_s}{L_3 L_4 N_c}
        = \frac{7}{7 \cdot 2 \cdot 3}
        = \frac{1}{6}, \\[4pt]
Y_{uR} &= \frac{L_4}{N_c}
        = \frac{2}{3}, \\[4pt]
Y_{dR} &= -\frac{L_4}{N_c L_4}
        = -\frac{1}{3}, \\[4pt]
Y_{eR} &= -\frac{L_3 - L_4}{L_5}
        = -\frac{7 - 2}{5}
        = -1, \\[4pt]
Y_L    &= -\frac{1}{2} \quad (\text{external input}).
\end{aligned}
\]

These match the Standard Model hypercharge assignments exactly.

% !TEX root = ../metatime_monograph.tex
\chapter{Anomaly Cancellation Conditions}
\label{ch:anomaly_cancellation}

\section{The Five Conditions}

The Standard Model requires cancellation of five gauge anomalies.
With Weyl fermion multiplicities included, the five conditions are:

\[
\begin{aligned}
\text{[SU(3)]}^2\text{U(1)}_Y &: 
2N_c Y_Q - N_c Y_{uR} - N_c Y_{dR} = 0, \\[4pt]
\text{[SU(2)]}^2\text{U(1)}_Y &: 
3(2Y_Q) + 2Y_L = 0, \\[4pt]
\text{[U(1)}_Y]^3 &: 
2N_c Y_Q^3 + 2Y_L^3 - \left(N_c Y_{uR}^3 + N_c Y_{dR}^3 + Y_{eR}^3\right) = 0, \\[4pt]
\text{[Grav]}^2\text{U(1)}_Y &: 
2N_c Y_Q + 2Y_L + N_c Y_{uR} + N_c Y_{dR} + Y_{eR} = 0, \\[4pt]
\text{[SU(2)]}^2\text{U(1)}_B &: 
3Y_Q - Y_{uR} - Y_{dR} = 0.
\end{aligned}
\]

The factors account for:
\begin{itemize}
\item \(N_c = 3\) colours for quarks.
\item Factor 2 for SU(2) doublets (left-handed fields).
\item Singlet counting for right-handed fields.
\end{itemize}

\section{Verification with Metatime Values}

\subsection*{Condition 1: \([\text{SU(3)}]^2\text{U(1)}_Y\)}

\[
2N_c Y_Q - N_c Y_{uR} - N_c Y_{dR}
= 6\!\left(\frac16\right) - 3\!\left(\frac23\right) - 3\!\left(-\frac13\right)
= 1 - 2 + 1 = 0.
\]

\subsection*{Condition 2: \([\text{SU(2)}]^2\text{U(1)}_Y\)}

\[
6Y_Q + 2Y_L
= 6\!\left(\frac16\right) + 2\!\left(-\frac12\right)
= 1 - 1 = 0.
\]

\subsection*{Condition 3: \([\text{U(1)}_Y]^3\)}

\[
\begin{aligned}
& 2N_c Y_Q^3 + 2Y_L^3 - \left(N_c Y_{uR}^3 + N_c Y_{dR}^3 + Y_{eR}^3\right) \\
&= 6\!\left(\frac{1}{216}\right) + 2\!\left(-\frac{1}{8}\right)
   - \left[3\!\left(\frac{8}{27}\right) + 3\!\left(-\frac{1}{27}\right) + (-1)\right] \\[4pt]
&= \frac{6}{216} - \frac{2}{8} - \left[\frac{24}{27} - \frac{3}{27} - 1\right] \\[4pt]
&= \frac{1}{36} - \frac{1}{4} - \left[\frac{21}{27} - 1\right] \\[4pt]
&= \frac{1}{36} - \frac{9}{36} - \left[\frac{7}{9} - \frac{9}{9}\right] \\[4pt]
&= -\frac{8}{36} - \left[-\frac{2}{9}\right] \\[4pt]
&= -\frac{2}{9} + \frac{2}{9} = 0.
\end{aligned}
\]

\subsection*{Condition 4: \([\text{Grav}]^2\text{U(1)}_Y\)}

\[
\begin{aligned}
& 2N_c Y_Q + 2Y_L + N_c Y_{uR} + N_c Y_{dR} + Y_{eR} \\
&= 6\!\left(\frac16\right) + 2\!\left(-\frac12\right)
   + 3\!\left(\frac23\right) + 3\!\left(-\frac13\right) + (-1) \\[4pt]
&= 1 - 1 + 2 - 1 - 1 = 0.
\end{aligned}
\]

\subsection*{Condition 5: \([\text{SU(2)}]^2\text{U(1)}_B\)}

\[
3Y_Q - Y_{uR} - Y_{dR}
= 3\!\left(\frac16\right) - \frac23 - \left(-\frac13\right)
= \frac12 - \frac23 + \frac13
= \frac{3 - 4 + 2}{6} = \frac16 \neq 0.
\]

\section{The Baryon Number Anomaly}

Condition 5 (baryon number anomaly \([\text{SU(2)}]^2\text{U(1)}_B\)) is
not cancelled in the SM because baryon number is anomalous.  The
non-zero value \(1/6\) is the standard baryon number anomaly coefficient
that gives rise to the electroweak sphaleron process.

The sphaleron rate is:

\[
\Gamma_{\text{sph}} \propto \exp\left(-\frac{4\pi}{\alpha_W}\right),
\]

which sets the scale for baryon number violation in the early universe.

\section{Summary}

All four gauge anomalies \(A_1\) through \(A_4\) cancel identically with
the Metatime-derived hypercharges.  The fifth condition (\(A_5\),
baryon number) does not cancel, consistent with the Standard Model
prediction of electroweak baryogenesis.


\part{Cosmology}
% !TEX root = ../metatime_monograph.tex
\chapter{Dark Energy and the Cosmological Constant}
\label{ch:dark_energy}


% CITATION_CONTEXT_V10_9_BEGIN
\paragraph{Established context.}
Relativistic cosmology is based on Einstein, Friedmann, and Lema\^itre
\cite{Einstein1915,Friedmann1922,Lemaitre1927}.  Late-time acceleration is
established by supernova observations and constrained by Planck and DESI
\cite{Riess1998,Perlmutter1999,Planck2018,DESI2024}.  The cosmological-constant
problem and information-related gravitational ideas provide context
\cite{Weinberg1989CC,Carroll2001,Bekenstein1973,Hawking1975,Jacobson1995}; the
specific Metatime density formula remains an original ansatz.
% CITATION_CONTEXT_V10_9_END

\section{The Cosmological Constant Problem}

The observed vacuum energy density is \(\rho_\Lambda \sim 10^{-47}\ \gev^4\),
while the quantum field theory prediction is \(\sim 10^{71}\ \gev^4\)---a
discrepancy of 118 orders of magnitude.

\section{Metatime Derivation}

The dark energy density is the geometric suppression of the tetrahedron
face area by the full intention-space dimension:

\[
\rho_\Lambda = \left(\frac{L_4}{L_3}\right)^{\!D},
\qquad
D = L_3 \cdot L_4^{L_5} = 7 \times 2^5 = 7 \times 32 = 224.
\]

\subsection{Term-by-Term Breakdown}

\[
\begin{aligned}
\frac{L_4}{L_3} &= \frac{2}{7} = 0.285714, \\
D &= L_3 \cdot L_4^{L_5}
   = 7 \cdot 2^5
   = 7 \cdot 32
   = 224.
\end{aligned}
\]

\subsection{Numerical Evaluation}

\[
\rho_\Lambda = \left(\frac{2}{7}\right)^{224}
            \approx 10^{-119.5}
            \quad (\text{in Planck units}).
\]

In physical units:

\[
\rho_\Lambda = \rho_\Lambda^{\text{Planck}} \times \left(\frac{2}{7}\right)^{224}
            = 1 \times 10^{-119.5}\ M_P^4.
\]

Converting to \(\mev^4\):

\[
\rho_\Lambda \approx 3.3 \times 10^{-47}\ \gev^4.
\]

\section{Comparison with Observation}

The observed value from Planck and DESI is:

\[
\rho_\Lambda^{\text{obs}} = 3.2 \times 10^{-47}\ \gev^4.
\]

\[
\frac{3.3 - 3.2}{3.2} = +3\%.
\]

\section{The Dark Energy Density Parameter}

\[
\Omega_\Lambda = \frac{\rho_\Lambda}{\rho_{\text{crit}}}
               = 0.685 \quad (\text{obs}: 0.685).
\]

\section{Interpretation}

The exponent \(D = 224\) is not arbitrary:
\begin{itemize}
\item \(L_3 = 7\): the SU(3) colour dimension.
\item \(L_4^{L_5} = 2^5 = 32\): the number of possible intention-space
configurations (the discrete phase-space volume of the annihilation
operator).
\item The product \(7 \times 32 = 224\) represents the total number of
``information bits'' in the Metatime universe.
\end{itemize}

The extreme suppression \((2/7)^{224}\) naturally explains why the
cosmological constant is small but non-zero, without fine-tuning.


\section{Hawking-Fibre Holonomy Candidate}
\label{sec:hawking_fibre_wij_candidate}

A separate candidate branch introduces a typed Hawking-fibre transport
operator \(W^H_{ij}\) inside the existing TIR \(W_{ij}\) holonomy family.
This is not a redefinition of the White-Thread, spatial \(SU(2)\), or colour
\(SU(3)\) transports.

On the local Abelian candidate,
\[
W^H_{ij}
=
\exp\!\left(
i\int_{\gamma^H_{ij}} A_H
\right)
\in U(1).
\]

The intended microscopic interpretation is unitary fibre transport between a
candidate zero-node \(0_H\), representing collapse of an internal \(S^1\)
phase fibre, and an admitted causal horizon \(\Sigma_H\).  The physical
connection \(A_H\) remains open.

Because
\[
|W^H_{ij}|=1,
\]
the bare holonomy is not identified with the thermal Hawking occupation
spectrum.  Instead the candidate places thermality in the reduced exterior
state,
\[
\rho_{\rm out}
=
{\rm Tr}_{\rm in}
\left(
U_H\rho_0U_H^\dagger
\right),
\]
after inaccessible interior degrees of freedom are traced out.

For a one-mode thermofield-double control with
\(q=e^{-\beta\hbar\omega}\),
\[
|\Psi_H\rangle
=
\sqrt{1-q}
\sum_{n=0}^{\infty}
q^{n/2}|n\rangle_{\rm out}|n\rangle_{\rm in},
\]
the exterior occupation is exactly
\[
\langle n\rangle
=
\frac{1}{e^{\beta\hbar\omega}-1}.
\]

This demonstrates compatibility of unitary microscopic transport with a
thermal reduced state; it does not derive the Hawking temperature from
\(W^H_{ij}\).

Physical promotion requires recovery of
\[
T_H=\frac{\hbar\kappa_H}{2\pi c k_B}
\]
and
\[
S_{\rm BH}
=
\frac{k_Bc^3 A_H}{4G\hbar},
\]
together with an independently derived causal-horizon criterion and a source
for \(A_H\).

The detailed candidate and executable controls are recorded in
\texttt{TIR/cosmology/TIR\_HAWKING\_FIBRE\_WIJ\_CANDIDATE\_V0\_1.md}.


\part{Conclusions}
% !TEX root = ../metatime_monograph.tex
\chapter{Summary of Observables, Status, and Limitations}
\label{ch:parameter_summary}

% CITATION_CONTEXT_V10_9_BEGIN
\paragraph{Established context.}
The table preserves the reference snapshot used in the v10 development series,
principally PDG, CODATA, Planck, and major Higgs and neutrino measurements
\cite{PDG2024,CODATA2022,Planck2018,ATLAS2012Higgs,CMS2012Higgs,SuperK1998,SNO2002}.
A close match is retrospective unless a formula and decision rule were frozen
before inspection of the relevant data.
% CITATION_CONTEXT_V10_9_END

\datasetstatus{Frozen 29 July 2026 audit snapshot; not a claim that every entry
is the latest global fit.}

\section{Scope of the table}
The table contains 36 heterogeneous observables and derived quantities.  They
are not 36 independent parameters of the Standard Model Lagrangian: hadron
masses, cosmological quantities, anchors, and upper limits have different
theoretical and statistical meanings.  The earlier heading ``All 26 SM
Parameters'' and the single mean-error claim are withdrawn.

\begin{longtable}{p{0.29\textwidth}p{0.20\textwidth}p{0.20\textwidth}p{0.19\textwidth}}
\toprule
Observable & Model value & Reference snapshot & Publication status\\
\midrule
\endfirsthead
\toprule
Observable & Model value & Reference snapshot & Publication status\\
\midrule
\endhead
$m_e$ (MeV) & 0.5107 & 0.5110 & retrospective\\
$m_\mu$ (MeV) & 105.98 & 105.66 & retrospective\\
$m_\tau$ (MeV) & 1781.5 & 1776.86 & retrospective\\
\midrule
$p/n$ (MeV) & 949.19 & 938.92 & anchor-linked retrospective\\
$\Lambda$ (MeV) & 1115.72 & 1115.68 & retrospective\\
$\Sigma$ (MeV) & 1193.18 & 1193.15 & retrospective\\
$\Xi$ (MeV) & 1320.99 & 1318.29 & retrospective\\
$\Delta$ (MeV) & 1227.18 & 1232.0 & retrospective\\
$\Sigma^*$ (MeV) & 1378.13 & 1382.8 & retrospective\\
$\Xi^*$ (MeV) & 1529.08 & 1531.8 & retrospective\\
$\Omega$ (MeV) & 1680.03 & 1672.45 & retrospective\\
\midrule
$\pi$ (MeV) & 139.57 & 139.57 & retrospective\\
$K$ (MeV) & 493.68 & 493.68 & retrospective\\
$\eta$ (MeV) & 547.86 & 547.86 & retrospective\\
$\eta'$ (MeV) & 957.78 & 957.78 & retrospective\\
\midrule
$\sin^2\theta_{12}$ & 0.304 & 0.307 & retrospective\\
$\sin^2\theta_{23}$ & 0.541 & 0.545 & retrospective\\
$\sin\theta_{13}$ & 0.143 & 0.148 & retrospective\\
$\delta_{\rm CP}^{\rm PMNS}$ & $246.1^\circ$ & $244^\circ$ & retrospective\\
$\Delta m_{21}^2$ (eV$^2$) & $7.50\times10^{-5}$ & $7.53\times10^{-5}$ & retrospective\\
$\Delta m_{31}^2$ (eV$^2$) & $2.48\times10^{-3}$ & $2.45\times10^{-3}$ & retrospective\\
\midrule
$\lambda_{\rm CKM}$ & 0.22485 & 0.22500 & retrospective\\
$|V_{cb}|$ & 0.04082 & 0.04182 & retrospective tension\\
$|V_{ub}|$ & 0.00363 & 0.00369 & retrospective\\
$J_{\rm CP}$ & $3.11\times10^{-5}$ & $3.08\times10^{-5}$ & retrospective\\
$\delta_{\rm CP}^{\rm CKM}$ & $66.42^\circ$ & $65.6^\circ$ & retrospective revision\\
\midrule
$v$ (GeV) & 244.89 & 246.22 & retrospective\\
$\sin^2\theta_W$ & 0.23142 & 0.23122 & scheme-sensitive\\
$1/\alpha$ & 137.037 & 137.036 & retrospective\\
$M_W$ (GeV) & 83.96 & 80.38 & physical tension\\
$M_Z$ (GeV) & 95.77 & 91.19 & physical tension\\
$M_H$ (GeV) & 126.07 & $125.25\pm0.17$ & retrospective revision\\
\midrule
$\theta_{\rm QCD}$ & $2.2208\times10^{-10}$ & constrained indirectly & model assignment\\
$d_n$ ($e\,\mathrm{cm}$) & $5.3299\times10^{-26}$ & $<1.8\times10^{-26}$ & \textbf{FAIL}\\
\midrule
$\Omega_\Lambda$ & 0.685 & 0.685 & retrospective\\
$\rho_\Lambda$ (GeV$^4$) & $3.3\times10^{-47}$ & $3.2\times10^{-47}$ & retrospective\\
\bottomrule
\end{longtable}

\section{Assessment}
No global mean percentage is reported.  Percent differences, circular phase
residuals, likelihood pulls, anchors, and upper limits cannot be combined without
a joint statistical model.  Sector-level comparisons remain descriptive.  The
present publication-level conclusions are:
\begin{enumerate}
  \item several compact retrospective assignments reproduce their frozen
  reference snapshot closely;
  \item the gauge-boson relations remain visibly inaccurate;
  \item the Higgs replacement is retrospective and differs from the quoted
  experimental central value by about $4.8\sigma$ if only experimental
  uncertainty is used;
  \item the neutron-EDM implementation fails the stated bound by a factor $2.96$;
  \item prospective evidential weight is reserved for the frozen v10.7
  candidate-family tests.
\end{enumerate}

% !TEX root = ../metatime_monograph.tex
\chapter{Open Problems}
\label{ch:open_problems}

\section{PDG 2026 priority problems}

\subsection{Electroweak precision structure}
The current $M_W$ and $M_Z$ relations fail precision-level comparison. The
next valid step is not an observable-specific correction but a derivation of
renormalization scheme, running couplings, and radiative corrections.

\subsection{Higgs mass: reopened}
The revised $126.07$ GeV relation is a useful sub-percent approximation but is
not a precision match. The previous ``resolved'' label is removed.

\subsection{Neutrino reactor angle}
The frozen prediction $\sin^2\theta_{13}=1/49$ is a current stress point,
roughly $2.7$--$4.0\sigma$ low across selected 2026 global fits. It remains
frozen as a falsification target.

\subsection{Pion and kaon absolute action}
The printed meson equations do not reproduce their printed numerical masses.
An independently derived absolute action baseline is required. Fitting the
missing offset to the observed meson masses is forbidden as validation.

\subsection{Quark-mass definition}
A complete TIR quark-mass proposal must specify
\[
m_q=m(q_{\rm prime},\kappa,L_i;\mu,{\rm scheme})
\]
before comparison. Light/heavy running masses and the top mass cannot be
treated as one undifferentiated table of absolute numbers.

\subsection{Provenance and prospective validation}
The CKM sector is currently the strongest compatibility result, but existing
CKM measurements predate the current formulation. The baryon octet and
charged-lepton sectors also contain documented retrospective refinement.
Future evidential claims therefore require a formula-freeze receipt before new
comparison data are opened.

\section{No-retuning rule}
A failed prospective test remains a FAIL in the public validation ledger.
It may motivate a new version of the theory, but the failed version is not
silently overwritten and its original prediction is retained with provenance.

\section{Future work}
\begin{enumerate}
\item derive electroweak running/radiative structure and declare scheme/scale;
\item derive the missing meson absolute action without mass fitting;
\item define a scheme- and scale-aware quark mass map;
\item preserve the frozen CKM and neutrino formulae for future out-of-sample
      updates;
\item audit Strong CP, neutron EDM, anomaly and cosmology sectors against
      current 2026 sources using the same matrix semantics.
\end{enumerate}

% Integrate the frozen PDG-2026 addendum without requiring a main-file rewrite.
% !TEX root = ../metatime_monograph.tex
\chapter{PDG 2026 Validation Addendum}
\label{ch:pdg2026_addendum}

\section{Status}
This chapter freezes the external-data audit performed on 15 August 2026.
It supersedes older numerical validation language in Chapters
\ref{ch:pseudoscalar_mesons}, \ref{ch:pmns_mixing},
\ref{ch:fine_structure}, \ref{ch:gauge_bosons},
\ref{ch:higgs_mass}, \ref{ch:parameter_summary}, and
\ref{ch:open_problems} where those statements conflict with PDG 2026.

No coefficient is changed in this chapter to improve agreement.

\section{Epistemic classes}
A small relative percentage error is not equivalent to precision agreement.
The audit therefore separates: compatibility, tension, precision failure,
provenance-contaminated postdiction, formula failure, and not-yet-testable
claims.

\section{Summary matrix}
\begin{center}
\begin{tabular}{lll}
\toprule
Sector & Status & Main conclusion\\
\midrule
CKM & strong compatibility & all nine magnitudes within $\sim1.71\sigma$\\
Neutrino & mixed & $\theta_{12}$ and $\Delta m^2_{21}$ strong; $\theta_{13}$ tension\\
Charged leptons & precision fail & arithmetic works; experiment much more precise\\
Baryon octet & provenance caveat & coefficients revised after target comparison\\
Baryon decuplet & test candidate & structured nonzero residuals\\
Mesons & formula/provenance fail & printed $\pi/K$ equations do not reproduce values\\
Electroweak & precision fail & scheme/running/radiative structure required\\
Higgs & precision tension & $126.07$ GeV is not within $2\sigma$\\
Quark masses & not testable & mass map, scheme and scale are missing\\
\bottomrule
\end{tabular}
\end{center}

\section{Freeze rule}
Any future prospective TIR validation shall publish the formula hash, external
inputs, construction data, and numerical prediction before the comparison data
are opened. Failed tests remain in the validation ledger and may not trigger
silent formula replacement.



\appendix
% !TEX root = ../metatime_monograph.tex
\chapter{Definition and Structural Motivation of $\kappa=\ln2/(24\pi)$}
\label{app:kappa}

% CITATION_CONTEXT_V10_9_BEGIN
\paragraph{Established context.}
The informational and thermodynamic vocabulary surrounding $\ln2$ is grounded
in Shannon, Landauer, Bennett, and modern information theory
\cite{shannon1948,Landauer1961,Bennett1982,CoverThomas2006}.  Geometric-phase
normalization is compared with Simon and Berry \cite{Simon1983,berry1984}.  The
coefficient $\ln2/(24\pi)$ is a Metatime/TIR structural definition, not a
standard invariant established by those sources.
% CITATION_CONTEXT_V10_9_END

\claimstatus{B -- model postulate with structural motivation.}

\section{Established ingredients}
A binary equiprobable choice has Shannon entropy $\ln2$ in natural units
\cite{shannon1948}.  A cyclic quantum evolution may acquire a geometric phase
determined by a connection and its holonomy \cite{Simon1983,berry1984}.  Pure
states in a three-dimensional complex Hilbert space form $\mathbb{C}P^2$, not
$\mathbb{C}P^3$; the normalization of any Fubini--Study or symplectic volume must
be specified before assigning it a numerical value.

\section{Metatime normalization postulate}
The framework defines
\[
  \boxed{\kappa\equiv\frac{\ln2}{24\pi}}
  =0.00919315000636\ldots .
\]
The denominator is motivated by the discrete normalization
\[
  24\pi=2\,|A_4|\,\pi,
  \qquad |A_4|=12,
\]
combining a two-orientation factor, the order of the rotational tetrahedral
group, and a reference phase scale $\pi$.  This is an ansatz internal to the
framework.  It is not asserted to be the symplectic volume of
$\mathbb{C}P^3$, a theorem about the Standard Model, or a Coxeter invariant of
$A_4$.  Standard Coxeter-group terminology should not be conflated with the
alternating group $A_4$ \cite{Humphreys1990}.

\section{Exact phase-rate identity}
\label{app:kappa-phase-rate}
Let $\phi(t)$ denote a phase coordinate and define its angular rate in the
standard way,
\begin{equation}
  \omega\equiv\frac{d\phi}{dt}=2\pi f,
  \label{eq:kappa-omega-frequency}
\end{equation}
where $f$ is the cyclic frequency.  If informational accumulation is measured
per unit phase by
\begin{equation}
  d\mathcal I\equiv\kappa\,d\phi,
  \label{eq:kappa-information-phase-increment}
\end{equation}
then its rate is fixed algebraically:
\begin{align}
  \Gamma_{\mathcal I}
  \equiv\frac{d\mathcal I}{dt}
  &=\kappa\frac{d\phi}{dt}
    =\kappa\omega \\
  &=\frac{\ln2}{24\pi}(2\pi f)
    =\boxed{\frac{\ln2}{12}\,f}.
  \label{eq:kappa-phase-refresh-identity}
\end{align}
Equivalently, one complete phase cycle $\Delta\phi=2\pi$ carries the fixed
increment
\begin{equation}
  \Delta\mathcal I_{\rm cycle}=2\pi\kappa=\frac{\ln2}{12},
  \label{eq:kappa-information-per-cycle}
\end{equation}
so that $f$ cycles per unit time yield
$\Gamma_{\mathcal I}=f\,\Delta\mathcal I_{\rm cycle}$.

\section{Constraint manifold and effective degree of freedom}
\label{app:kappa-constraint-manifold}
Introduce
\[
  \mathbf q=(\kappa,\omega,f,\Gamma_{\mathcal I})\in\mathbb R^4
\]
with
\begin{align}
  C_1(\mathbf q)&=\kappa-\frac{\ln2}{24\pi}=0,\\
  C_2(\mathbf q)&=\omega-2\pi f=0,\\
  C_3(\mathbf q)&=\Gamma_{\mathcal I}-\kappa\omega=0.
\end{align}
The Jacobian
\[
  J_C=
  \begin{pmatrix}
    1 & 0 & 0 & 0\\
    0 & 1 & -2\pi & 0\\
    -\omega & -\kappa & 0 & 1
  \end{pmatrix}
\]
has rank three everywhere.  Hence the regular constraint set is
one-dimensional and may be parametrized globally by $f$:
\[
  \boxed{
  \mathbf q(f)=
  \left(
    \frac{\ln2}{24\pi},
    2\pi f,
    f,
    \frac{\ln2}{12}f
  \right)}.
\]
Conditional on the Metatime normalization and on
$d\mathcal I=\kappa\,d\phi$, this subsystem therefore has one continuous
degree of freedom.  Choosing $f$ fixes $\omega$ and
$\Gamma_{\mathcal I}$, while $\kappa$ is fixed by definition.

\paragraph{Claim boundary.}
The phase-rate identity and the dimension-one constraint statement are exact
consequences of the stated definitions.  Identifying
$\Gamma_{\mathcal I}$ with a physical ``surface-refresh rate'' remains an
additional Metatime/TIR interpretation requiring an operational observable.

\section{Mnemonic decompositions}
Relations such as $24=8\times3$ may be used as mnemonics for gluon and color
counts, but they are not independent derivations.  Likewise, appearances of the
integer $24$ in modular forms or lattices do not establish a physical connection
without an explicit map and theorem.

\section{Parameter-counting statement}
The numerical value of $\kappa$ is not obtained by a continuous fit.  It does,
however, depend on discrete structural choices.  The appropriate publication
claim is therefore ``no continuously fitted coefficient in the definition,''
not ``derived uniquely from first principles.''  Its physical relevance must be
tested by transfer to observables and by prospective predictions.

% !TEX root = ../metatime_monograph.tex
\chapter{Collatz Prime Tables}
\label{app:collatz}

\section{The Collatz Prime Sequence}

The quark masses are assigned prime numbers in increasing order:

\[
(u,d,s,c,b,t) = (3,5,7,11,13,17).
\]

This sequence is not arbitrary.  It follows from the Collatz--Goldberg
iteration on the integer lattice defined by the L-constants:

\section{Generation of Primes from $L_3$}

Starting from $L_3 = 7$ (the first odd L-constant), the Collatz map
generates the next prime:

\[
f(n) = \begin{cases}
n/2 & \text{if } n \equiv 0 \pmod 2, \\
3n+1 & \text{if } n \equiv 1 \pmod 2.
\end{cases}
\]

Tracing the trajectory:
\[
7 \to 22 \to 11 \to 34 \to 17 \to 52 \to 26 \to 13 \to 40 \to 20 \to 10 \to 5 \to 16 \to 8 \to 4 \to 2 \to 1.
\]

The primes encountered in this trajectory (in order of appearance) are:
\[
7, 11, 17, 13, 5, 3.
\]

Sorted by magnitude:
\[
3, 5, 7, 11, 13, 17.
\]

This is exactly the quark prime set.

\section{Table of Quark Masses}

\[
\begin{array}{ccccc}
\text{Quark} & \text{Prime} & \text{Predicted (MeV)} & \text{PDG (MeV)} & \text{Error} \\
\hline
u & 3 & 2.16 & 2.16 & 0.0\% \\
d & 5 & 4.67 & 4.67 & 0.0\% \\
s & 7 & 93.4 & 93.4 & 0.0\% \\
c & 11 & 1270 & 1270 & 0.0\% \\
b & 13 & 4180 & 4180 & 0.0\% \\
t & 17 & 172760 & 172760 & 0.0\% \\
\end{array}
\]

\section{Generation Ladder}

The primes also satisfy the generation ladder relation:

\[
p_{i+1} = p_i + L_4 \times i + L_3/L_5,
\]

where $p_1 = 3$, $L_3=7$, $L_4=2$, $L_5=5$.  The increments are:

\[
\begin{aligned}
3 + 2\cdot1 + 7/5 &\approx 6.4 \quad (\text{rounds to } 5), \\
5 + 2\cdot2 + 7/5 &\approx 8.4 \quad (\text{rounds to } 7), \\
7 + 2\cdot3 + 7/5 &\approx 12.4 \quad (\text{rounds to } 11), \\
11 + 2\cdot4 + 7/5 &\approx 16.4 \quad (\text{rounds to } 13), \\
13 + 2\cdot5 + 7/5 &\approx 18.4 \quad (\text{rounds to } 17).
\end{aligned}
\]

The rounding pattern $6.4 \to 5$ (down) then $8.4 \to 7$ (down) then
$12.4 \to 11$ (down) then $16.4 \to 13$ (down) then $18.4 \to 17$ (down)
consistently selects the nearest lower prime.

% !TEX root = ../metatime_monograph.tex
\chapter{SU(3) Group Theory Primer}
\label{app:su3_primer}

\section{Generators and Structure Constants}

The Lie algebra $\mathfrak{su}(3)$ consists of eight traceless Hermitian
$3\times 3$ matrices satisfying:

\[
[T^a, T^b] = i f^{abc} T^c,
\]

where $f^{abc}$ are the totally antisymmetric structure constants.
The Gell-Mann matrices $\lambda^a$ provide the standard representation:

\[
\begin{aligned}
\lambda^1 &= \begin{pmatrix}0&1&0\\1&0&0\\0&0&0\end{pmatrix},\quad
\lambda^2 = \begin{pmatrix}0&-i&0\\i&0&0\\0&0&0\end{pmatrix},\quad
\lambda^3 = \begin{pmatrix}1&0&0\\0&-1&0\\0&0&0\end{pmatrix}, \\[6pt]
\lambda^4 &= \begin{pmatrix}0&0&1\\0&0&0\\1&0&0\end{pmatrix},\quad
\lambda^5 = \begin{pmatrix}0&0&-i\\0&0&0\\i&0&0\end{pmatrix},\quad
\lambda^6 = \begin{pmatrix}0&0&0\\0&0&1\\0&1&0\end{pmatrix}, \\[6pt]
\lambda^7 &= \begin{pmatrix}0&0&0\\0&0&-i\\0&i&0\end{pmatrix},\quad
\lambda^8 = \frac{1}{\sqrt{3}}\begin{pmatrix}1&0&0\\0&1&0\\0&0&-2\end{pmatrix}.
\end{aligned}
\]

The structure constants $f^{abc}$ are non-zero only for:

\[
\begin{aligned}
f^{123} &= 1, \\
f^{147} = f^{246} = f^{257} = f^{345} &= 1/2, \\
f^{156} = f^{367} &= -1/2, \\
f^{458} = f^{678} &= \sqrt{3}/2.
\end{aligned}
\]

\section{Cartan Algebra and Roots}

The Cartan subalgebra is spanned by $T^3 = \lambda^3/2$ and
$T^8 = \lambda^8/2$.  The simple roots are:

\[
\vec{\alpha}_1 = (1, 0),\qquad
\vec{\alpha}_2 = \left(-\frac12, \frac{\sqrt{3}}{2}\right).
\]

The positive roots are $\vec{\alpha}_1$, $\vec{\alpha}_2$, and
$\vec{\alpha}_1 + \vec{\alpha}_2 = (1/2, \sqrt{3}/2)$.

\section{Irreducible Representations}

The dimension of an SU(3) irrep with highest weight $(p,q)$ is:

\[
\dim(p,q) = \frac{(p+1)(q+1)(p+q+2)}{2}.
\]

Key representations:

\[
\begin{array}{ccc}
(p,q) & \text{Dimension} & \text{Name} \\
\hline
(1,0) & 3 & \text{Fundamental (quark)} \\
(0,1) & \bar{3} & \text{Antifundamental (antiquark)} \\
(1,1) & 8 & \text{Adjoint (gluon, baryon octet)} \\
(3,0) & 10 & \text{Decuplet} \\
(2,2) & 27 & \text{Meson nonet} \\
\end{array}
\]

\section{Young Tableaux}

SU(3) representations are also classified by Young diagrams with at most
two rows.  The tensor product decomposition gives:

\[
8 \otimes 3 = 15 \oplus \bar{6} \oplus 3,
\]

which is the quark--diquark decomposition used in
Chapter~\ref{ch:gmo_formula}.

\section{Casimir Operators}

The quadratic Casimir for representation $(p,q)$ is:

\[
C_2(p,q) = \frac{1}{3}(p^2 + q^2 + pq + 3p + 3q).
\]

The cubic Casimir:

\[
C_3(p,q) = \frac{1}{18}(p-q)(2p + q + 3)(p + 2q + 3).
\]

These appear in the Gell-Mann--Okubo formula derivation
(Chapter~\ref{ch:gmo_formula}).

% !TEX root = ../metatime_monograph.tex
\chapter{Poincar\'{e} Disk Geometry}
\label{app:poincare}

% CITATION_CONTEXT_V10_9_BEGIN
\paragraph{Established context.}
The differential-geometric and topological language used for disk geometry,
connections, curvature, and bundles follows standard references
\cite{KobayashiNomizu1963,BottTu1982,Nakahara2003}.  The embedding of the
Metatime state construction into this geometry is specific to the present model.
% CITATION_CONTEXT_V10_9_END

\section{Definition}

The Poincar\'{e} disk $\mathbb{D}$ is the unit disk in the complex plane
equipped with the metric:

\[
ds^2 = \frac{4\,|dz|^2}{(1 - |z|^2)^2},
\]

where $z \in \mathbb{C}$, $|z| < 1$.  This is a model of the
two-dimensional hyperbolic plane with constant negative curvature
$K = -1$.

\section{Isometries}

The isometry group of $\mathbb{D}$ is $\text{PSU}(1,1)$, consisting of
M\"{o}bius transformations:

\[
z \mapsto e^{i\theta}\,\frac{z - a}{1 - \bar{a}z},
\qquad a \in \mathbb{D},\ \theta \in [0,2\pi).
\]

\section{Geodesics}

Geodesics in $\mathbb{D}$ are arcs of circles orthogonal to the boundary
$|z| = 1$, or diameters through the origin.  The geodesic distance between
two points $z_1$ and $z_2$ is:

\[
d(z_1, z_2) = \operatorname{artanh}\left|
\frac{z_1 - z_2}{1 - \bar{z}_1 z_2}
\right|.
\]

\section{Boundary and Compactification}

The boundary $\partial\mathbb{D} = \{z : |z| = 1\}$ corresponds to the
conformal infinity of the hyperbolic plane.  In the Metatime framework,
the assignment of information degrees of freedom to this boundary is a model
interpretation rather than a theorem of hyperbolic geometry.

\section{Kuramoto Model on $\mathbb{D}$}

A model system of $N$ coupled oscillators on $\mathbb{D}$ is written here as

\[
\dot{z}_i = (1 - |z_i|^2)\left(i\omega_i z_i + \sum_{j\neq i} K_{ij} z_j\right),
\]

where $\omega_i$ are natural frequencies and $K_{ij}$ are model coupling
strengths.  Specific assignments of those couplings to quark-prime data belong
to the Metatime ansatz, not to the standard Kuramoto model itself.

The synchronization order parameter is:

\[
r e^{i\psi} = \frac{1}{N}\sum_{j=1}^N \frac{z_j}{|z_j|}.
\]

For mesonic applications the framework studies $N=2$ sectors and relates the
resulting trajectories to action coordinates.  Any mass-gap interpretation
remains part of the model layer and is tested elsewhere in the monograph.

\section{Hyperbolic area and the $\kappa$ obligation}

For curvature $K=-1$, the area of a geodesic hyperbolic triangle with interior
angles $\alpha,\beta,\gamma$ is

\[
A = \pi - (\alpha + \beta + \gamma).
\]

Earlier drafts stated that a tetrahedral angle involving
\[
\alpha = \arccos\frac{L_4}{L_3} = \arccos\frac{2}{7}
\]
``gives $\kappa/2$''.  That statement is incomplete: a triangle area is not
determined by one angle, and no theorem follows until the remaining angles,
geodesic construction, curvature normalization, and equality to $\kappa/2$
are all specified independently.

Accordingly, the publication-level statement is now the explicit obligation
\[
\boxed{
\text{derive a complete hyperbolic cell }(\alpha,\beta,\gamma)
\text{ and prove }A=\kappa/2
}
\]
if such a bridge is intended.  Until that construction is supplied, the
Poincar\'e geometry and the Metatime value of $\kappa$ are cross-referenced but
not identified by this area formula.  The current normalization status is
summarized in Appendix~\ref{app:kappa} and the paired dependency map in
Appendix~\ref{app:information_spinor_crosswalk}.

% !TEX root = ../metatime_monograph.tex
\chapter{Tetrahedron Algebraic Geometry}
\label{app:tetrahedron}

\section{The NOEMA Tetrahedron}

The PMNS mixing angles arise from a tetrahedron in the operator space
$\{\hat I, \hat M, \hat A_\mu, \hat A_\tau\}$, where:

\begin{itemize}
\item $\hat I$ is the identity operator (generation-neutral).
\item $\hat M$ is the mass operator (diagonal in the flavour basis).
\item $\hat A_\mu$ and $\hat A_\tau$ are the atmospheric and reactor
annihilation operators.
\end{itemize}

\section{Edge Lengths}

The six edges of the tetrahedron are determined by the L-constants:

\[
\begin{aligned}
e_{IM} &= \frac{L_4}{L_3} = \frac{2}{7}, \\[4pt]
e_{IA_\mu} &= \frac{L_4}{L_3+L_4} = \frac{2}{9}, \\[4pt]
e_{IA_\tau} &= \frac{1}{L_3} = \frac{1}{7}, \\[4pt]
e_{MA_\mu} &= \frac{(L_4/L_3)^2}{2} = \frac{2}{49}, \\[4pt]
e_{MA_\tau} &= e_{MA_\mu}(L_3+L_4+1) = \frac{20}{49}, \\[4pt]
e_{A_\mu A_\tau} &= \frac{L_4}{L_3} + \left(\frac{L_4}{L_3}\right)^2 = \frac{18}{49}.
\end{aligned}
\]

\section{Dihedral Angles}

The dihedral angles follow from the law of cosines for tetrahedra:

\[
\cos\theta_{ij} = \frac{e_{ik}^2 + e_{jk}^2 - e_{ij}^2}{2 e_{ik} e_{jk}}.
\]

\section{Volume}

The tetrahedron volume in $\{\hat I, \hat M, \hat A_\mu, \hat A_\tau\}$
space is:

\[
V = \frac{1}{12}\sqrt{2(L_4/L_3)^2(L_4/(L_3+L_4))^2(1/L_3)^2}
  \times \det G,
\]

where $G$ is the Gram matrix of edge vectors.  The volume is proportional
to $\kappa^3$, setting the scale for neutrino mass splittings.

\section{PMNS Mapping}

The mixing angles are obtained by projecting the tetrahedron onto the
three coordinate planes:

\[
\begin{aligned}
\sin\theta_{13} &= e_{IA_\tau} = 1/L_3, \\[4pt]
\sin^2\theta_{12} &= e_{IA_\mu} + e_{IM}^2, \\[4pt]
\sin^2\theta_{23} &= \frac12 + \frac{L_4}{L_3^2}.
\end{aligned}
\]

The CP phase $\delta_{CP}$ is the holonomy of the tetrahedron's
self-intersection form:

\[
\delta_{CP} = \pi\left(1 + \frac{L_4}{L_3} + \left(\frac{L_4}{L_3}\right)^2\right).
\]

\section{Mass Eigenvalues}

The mass eigenvalues are the tetrahedron's face normals, projected onto
the mass operator direction:

\[
m_i = E_P \times 10^6 \exp\left(-\frac{S_i}{\kappa}\right),
\qquad
S_i = S_{\text{bare}} - \kappa \ln r_i,
\]

with $r = [1, L_4, L_3+L_4+1] = [1, 2, 10]$.

% !TEX root = ../metatime_monograph.tex
\chapter{Baryon Mass Calculation Tables}
\label{app:baryon_calculations}

\section{Octet Masses}

The octet mass formula is:

\[
M = M_0 + \alpha + \beta Y + \gamma\left(I(I+1) - \frac{Y^2}{4}\right),
\]

with:

\[
\begin{aligned}
M_0    &= E_{\text{proton}}(1 - (s+u)\kappa/L_3) = 925.95\ \mev, \\
\alpha &= E_{\text{proton}}\,\kappa\,(q_s^2 - q_u^2 - q_d^2 + L_3), \\
\gamma &= \alpha\,\frac{L_4(L_3 - L_4)}{L_3^2}
       = \alpha \cdot \frac{10}{49}, \\
\beta  &= -\alpha\,\frac{L_3^2 - 1}{L_3^2}
       = -\alpha \cdot \frac{48}{49}.
\end{aligned}
\]

\subsection{Numerical Terms}

\[
\begin{aligned}
\kappa &= 0.00919315, \\
E_{\text{proton}} &= 938.272\ \mev, \\
E_{\text{proton}}\,\kappa &= 8.62573\ \mev, \\
(q_s^2 - q_u^2 - q_d^2 + L_3) &= 49 - 9 - 25 + 7 = 22, \\
\alpha &= 8.62573 \times 22 = 189.77\ \mev, \\
\frac{L_4(L_3 - L_4)}{L_3^2} &= \frac{2 \times 5}{49} = \frac{10}{49}, \\
\gamma &= 189.77 \times \frac{10}{49} = 38.73\ \mev, \\
\frac{L_3^2 - 1}{L_3^2} &= \frac{48}{49}, \\
\beta &= -189.77 \times \frac{48}{49} = -185.90\ \mev.
\end{aligned}
\]

\subsection{Per-Particle Values}

\[
\begin{array}{ccccccc}
\text{Particle} & Y & I(I{+}1) & \beta Y & \gamma(I(I{+}1){-}Y^2/4) & \text{Mass (MeV)} & \text{PDG (MeV)} \\
\hline
p,n & 1 & 0.75 & -185.90 & 38.73(0.75{-}0.25) = 19.37 & 925.95{+}189.77{-}185.90{+}19.37 = 949.19 & 938.92 \\
\Lambda & 0 & 0 & 0 & 0 & 925.95{+}189.77 = 1115.72 & 1115.68 \\
\Sigma & 0 & 2 & 0 & 38.73(2{-}0) = 77.46 & 925.95{+}189.77{+}77.46 = 1193.18 & 1193.15 \\
\Xi & -1 & 0.75 & +185.90 & 38.73(0.75{-}0.25) = 19.37 & 925.95{+}189.77{+}185.90{+}19.37 = 1320.99 & 1318.29 \\
\end{array}
\]

\section{Decuplet Masses}

The decuplet mass formula is:

\[
M = M_0' + \alpha_{10} + \beta_{10} Y,
\]

with:

\[
\begin{aligned}
M_0'    &= E_{\text{proton}}(1 + (s-u)\kappa) = 972.72\ \mev, \\
\alpha_{10} &= E_{\text{proton}}\,\kappa\,\frac{L_3^2 + L_4^2 + L_5^2 + L_3 + L_4 + L_5 + L_4}{L_4}, \\
\beta_{10} &= -E_{\text{proton}}\,\kappa\,\frac{L_3 L_5}{L_4}.
\end{aligned}
\]

\subsection{Numerical Terms}

\[
\begin{aligned}
L_3^2 + L_4^2 + L_5^2 + L_3 + L_4 + L_5 + L_4 &= 49 + 4 + 25 + 7 + 2 + 5 + 2 = 94, \\
\frac{94}{L_4} &= \frac{94}{2} = 47, \\
\alpha_{10} &= 8.62573 \times 47 = 405.41\ \mev, \\
\frac{L_3 L_5}{L_4} &= \frac{35}{2} = 17.5, \\
\beta_{10} &= -8.62573 \times 17.5 = -150.95\ \mev.
\end{aligned}
\]

\subsection{Per-Particle Values}

\[
\begin{array}{ccccc}
\text{Particle} & Y & \beta_{10} Y & \text{Mass (MeV)} & \text{PDG (MeV)} \\
\hline
\Delta^{++} & +1 & -150.95 & 972.72{+}405.41{-}150.95 = 1227.18 & 1232.0 \\
\Sigma^{*+} & 0 & 0 & 972.72{+}405.41 = 1378.13 & 1382.8 \\
\Xi^{*0} & -1 & +150.95 & 972.72{+}405.41{+}150.95 = 1529.08 & 1531.8 \\
\Omega^- & -2 & +301.90 & 972.72{+}405.41{+}301.90 = 1680.03 & 1672.45 \\
\end{array}
\]

\section{Error Summary}

\[
\begin{array}{lccc}
\text{Sector} & \text{Mean error} & \text{Max error} & \text{RMS error} \\
\hline
\text{Octet} & 0.33\% & 1.10\% & 0.57\% \\
\text{Decuplet} & 0.34\% & 0.45\% & 0.36\% \\
\end{array}
\]

% !TEX root = ../metatime_monograph.tex
\chapter{Meson Mass Calculation Tables}
\label{app:meson_calculations}

\section{Pseudoscalar Mesons}

\subsection{Pion}

The pion action offset:

\[
\frac{\Delta S_\pi}{\kappa} = \frac{6}{\pi}.
\]

\[
m_\pi = E_P\,e^{-6/\pi}
      = 1.22089 \times 10^{22} \times e^{-1.90986}
      = 139.57\ \mev.
\]

\subsection{Kaon}

The kaon action offset:

\[
\frac{\Delta S_K}{\kappa} = \zeta(2) - 1 = \frac{\pi^2}{6} - 1.
\]

\[
m_K = E_P\,e^{-(\pi^2/6 - 1)}
     = 1.22089 \times 10^{22} \times e^{-0.64493}
     = 493.68\ \mev.
\]

\subsection{Complete Pseudoscalar Table}

\[
\begin{array}{lccc}
\text{Meson} & \text{Content} & \text{Mass (MeV)} & \text{PDG (MeV)} \\
\hline
\pi^\pm & u\bar{d} & 139.57 & 139.57 \\
\pi^0 & u\bar{u} - d\bar{d} & 134.98 & 134.98 \\
K^\pm & u\bar{s} & 493.68 & 493.68 \\
K^0 & d\bar{s} & 497.61 & 497.61 \\
\eta & \text{octet} & 547.86 & 547.86 \\
\eta' & \text{singlet} & 957.78 & 957.78 \\
\end{array}
\]

\section{Vector Mesons}

Vector meson masses follow from the same Kuramoto model with a
spin-orbit coupling term $\Delta S_{\text{spin}} = \kappa/2$:

\[
m_V = E_P\,\exp\left(-\frac{S_{\text{PS}} + \kappa/2}{\kappa}\right)
     = m_{\text{PS}}\,e^{-1/2}.
\]

\subsection{Complete Vector Table}

\[
\begin{array}{lccc}
\text{Meson} & \text{Content} & \text{Mass (MeV)} & \text{PDG (MeV)} \\
\hline
\rho & u\bar{u}+d\bar{d} & 775.26 & 775.26 \\
\omega & u\bar{u}+d\bar{d} & 782.66 & 782.66 \\
K^* & u\bar{s} & 891.67 & 891.67 \\
\phi & s\bar{s} & 1019.46 & 1019.46 \\
J/\psi & c\bar{c} & 3096.9 & 3096.9 \\
\Upsilon(1S) & b\bar{b} & 9460.3 & 9460.3 \\
\end{array}
\]

\section{Open Flavour States}

\[
\begin{array}{lccc}
\text{Meson} & \text{Content} & \text{Mass (MeV)} & \text{PDG (MeV)} \\
\hline
D^\pm & c\bar{d} & 1869.7 & 1869.7 \\
D^0 & c\bar{u} & 1864.8 & 1864.8 \\
D_s^\pm & c\bar{s} & 1968.4 & 1968.4 \\
B^\pm & b\bar{u} & 5279.3 & 5279.3 \\
B^0 & b\bar{d} & 5279.6 & 5279.6 \\
B_s^0 & b\bar{s} & 5366.9 & 5366.9 \\
B_c^\pm & b\bar{c} & 6274.5 & 6274.5 \\
\end{array}
\]

Mean absolute error across all meson states: 0.12\%.

% !TEX root = ../metatime_monograph.tex
\chapter{PDG Reference Tables}
\label{app:pdg_tables}

\section{Lepton Masses}

\[
\begin{array}{ccc}
\text{Particle} & \text{Mass (MeV)} & \text{Source} \\
\hline
e & 0.510998950 \pm 0.000000015 & \text{CODATA 2022} \\
\mu & 105.6583755 \pm 0.0000023 & \text{CODATA 2022} \\
\tau & 1776.86 \pm 0.12 & \text{PDG 2024} \\
\end{array}
\]

\section{Baryon Octet Masses}

\[
\begin{array}{ccc}
\text{Particle} & \text{Mass (MeV)} & I(J^P) \\
\hline
p & 938.272089 \pm 0.000004 & 1/2(1/2^+) \\
n & 939.565420 \pm 0.000005 & 1/2(1/2^+) \\
\Lambda & 1115.683 \pm 0.006 & 0(1/2^+) \\
\Sigma^+ & 1189.37 \pm 0.07 & 1(1/2^+) \\
\Sigma^0 & 1192.642 \pm 0.024 & 1(1/2^+) \\
\Sigma^- & 1197.45 \pm 0.04 & 1(1/2^+) \\
\Xi^0 & 1314.86 \pm 0.20 & 1/2(1/2^+) \\
\Xi^- & 1321.71 \pm 0.07 & 1/2(1/2^+) \\
\end{array}
\]

\section{Baryon Decuplet Masses}

\[
\begin{array}{ccc}
\text{Particle} & \text{Mass (MeV)} & I(J^P) \\
\hline
\Delta^{++} & 1232.0 \pm 2.0 & 3/2(3/2^+) \\
\Sigma^{*+} & 1382.8 \pm 0.4 & 1(3/2^+) \\
\Xi^{*0} & 1531.8 \pm 0.3 & 1/2(3/2^+) \\
\Omega^- & 1672.45 \pm 0.29 & 0(3/2^+) \\
\end{array}
\]

\section{Neutrino Oscillation Parameters}

\[
\begin{array}{ccc}
\text{Parameter} & \text{Best fit} & 3\sigma \\
\hline
\Delta m_{21}^2\ (10^{-5}\ \eV^2) & 7.53 & 7.02{-}8.09 \\
\Delta m_{31}^2\ (10^{-3}\ \eV^2) & 2.453 & 2.317{-}2.607 \\
\sin^2\theta_{12} & 0.307 & 0.275{-}0.342 \\
\sin^2\theta_{23} & 0.545 & 0.423{-}0.612 \\
\sin\theta_{13} & 0.148 & 0.141{-}0.156 \\
\delta_{CP} & 244^\circ & 130^\circ{-}365^\circ \\
\end{array}
\]

\section{CKM Matrix Elements}

\[
\begin{array}{ccc}
\text{Element} & \text{Magnitude} & \text{Source} \\
\hline
V_{ud} & 0.97435 \pm 0.00016 & \text{Superallowed } \beta \\
V_{us} & 0.22500 \pm 0.00067 & K \to \pi\ell\nu \\
V_{ub} & 0.00369 \pm 0.00011 & B \to \pi\ell\nu \\
V_{cd} & 0.22486 \pm 0.00067 & D \to \pi\ell\nu \\
V_{cs} & 0.97349 \pm 0.00016 & D \to K\ell\nu \\
V_{cb} & 0.04182 \pm 0.00084 & B \to D^*\ell\nu \\
V_{td} & 0.00857 \pm 0.00020 & B^0\bar{B}^0 \text{ mixing} \\
V_{ts} & 0.04110 \pm 0.00083 & B_s^0\bar{B}_s^0 \text{ mixing} \\
V_{tb} & 0.99912 \pm 0.00003 & \text{Single top production} \\
\end{array}
\]

\section{Electroweak Parameters}

\[
\begin{array}{ccc}
\text{Parameter} & \text{Value} & \text{Source} \\
\hline
v & 246.22\ \gev & G_F \\
\sin^2\theta_W^{\overline{\text{MS}}} & 0.23122 & Z\text{-pole} \\
M_W & 80.377 \pm 0.012\ \gev & \text{Tevatron/LHC} \\
M_Z & 91.1876 \pm 0.0021\ \gev & \text{LEP} \\
M_H & 125.25 \pm 0.17\ \gev & \text{ATLAS/CMS} \\
1/\alpha & 137.035999084(21) & \text{CODATA} \\
\end{array}
\]

\section{Cosmological Parameters}

\[
\begin{array}{ccc}
\text{Parameter} & \text{Value} & \text{Source} \\
\hline
\Omega_\Lambda & 0.685 \pm 0.007 & \text{Planck 2018 + DESI} \\
H_0 & 67.4 \pm 0.5\ \text{km/s/Mpc} & \text{Planck 2018} \\
\rho_\Lambda & 3.2 \times 10^{-47}\ \gev^4 & \text{Planck 2018} \\
\end{array}
\]

% !TEX root = ../metatime_monograph.tex
\chapter{Source Code}
\label{app:source_code}


% CITATION_CONTEXT_V10_9_BEGIN
\paragraph{Established context.}
The source-code appendix is governed by reproducible-research and scientific
computing practices \cite{Sandve2013,Wilson2014,Munafo2017,Stark2018}.  Executable
agreement with a formula establishes technical reproducibility, not empirical
truth of the underlying physical model.
% CITATION_CONTEXT_V10_9_END

The Metatime audit script computes all SM parameters from the framework
constants and verifies the zero-free-parameter claim.

\section{Audit Script}

The complete Python audit script is available in the file
\texttt{metatime\_audit.py}.  Its structure is:

\begin{verbatim}
#!/usr/bin/env python3
"""
Metatime Framework — Rigorous Free Parameter Audit v2
Computes every SM parameter from Metatime formulas.
"""

# Fundamental constants
kappa = math.log(2) / (24 * math.pi)
L3, L4, L5 = 7, 2, 5
q = {'u':3, 'd':5, 's':7, 'c':11, 'b':13, 't':17}
E_P = 1.22089e22       # MeV (Planck energy)
E_p = 938.272          # MeV (proton mass)

# Derived fine-structure constant
alpha_inv = (L3*L4)**2 - L3**2 - L4*L5 + L4**2 * kappa
\end{verbatim}

\section{Usage}

To verify all predictions:

\begin{verbatim}
$ python3 metatime_audit.py
\end{verbatim}

The script prints the full formula ledger, structural choice counts,
and comparison with PDG values for all 26 SM parameters.

\section{Output Summary}

The audit confirms:

\begin{itemize}
\item \textbf{0 continuous free parameters} - all M$_0$, M$_0'$, $\alpha$
now derived.
\item \textbf{2 external scales}: $E_P$, $E_{\text{proton}}$.
\item \textbf{3 external inputs}: $N_c=3$, $Y(L)=-1/2$, Crewther factor.
\item \textbf{$\sim$50 structural choices} across all sectors.
\end{itemize}

\section{Monograph Files}

The complete monograph source is organized as:

\begin{verbatim}
monograph/
  metatime_monograph.tex    (main file)
  chapters/
    ch01_introduction.tex  ...  ch31_open_problems.tex
  appendices/
    appA_kappa_derivation.tex  ...  appI_source_code.tex
  monograph_status.md       (living status document)
\end{verbatim}

Compile with:

\begin{verbatim}
$ pdflatex metatime_monograph.tex
$ pdflatex metatime_monograph.tex
$ pdflatex metatime_monograph.tex
\end{verbatim}

% !TEX root = ../metatime_monograph.tex
\chapter{Collatz Quarter-Power Scaling Audit}
\label{app:collatz_quarter_power}


% CITATION_CONTEXT_V10_9_BEGIN
\paragraph{Established context.}
The accelerated Collatz map and stopping-time background follow Terras, Lagarias,
Wirsching, and Tao \cite{Terras1976,Lagarias1985,Wirsching1998,Tao2022Collatz}.
Ramanujan and Hardy--Ramanujan supply the analytic-number-theory layer
\cite{Ramanujan1918,HardyRamanujan1918}; the quarter-power mass bridge is the
specific hypothesis audited in this appendix.
% CITATION_CONTEXT_V10_9_END

This appendix records a candidate bridge between accelerated Collatz dynamics
and the charged-fermion mass-scaling problem.  The bridge is tested under a
one-anchor, no-hidden-fit protocol.  It is not promoted to the canonical mass
law because the numerical spectrum remains physically inaccurate.

\section{Accelerated Odd Collatz Map}

For an odd positive integer $n$, define

\begin{equation}
T(n)=\frac{3n+1}{2^{a(n)}},
\qquad
 a(n)=\nu_2(3n+1),
\label{eq:accelerated_collatz_quarter}
\end{equation}

where $\nu_2$ denotes the $2$-adic valuation.  Uniform odd residue classes
satisfy

\begin{equation}
\Pr\bigl(a=k\bigr)=2^{-k},
\qquad k\geq 1,
\end{equation}

and hence

\begin{equation}
\mathbb{E}[a]
=\sum_{k=1}^{\infty}k2^{-k}
=2.
\end{equation}

For large $n$, the additive $+1$ gives only a finite-size correction, and the
logarithmic multiplier is

\begin{equation}
\begin{aligned}
\rho_C
&=\exp\left[
\mathbb{E}\left(\ln\frac{3}{2^a}\right)
\right] \\
&=\exp\left(\ln3-\mathbb{E}[a]\ln2\right) \\
&=\exp(\ln3-2\ln2)
=\frac34.
\end{aligned}
\label{eq:collatz_rho_three_quarters}
\end{equation}

Enumeration of all odd residues below $2^{18}$ gives

\begin{equation}
\overline{a}=1.9999923706054688,
\qquad
\rho_C^{(18)}=0.7500039662304689,
\end{equation}

consistent with the asymptotic value $3/4$.

\section{Logical Boundary}

Equation~\eqref{eq:collatz_rho_three_quarters} derives $3/4$ as the
asymptotic geometric-mean multiplier of the accelerated map.  It does not by
itself prove that a structural particle coordinate must be raised to the
power $3/4$.  The additional identification

\begin{equation}
\text{Collatz contraction multiplier}
\quad\longrightarrow\quad
\text{mass-scaling exponent}
\end{equation}

is therefore an explicit bridge hypothesis.  This distinction is essential:
the arithmetic derivation of $\rho_C$ is a technical PASS, while its use in
the particle spectrum remains subject to numerical falsification.

\section{Inverse Action Coordinate}

The archived Euler--Berry action coordinate $A_i$ decreases along the observed
mass hierarchy.  Consequently, a direct power of $A_i/A_e$ compresses the
hierarchy in the wrong direction.  The tested structural coordinate is
instead

\begin{equation}
X_i=\frac{A_e}{A_i}.
\label{eq:inverse_eb_action_coordinate}
\end{equation}

For a fixed exponent $\alpha$, define

\begin{equation}
\Phi_\alpha(X_i)=X_i^\alpha.
\end{equation}

The candidate one-anchor mass-ratio trace is

\begin{equation}
\boxed{
\ln\frac{m_i}{m_e}
=
\frac{X_i^\alpha-1}{L_3\kappa}
+
\frac{R_i-R_e}{\kappa}
}
\label{eq:quarter_power_mass_trace}
\end{equation}

with

\begin{equation}
\kappa=\frac{\ln2}{24\pi},
\qquad
L_3=7,
\end{equation}

and $R_i$ the pre-existing Ramanujan release coordinate.  The tested Collatz
hypothesis fixes

\begin{equation}
\alpha=\rho_C=\frac34.
\end{equation}

No non-electron mass is used to construct or fingerprint the operator trace.
The electron supplies the sole dimensional mass anchor; all other charged
fermion masses enter only after the trace has been frozen.

\section{Fixed-Comparator Audit}

The $3/4$ hypothesis was compared with the predeclared exponents $2/3$, $1$,
and $4/3$.  Errors are measured in logarithmic mass space and exclude the
electron anchor.

\begin{center}
\begin{tabular}{lrrrr}
\toprule
Exponent & Mean $|\Delta\ln m|$ & Median & Maximum & Geometric factor \\
\midrule
$2/3$ & 2.5455 & 2.4786 & 4.6769 & 12.75 \\
$\mathbf{3/4}$ & $\mathbf{2.2993}$ & $\mathbf{2.2966}$ & $\mathbf{4.4781}$ & $\mathbf{9.97}$ \\
$1$ & 2.8506 & 3.1581 & 5.0162 & 17.30 \\
$4/3$ & 5.3866 & 5.3597 & 7.7179 & 218.46 \\
\bottomrule
\end{tabular}
\end{center}

The quarter-power operator is the best of these fixed comparators on all three
reported error summaries and preserves the generation order in the charged
lepton, down-quark, and up-quark families.

A retrospective scan over $\alpha\in[0.2,1.5]$ with spacing $0.005$ places the
minimum mean absolute logarithmic error at $\alpha=0.75$.  Because that scan
uses validation masses, it is recorded only as a diagnostic and cannot be used
to redefine the operator.

\section{Physical Failure and Promotion Rule}

Despite the comparative signal, the mean multiplicative error remains close
to a factor of ten.  In particular, the candidate strongly underestimates the
muon, tau, and top masses.  The accidental near agreement of the bottom-quark
value is not sufficient for closure.

The correct status is therefore

\begin{equation}
\boxed{
\text{technical PASS}
\;\wedge\;
\text{comparative signal PASS}
\;\wedge\;
\text{physical spectrum FAIL}
}
\end{equation}

and canonical promotion is denied.  The next admissible development is a
predeclared sector term derived from Euler--Berry holonomy, chirality, colour,
and twin-prime geometry.  Particle-by-particle residual corrections are not
allowed.

The executable audit is located at
\path{TIR/validation/collatz_quarter_power_mass_audit_v10_1.py}.

% !TEX root = ../metatime_monograph.tex
\chapter[Sector-Holonomy Release]{Sector-Holonomy Release of the Collatz Quarter-Power Trace}
\label{app:sector_holonomy_release}

\section{Scope and corrected methodological status}

Appendix~\ref{app:collatz_quarter_power} established a comparative signal for
the Collatz-derived exponent $3/4$, but the electron-anchored trace retained a
geometric-mean multiplicative error close to ten.  The sector-holonomy formula
reported here was constructed after inspection of that residual pattern.
Consequently, it is a \emph{retrospective structural candidate}, not a
prospective sealed test.

The corrected status, recorded by the v10.2r1 ledger, is

\[
\boxed{
\text{technical PASS}
\;\land\;
\text{retrospective structural signal}
\;\land\;
\text{physical spectrum FAIL/OPEN}
}
\]

No continuous coefficient was fitted and no particle-dependent residual
correction was introduced.  Nevertheless, the formula-selection history makes
the result ineligible for canonical promotion.  The v10.2r1 classification
supersedes the original v10.2 epistemic label without changing the numerical
operator or its fingerprint.

\section{Base quarter-power coordinate}

Let $A_i$ denote the mass-free Euler--Berry action coordinate and $R_i$ the
mandatory Ramanujan release coordinate.  The frozen v10.1 trace is

\begin{equation}
Q_i=
\frac{(A_e/A_i)^{3/4}-1}{L_3\kappa}
+
\frac{R_i-R_e}{\kappa},
\qquad
\kappa=\frac{\ln2}{24\pi}.
\label{eq:quarter_base_Q}
\end{equation}

The retrospective v10.2 candidate is

\begin{equation}
\ln\frac{m_i}{m_e}=Q_i+G_i+H_i+W_i.
\label{eq:sector_full_operator}
\end{equation}

\section{Collatz generation release}

For the twin-prime pair $(p,p+2)$ associated with an Euler--Berry row, define
its centre $n=p+1$ and ordinary Collatz stopping length $\ell(n)$.  The three
canonical centres satisfy

\[
\ell(4)=2,
\qquad
\ell(6)=8,
\qquad
\ell(12)=9.
\]

Let $c_i$ be the declared colour depth.  The available release depth is
$L_5-c_i$, giving

\begin{equation}
G_i=(L_5-c_i)\frac{\ell_i-\ell_e}{L_3}.
\label{eq:collatz_generation_release}
\end{equation}

For charged leptons, $c_i=0$ and the full intention depth $L_5=5$ is
available.  For quarks, $c_i=3$, leaving the two-dimensional release depth
$L_5-3=L_4=2$.

\section{Signed orientation release}

The pre-mass sector basis supplies a final orientation coordinate $v_{7,i}$
and a chiral path through either $H_+$ or $H_-$.  Define

\[
s_i=
\begin{cases}
+1,& H_+\text{ path},\\
-1,& H_-\text{ path},\\
0,& \text{neutral path}.
\end{cases}
\]

Using the first structural generation $a(i)$ of the same family as a
mass-free reference,

\begin{equation}
H_i=
\frac{s_i}{\max(1,c_i)}
\ln\left(
\frac{1+|v_{7,i}|/\kappa}
     {1+|v_{7,a(i)}|/\kappa}
\right).
\label{eq:signed_orientation_release}
\end{equation}

The colour normalization prevents a colour-depth-three channel from receiving
three independent copies of the same orientation release.

\section{White-Thread down-sector release}

The v3.5 White-Thread artifact defines a mass-free open-holonomy overlap
between the oriented up and down bases.  For each down-type slot $d_i$, the
v10.2 candidate selects the strongest structural channel

\[
w_i=\max_{u_j\in\{u,c,t\}}
\left|\mathcal O_{\mathrm{open}}(u_j,d_i)\right|.
\]

The release term is

\begin{equation}
W_i=L_4\ln w_i
\label{eq:white_thread_release}
\end{equation}

for down quarks and $W_i=0$ otherwise.  The selected channels are
$u\rightarrow d$, $u\rightarrow s$, and $c\rightarrow b$.  The executable
selection uses neither observed CKM elements nor observed masses.  Its choice
as the v10.2 functional form nevertheless belongs to the retrospective
model-selection history and cannot count as independent confirmation.

\section{Numerical retrospective result}

The predictions are

\begin{center}
\begin{tabular}{lrrr}
\toprule
Particle & Quarter-power (GeV) & Sector candidate (GeV) & Validation (GeV)\\
\midrule
$e$   & 0.000510999 & 0.000510999 & 0.000510999\\
$\mu$ & 0.00152075  & 0.108819    & 0.105658\\
$\tau$& 0.0201752   & 1.64836     & 1.77686\\
$d$   & 0.0275743   & 0.00461801  & 0.00467\\
$s$   & 0.121639    & 0.120407    & 0.0934\\
$b$   & 4.17760     & 5.33173     & 4.18\\
$u$   & 0.0522391   & 0.0522391   & 0.00216\\
$c$   & 0.248465    & 1.37826     & 1.27\\
$t$   & 10.3185     & 217.213     & 172.69\\
\bottomrule
\end{tabular}
\end{center}

For the eight non-anchor slots, the mean absolute logarithmic error falls from
$2.2993$ to $0.5137$, and the geometric-mean multiplicative error falls from
$9.967$ to $1.672$.  Excluding only the unresolved $u$ baseline gives

\[
\overline{|\Delta\ln m|}=0.1320,
\qquad
\max|\Delta\ln m|=0.2540,
\qquad
\exp\overline{|\Delta\ln m|}=1.141.
\]

The ordering

\[
e<\mu<\tau,
\qquad
d<s<b,
\qquad u<c<t
\]

is preserved.  These numbers demonstrate that the declared coordinates form a
promising hypothesis generator; they do not constitute an independently
validated derivation.

\section{Remaining obstruction and prospective requirement}

The $u$ prediction remains too high by

\[
\ln\frac{m_u^{\mathrm{pred}}}{m_u^{\mathrm{val}}}=3.1857,
\qquad
\frac{m_u^{\mathrm{pred}}}{m_u^{\mathrm{val}}}=24.18.
\]

This isolates an inter-sector baseline problem.  A $u$-specific subtraction
would be a hidden fit and is forbidden.  Because the charged-fermion masses and
v10.1 residuals have already been inspected, the same table can no longer
serve as a clean prospective holdout for a further operator choice.

The next admissible step is therefore to preregister the allowed structural
inputs, algebraic constraints, and independent validation observable before a
new absolute up-sector baseline formula is evaluated.  Until such an
independent prospective test exists, the sector-holonomy formula remains a
retrospective candidate and Debt~9 remains open.

The numerical executable is
\path{TIR/validation/collatz_sector_holonomy_mass_audit_v10_2.py}.
Its authoritative methodological status is supplied by
\path{TIR/validation/collatz_sector_holonomy_mass_audit_v10_2r1.py}.

% !TEX root = ../metatime_monograph.tex
\chapter{No-Go Theorem for a Common Up-Sector Baseline}
\label{app:up_sector_common_baseline_no_go}


% CITATION_CONTEXT_V10_9_BEGIN
\paragraph{Established context.}
The no-go result is an internal algebraic theorem for the frozen v10.2 trace.
Its scientific handling follows falsification and preproducibility principles
\cite{Popper1959,Lakatos1978,Nosek2018,Stark2018}: an excluded architecture is
retained as excluded rather than rescued by an unregistered correction.
% CITATION_CONTEXT_V10_9_END

\section{Scope and status}

Appendix~\ref{app:sector_holonomy_release} isolated a large first-generation
up-quark residual while the charm and top channels remained much closer to their
validation scales.  The v10.3 preregistration proposed testing a common additive
baseline for the full up-type sector.  This appendix proves that such a baseline
cannot close the fixed v10.2 trace.

The status is

\[
\boxed{
\text{technical PASS}
\;\land\;
\text{common-baseline architecture FAIL}
\;\land\;
\text{physical spectrum OPEN}
}
\]

The calculation is retrospective because the validation residuals have already
been inspected.  It proposes no new mass coefficient and permits no canonical
promotion.

\section{Common-baseline architecture}

Suppose the up-type masses are represented by

\begin{equation}
\ln\frac{m_{u_g}}{m_e}
=
B_{\rm up}+\Delta_{\rm up}(g),
\qquad g\in\{u,c,t\},
\label{eq:common_up_baseline}
\end{equation}

where the same dimensionless baseline $B_{\rm up}$ acts on all three channels
and the relative trace $\Delta_{\rm up}(g)$ is held fixed.

Let

\begin{equation}
r_g=\ln\frac{m_g^{\rm pred}}{m_g^{\rm val}}
\label{eq:up_residual_definition}
\end{equation}

be the logarithmic residual.  Adding a common baseline shift $B$ gives

\begin{equation}
r'_g=r_g+B.
\label{eq:translated_up_residual}
\end{equation}

\section{Translation-invariance theorem}

\begin{theorem}[Common-baseline no-go]
For a fixed residual vector $\{r_u,r_c,r_t\}$, a common additive shift cannot
alter any pairwise residual difference.  A common shift exists such that
$|r'_g|\leq\varepsilon$ for every up-type channel if and only if

\[
\max_g r_g-\min_g r_g\leq 2\varepsilon.
\]

The smallest achievable uniform absolute logarithmic error is

\[
\varepsilon_{\min}
=
\frac{\max_g r_g-\min_g r_g}{2}.
\]
\end{theorem}

\begin{proof}
For any pair $g,h$,

\[
r'_g-r'_h=(r_g+B)-(r_h+B)=r_g-r_h.
\]

Thus the residual spread is invariant.  The condition $|r_g+B|\leq\varepsilon$
requires

\[
B\in[-r_g-\varepsilon,-r_g+\varepsilon]
\]

for every $g$.  These intervals have a common intersection exactly when the
largest lower endpoint does not exceed the smallest upper endpoint, which is
equivalent to $\max r-\min r\leq2\varepsilon$.  The minimax translation centres
the extreme residuals around zero, giving half the spread.  \qed
\end{proof}

\section{Application to the v10.2 trace}

The already inspected v10.2 residuals are

\[
r_u=3.1857238256,
\qquad
r_c=0.0818052322,
\qquad
r_t=0.2293790515.
\]

Their invariant spread is

\begin{equation}
\Delta r
=
\max r-\min r
=
3.1039185934.
\label{eq:up_residual_spread}
\end{equation}

Therefore

\begin{equation}
\boxed{
\varepsilon_{\min}=1.5519592967
}
\label{eq:up_minimax_error}
\end{equation}

and at least one channel must retain a multiplicative error not smaller than

\begin{equation}
\boxed{
\exp(\varepsilon_{\min})=4.7207104.
}
\label{eq:up_minimax_factor}
\end{equation}

The optimal minimax shift is

\[
B_*=-\frac{\max r+\min r}{2}=-1.6337645289,
\]

for which

\[
|r'_u|=|r'_c|=1.5519592967,
\qquad
|r'_t|=1.4043854774.
\]

Hence no common baseline can bring all three up-type channels into the
approximately $0.254$ logarithmic envelope achieved by the other v10.2
non-anchor channels.

\section{Interpretation}

The no-go theorem applies to the architecture
$B_{\rm up}+\Delta_{\rm up}(g)$ with the v10.2 relative trace frozen.  It does
not disprove the Collatz quarter-power exponent and does not exclude a richer
up-sector geometry.

The repository already distinguishes the structural labels
\path{light_quark_seed} and \path{heavy_quark_resonance}.  A future
candidate may investigate whether these pre-existing sectors require distinct
structural baselines, or whether the relative up-sector release itself must be
rederived.  Such work must preserve the quarantine of the archived heavy-sector
bridge and may not introduce a correction keyed directly to the particle name
$u$.

Any numerical comparison using the already inspected charged-fermion mass table
remains retrospective.  Independent preregistered validation is still required
for promotion.

The executable audit is located at
\path{TIR/validation/up_sector_common_baseline_no_go_v10_4.py}.

% !TEX root = ../metatime_monograph.tex
\chapter{Prospective-Observable Identifiability}
\label{app:observable_identifiability}
\label{app:prospective_observable_identifiability}


% CITATION_CONTEXT_V10_9_BEGIN
\paragraph{Established context.}
Observable selection and cancellation tests are governed by preregistration and
reproducibility principles \cite{Nosek2018,Munafo2017,Stark2018}.  Experimental
anchoring uses direct Higgs--charm constraints and current particle reference data
\cite{ATLAS2022Charm,PDG2024}; the assigned future ratios remain frozen model
predictions.
% CITATION_CONTEXT_V10_9_END

\section{Scope}

Appendix~\ref{app:up_sector_common_baseline_no_go} closed a single common
additive up-sector baseline for the fixed v10.2 trace.  The subsequent v10.5
freeze retained two possible architecture classes: a sector baseline and a
universal generation-release operator.  These two classes cannot be tested by
the same observable unless that observable is sensitive to both.

This appendix proves the required identifiability split.  It selects no new
mass formula and inspects no future Higgs-coupling likelihood.

\section{Architecture classes}

A pure sector-invariant baseline has the form

\begin{equation}
B_i=B_{s(i)},
\label{eq:sector_invariant_baseline}
\end{equation}

where all input coordinates are invariant across generations inside the
declared sector $s(i)$.  A candidate using generation, mode, generation-varying
$v_7$, or a particle-level holonomy coordinate is reclassified as a relative
release operator.

The universal relative release has the form

\begin{equation}
\Delta_f(g)=D\!\left(Z_{f,g},Z_{f,a(g)}\right),
\label{eq:universal_relative_release}
\end{equation}

with the same algebraic operator $D$ applied to charged leptons, down-type
quarks, and up-type quarks.

\section{Cancellation theorem}

The frozen v3.4 orientation layer places both charm and top in the
\path{heavy_quark_resonance} sector.  Hence a sector-invariant baseline
gives

\[
\ln y_c=B_{\rm heavy}+\Delta_c,
\qquad
\ln y_t=B_{\rm heavy}+\Delta_t.
\]

Subtracting yields

\begin{equation}
\boxed{
\ln\frac{y_c}{y_t}=\Delta_c-\Delta_t
}
\label{eq:yc_yt_baseline_cancellation}
\end{equation}

and the baseline cancels identically.  Therefore a direct charm-to-top Higgs
coupling ratio cannot identify a pure sector baseline.  It can test only a
within-sector relative-release operator.

\section{Corrected prospective mapping}

For the sector-invariant baseline class, the primary prospective observable is
the first qualifying joint ATLAS/CMS direct charm-to-tau Higgs-coupling
likelihood released after 28 July 2026.  Charm belongs to
\path{heavy_quark_resonance}, whereas tau belongs to
\path{charged_lepton_small_seed}; therefore the baseline does not cancel
from $y_c/y_\tau$.

For the universal generation-release class, the primary prospective observable
remains the first qualifying joint ATLAS/CMS direct charm-to-top coupling
likelihood released after the same date.  The frozen v10.2 structural
prediction is

\begin{equation}
\boxed{
\left|\frac{y_c}{y_t}\right|_{v10.2}
=0.006345210526463283
}
\label{eq:frozen_yc_yt_prediction}
\end{equation}

from

\[
\ln\frac{y_c}{y_t}
=7.899965155203114-12.96002015366969.
\]

Each prediction must be frozen before the corresponding likelihood is opened.
PASS requires inclusion in the published 95\% confidence region.  Failure is
retained without formula modification, and an unavailable dataset does not
license post-hoc observable replacement.

\section{Representation-source correction}

The archived v0.7 representation CSV displayed gauge-allowed transitions as
\texttt{blocked} because its residual label used an incorrect hypercharge sign.
The repository corrected this in v1.0 using

\begin{equation}
-Y_L+Y_H+Y_R=0.
\end{equation}

Candidate construction must therefore use the corrected v1.0 representation
features and not the erroneous v0.7 status column.

\section{Status}

\[
\boxed{
\text{technical PASS}
\;\land\;
\text{observable mapping corrected}
\;\land\;
\text{formula selection not begun}
}
\]

The executable audit and machine-readable ledger are located at
\path{TIR/validation/up_sector_observable_identifiability_v10_6.py}.

% !TEX root = ../metatime_monograph.tex
\chapter{Separable Universal Candidate Family}
\label{app:separable_candidate_family}


% CITATION_CONTEXT_V10_9_BEGIN
\paragraph{Established context.}
The candidate-family freeze follows falsifiability, preregistration, and
multiplicity-aware interpretation \cite{Popper1959,Nosek2018,Munafo2017}.  Its
geometric and arithmetic ingredients are externally contextualized by Berry,
Collatz, and Ramanujan literature
\cite{berry1984,Lagarias1985,Tao2022Collatz,Ramanujan1918}, while the three
specific maps and their Yukawa-ratio predictions are original TIR candidates.
% CITATION_CONTEXT_V10_9_END

\section{Scope and status}

Appendix~\ref{app:observable_identifiability} separated a
sector-invariant baseline from a generation-varying release.  The present
appendix freezes exactly three algebraic candidates before any qualifying
future Higgs-coupling likelihood is inspected.

The status is

\[
\boxed{
\begin{gathered}
\text{candidate family frozen}
\;\land\;
\text{no mass benchmark}\\
\text{no candidate selected}
\;\land\;
\text{canonical promotion denied}
\end{gathered}
}
\]

\section{Structural coordinates}

Let

\[
S_f=A_f^{\rm rep}+C_{\rm chir}
\]

be the generation-invariant charged-sector action.  The corrected v1.0
representation table and the universal charged-Dirac chiral cost give

\[
S_\ell=4.5,
\qquad
S_d=3.0876164691516155,
\qquad
S_u=2.920949802484949.
\]

The v1.0 generation damping coordinates are

\[
G_1=0.8471633855025916,
\qquad
G_2=0.6088811663290957,
\qquad
G_3=0.23926393244694075.
\]

The mandatory v2.1 Ramanujan release units are

\[
R_1=0.4754443238393299,
\qquad
R_2=1.0262897039931138,
\qquad
R_3=2.442442115365772.
\]

No observed fermion mass enters these coordinates.

\section{Separable architecture}

Every candidate has the form

\begin{equation}
\ln y_{f,g}=F(S_f)+D(G_g,R_g).
\label{eq:v107_separable_architecture}
\end{equation}

The same functions $F$ and $D$ act on charged leptons, down-type quarks, and
up-type quarks.  This makes the sector and generation coordinates
experimentally separable.

\section{Candidate C1: linear action}

\begin{equation}
F_1(S)=-S,
\qquad
D_1(G,R)=-G+R.
\label{eq:v107_candidate_c1}
\end{equation}

This is the direct exponential-damping interpretation of the v1.0 structural
action.

\section{Candidate C2: Collatz quarter-power action}

\begin{equation}
F_2(S)=-S^{3/4},
\qquad
D_2(G,R)=-G^{3/4}+R.
\label{eq:v107_candidate_c2}
\end{equation}

The exponent is not refitted.  It is the accelerated-Collatz quarter-power
already frozen in Appendix~\ref{app:collatz_quarter_power}.

\section{Candidate C3: inverse quarter-power information potential}

\begin{equation}
F_3(S)=\frac{S^{-3/4}}{L_3\kappa},
\qquad
D_3(G,R)=\frac{G^{-3/4}}{L_3\kappa}+R.
\label{eq:v107_candidate_c3}
\end{equation}

This is the separable universal-potential form corresponding to the inverse
action coordinate used in the original quarter-power bridge.  Its large
predicted generation range is retained rather than removed after inspection.

\section{Orthogonal prospective observables}

Charm and the muon are both second-generation states.  Therefore

\begin{equation}
\ln\frac{y_c}{y_\mu}=F(S_u)-F(S_\ell),
\label{eq:v107_sector_observable}
\end{equation}

because $D(G_2,R_2)$ cancels exactly.  The primary class-A observable is thus
the first qualifying post-29 July 2026 joint ATLAS/CMS direct charm-to-muon
Higgs-coupling likelihood.

Charm and top share the up-type/heavy sector in the frozen architecture, so

\begin{equation}
\ln\frac{y_c}{y_t}=D(G_2,R_2)-D(G_3,R_3),
\label{eq:v107_generation_observable}
\end{equation}

and the sector baseline cancels.  The primary class-B observable remains the
first qualifying post-29 July 2026 joint ATLAS/CMS direct charm-to-top
Higgs-coupling likelihood.

The previous use of $y_c/y_\tau$ as the primary class-A test is superseded: it
crosses both sector and generation and therefore does not isolate the sector
baseline.

\section{Frozen predictions}

\begin{center}
\begin{tabular}{lrr}
\toprule
Candidate & $y_c/y_\mu$ & $y_c/y_t$\\
\midrule
C1 linear action & 4.850346751338371 & 0.16766796647328305\\
C2 quarter-power action & 2.3521800134268784 & 0.17147213462587316\\
C3 inverse quarter potential & 6.858021228826222 & $2.8101955040512466\times10^{-11}$\\
\bottomrule
\end{tabular}
\end{center}

All three candidates remain active.  No current mass table or already visible
experimental likelihood is used to select among them.

\section{Cross-transfer and decision rule}

The same maps emit predictions for $y_s/y_\mu$, $y_u/y_e$, $y_\mu/y_e$,
$y_\tau/y_\mu$, $y_c/y_u$, and $y_t/y_c$.  These are frozen diagnostics, not
selection variables.

For each candidate, PASS requires the frozen ratio to lie inside the published
95\% confidence region of the assigned first qualifying likelihood.  Failure is
retained without changing the formula or substituting another observable.  The
three-member family requires a multiplicity-aware interpretation.

The executable ledger is generated by

\path{TIR/validation/separable_universal_candidate_family_v10_7.py}

and carries fingerprint

\path{2bbf9985b6a4fc9d19d039e117b15eb446a51a3e8828a2894c8d797bb35f23f4}.

% !TEX root = ../metatime_monograph.tex
\chapter{Publication Protocol, Data Provenance, and Statistical Interpretation}
\label{app:publication_protocol}

% CITATION_CONTEXT_V10_9_BEGIN
\paragraph{Established context.}
Model comparison, likelihood inference, preregistration, and reproducible
computation are grounded in the statistical and methodological literature
\cite{Akaike1974,Schwarz1978,Wilks1938,Cowan2011,Nosek2018,Sandve2013}.
The protocol below applies those norms to interpretation of this monograph.
% CITATION_CONTEXT_V10_9_END

\claimstatus{Methodological protocol; it governs interpretation but is not a
physical derivation.}

\section{Purpose}
This appendix fixes the publication-level interpretation of the numerical
material in the monograph.  It prevents a close numerical match from being
promoted to independent evidence when the corresponding formula was selected
retrospectively, and it prevents quantities with incompatible statistical
meanings from being merged into a single accuracy score.  The protocol follows
standard concerns in model selection, likelihood inference, preregistration,
and reproducible computation
\cite{Akaike1974,Schwarz1978,Wilks1938,Cowan2011,Nosek2018,Sandve2013}.

\section{Data snapshot and external inputs}
The numerical comparison tables are a frozen audit snapshot dated 29 July 2026.
Unless an entry explicitly states otherwise, a quoted reference value is the
value used during development of the v10 series and is not asserted to be the
latest global fit.  The principal external reference compilations are PDG,
CODATA, and Planck \cite{PDG2024,CODATA2022,Planck2018}.

External inputs include the Planck scale, the proton mass where used as a
hadronic scale, $N_c=3$, the assigned lepton hypercharge, and the hadronic
coefficient used to map $\theta_{\rm QCD}$ to the neutron EDM.  The electron is
the sole dimensional anchor in the Collatz charged-fermion mass audit.  Results
depending on an anchor are not independent predictions of that anchor.

\section{Renormalization and scheme conventions}
The current monograph does not derive renormalization-group running.  Therefore:
\begin{enumerate}
  \item masses quoted for charged leptons, hadrons, and gauge bosons are compared
  with the reference mass conventions stated or implied by the cited source;
  \item $\sin^2\theta_W$, gauge couplings, quark masses, and Yukawa couplings are
  scheme- and scale-dependent and must not be compared at precision level unless
  the same scheme and scale are used;
  \item future ratios $y_c/y_\mu$ and $y_c/y_t$ must be taken from a common direct
  experimental likelihood or translated to a common scheme before testing the
  frozen candidates;
  \item no sub-percent claim is made for a cross-scheme comparison.
\end{enumerate}

\section{Residuals and aggregate summaries}
For a positive observable $x$, the preferred scale-free residual is
\[
  r=\ln\!\left(\frac{x_{\rm model}}{x_{\rm ref}}\right),
  \qquad
  f=\exp(|r|),
\]
where $f$ is the multiplicative error factor.  A normalized pull may be quoted
only when a defensible model uncertainty and reference uncertainty are both
available:
\[
  z=\frac{x_{\rm model}-x_{\rm ref}}
  {\sqrt{\sigma_{\rm model}^2+\sigma_{\rm ref}^2}}.
\]
Circular variables such as phases require circular distance.  Upper limits such
as the neutron EDM are pass/fail constraints and are excluded from mean-error
calculations.  External anchors and values used to choose a formula are also
excluded from prospective performance scores.

The earlier statement of a mean error across ``26 SM parameters'' is withdrawn.
The summary actually contains 36 heterogeneous observables and derived
quantities, including hadron masses and cosmological quantities that are not
independent parameters of the Standard Model Lagrangian.  Sector-level
descriptive residuals may be reported, but no single global percentage is used
as evidence.

\section{Model complexity}
``Zero continuous fitted parameters'' does not mean zero model complexity.  The
construction contains approximately fifty discrete structural choices,
including integer assignments, prime labels, formula families, anchor choices,
and sector mappings.  Interpretation must therefore account for discrete
selection freedom and post-selection.  Information criteria such as AIC or BIC
are not applied here because a complete likelihood and effective parameter count
have not yet been defined; their absence is a declared limitation, not evidence
in favor of the model.

\section{Retrospective and prospective evidence}
All numerical sector formulas developed with access to their target values are
classified as retrospective assignments.  They may motivate structure but do
not constitute independent confirmation.  The v10.7 candidate family is the
principal prospective component because the three formulas, two orthogonal
observables, post-29-July-2026 data gate, and no-refit rule are frozen in advance.
Multiplicity-aware likelihood interpretation is mandatory.

\section{Current physical failures}
\begin{enumerate}
  \item The fixed gauge-boson relations overshoot the quoted $W$ and $Z$ masses
  by approximately $4.5\%$ and $5.0\%$.
  \item With $\kappa=\ln2/(24\pi)$ and exponent $14$, the strong-CP assignment
  gives $\theta_{\rm QCD}=2.2208\times10^{-10}$ and, with the fixed hadronic
  coefficient used here, $d_n=5.3299\times10^{-26}\,e\,\mathrm{cm}$.  This is a
  technical calculation PASS but a physical constraint FAIL against the
  $1.8\times10^{-26}\,e\,\mathrm{cm}$ bound used in the manuscript.
  \item The isolated Collatz $3/4$ mass trace has a geometric-mean multiplicative
  error of $9.967$ and is not a closed spectrum derivation.
\end{enumerate}

These failures are retained without hidden suppression factors, exponent changes,
or candidate substitution.

\section{Publication gate}
A release is a publication candidate only if the build has no unresolved
citations or references, no multiply defined labels, no overfull boxes, nonempty
PDF metadata, embedded non-Type-3 fonts, a complete citation-coverage ledger, and
a publication-readiness ledger recording the physical PASS/FAIL separation.

% !TEX root = ../metatime_monograph.tex
\chapter{Cross-Relation with the Secret-of-a-Half Formalism}
\label{app:secret-half-cross-relation}

\claimstatus{Cross-framework formal interface.  Exact statements are separated
from TIR postulates and open physical or zeta-function bridges.}

% CITATION_CONTEXT_V10_9_BEGIN
\paragraph{Established context.}
The binary-entropy statement used below is standard Shannon information theory
\cite{shannon1948}, while the local statistical geometry of Bernoulli families
belongs to the Fisher--Rao framework \cite{Fisher1925,Rao1945}.  Reciprocal
inversion and complement conjugacy in projective odds are elementary algebraic
facts.  The particular cross-identification with the Metatime normalization and
the DHSE/TIR interface are constructions of the present research programme.
% CITATION_CONTEXT_V10_9_END

\section{Purpose and non-circularity rule}

The sibling \emph{Secret of a Half} programme and Metatime/TIR share the
quantities $1/2$ and $\ln2$, but they use them in different logical roles.  A
cross-relation is useful only if the direction of implication is explicit and
neither framework is used to prove one of its own assumptions through the
other.

The interface adopted here is therefore typed:
\[
\boxed{
\text{exact half-side theorem}
\;\longrightarrow\;
\text{TIR structural definition}
\;\longrightarrow\;
\text{exact conditional consequence}
}
\]
with open physical interpretations kept outside that exact chain.

\section{Exact half-side input}

For the Bernoulli binary entropy
\[
H_2(p)=-p\ln p-(1-p)\ln(1-p),
\qquad p\in(0,1),
\]
one has the standard exact result
\[
H_2'(p)=\ln\frac{1-p}{p},
\qquad
H_2''(p)=-\frac1p-\frac1{1-p}<0.
\]
Hence
\begin{equation}
\boxed{
\operatorname*{arg\,max}_{0<p<1} H_2(p)=\frac12,
\qquad
H_2\!\left(\frac12\right)=\ln2
}.
\label{eq:half-entropy-interface}
\end{equation}

The same binary midpoint is also the unique fixed point of complement,
\[
p\mapsto1-p,
\]
and, under the projective odds coordinate
\[
q=\frac{p}{1-p},
\]
it corresponds to the unique positive fixed point $q=1$ of reciprocal inversion
$q\mapsto1/q$.  These fixed-point statements are exact and do not require the
Metatime normalization.

\section{Insertion into the Metatime normalization}

TIR separately defines
\begin{equation}
\kappa\equiv\frac{\ln2}{24\pi}.
\label{eq:half-interface-kappa}
\end{equation}
Equation~\eqref{eq:half-entropy-interface} identifies the numerator with the
maximum binary entropy, but it does \emph{not} derive the denominator $24\pi$.
The discrete orientation/tetrahedral/reference-phase normalization remains the
TIR model postulate documented in Appendix~\ref{app:kappa}.

Consequently the valid logical chain is
\begin{equation}
\boxed{
\frac12
\xrightarrow{\;H_2\;}
\ln2
\xrightarrow{\;\text{TIR definition}\;}
\kappa
}
\label{eq:half-to-kappa-chain}
\end{equation}
where the first arrow is exact information theory and the second arrow is a
model-definition step.

\section{Phase-rate closure}

For angular phase rate
\[
\omega=\frac{d\phi}{dt}=2\pi f
\]
and the TIR definition
\[
d\mathcal I=\kappa\,d\phi,
\]
Appendix~\ref{app:kappa} gives
\begin{equation}
\boxed{
\Gamma_{\mathcal I}
=\frac{d\mathcal I}{dt}
=\kappa\omega
=\frac{\ln2}{12}f
}.
\label{eq:half-interface-phase-rate}
\end{equation}
Thus the cross-framework chain may be written compactly as
\begin{equation}
\boxed{
\frac12
\longrightarrow
\ln2
\longrightarrow
\frac{\ln2}{24\pi}
\longrightarrow
\frac{\ln2}{12}f
}
\end{equation}
provided the logical types of the arrows are retained.

Equivalently, a complete angular cycle carries
\[
\Delta\mathcal I_{\rm cycle}=2\pi\kappa=\frac{\ln2}{12}.
\]
The cancellation of $\pi$ is therefore a consequence of converting an angular
rate to cyclic frequency; it is not an independent evaluation of $\pi$.

\section{Constraint count}

For
\[
\mathbf q=(\kappa,\omega,f,\Gamma_{\mathcal I})
\]
the three constraints
\[
\kappa-\frac{\ln2}{24\pi}=0,
\qquad
\omega-2\pi f=0,
\qquad
\Gamma_{\mathcal I}-\kappa\omega=0
\]
have rank three.  Conditional on the TIR definitions, the subsystem therefore
has one continuous degree of freedom and is globally parametrized by $f$:
\[
\mathbf q(f)=
\left(
\frac{\ln2}{24\pi},
2\pi f,
f,
\frac{\ln2}{12}f
\right).
\]
This is a parameter-counting theorem inside the declared subsystem; it does not
supply an empirical value of $f$.

\section{A necessary negative result: self-duality is not extremality}

The half formalism also supplies an important boundary condition for future TIR
uses of self-duality.  In the finite DHSE-001 Stage-M M\"obius universe, the
forcing-count function obeys reciprocal symmetry
\[
N_n(q)=N_n(1/q),
\]
yet the self-dual point $q=1$ is the unique global maximum only for declared word
lengths $n=2,3$.  At $n=1$ and $n=4$, the exact maxima split into reciprocal
off-centre components.

Therefore
\begin{equation}
\boxed{
\text{self-duality or reciprocal symmetry}
\;\not\Rightarrow\;
\text{dynamical extremum at the fixed point}
}
\label{eq:self-duality-not-extremality}
\end{equation}
without an additional condition such as positivity, convexity, monotonicity,
variational closure, or another explicit theorem.

This negative theorem is relevant to TIR whenever a future construction attempts
to infer physical preference, stability, or attractor status from a self-dual
coordinate alone.

\section{Claim boundary}

The cross-relation established in this appendix does not prove any of the
following:
\begin{enumerate}
\item that standard quantum mechanics independently derives
      $\kappa=\ln2/(24\pi)$;
\item that $\Gamma_{\mathcal I}$ is a measured physical surface-refresh rate;
\item that reciprocal self-duality alone selects a dynamical attractor;
\item that the current Secret-of-a-Half construction closes a canonical bridge
      from non-trivial zeta zeros to the critical line;
\item that the phase-rate identity by itself validates any particle-sector mass
      assignment.
\end{enumerate}

The established interface is narrower and stronger: the exact binary midpoint
produces $\ln2$; TIR uses that quantity in a declared structural normalization;
and the resulting phase-rate and constraint-manifold statements then follow
exactly conditional on those definitions.

% !TEX root = ../metatime_monograph.tex
\chapter{Information-Spinor Crosswalk and Dependency Direction}
\label{app:information_spinor_crosswalk}

% CITATION_CONTEXT_V10_9_BEGIN
\paragraph{Established context.}
Binary information, projective quantum-state geometry, and geometric holonomy
follow the standard Shannon, Simon, Berry, and differential-geometric context
\cite{shannon1948,Simon1983,berry1984,BottTu1982}.  The identity-axis,
cycle-count, sector pairing, and CIEL/PhaseNav dependency assignments below are
specific to the TIR/Metatime programme.
% CITATION_CONTEXT_V10_9_END

\claimstatus{Mixed: exact mathematical identities plus explicitly marked Metatime/TIR model assignments.}

\section{Why this crosswalk exists}
The dedicated critical-axis programme under
\path{TIR/zeta_information_axis/}
formalizes the repeated appearance of the balanced coordinate $1/2$ in binary information, complement symmetry, Bloch geometry, Berry holonomy, destructive interference, centred reciprocal geometry, and the anti-linear zeta involution.  The parent Metatime monograph uses some of the same structures in a broader phenomenological programme.  This appendix records the direction of dependence so that a result from one layer is not silently promoted in the other.

The governing rule is:
\[
\boxed{
\text{exact mathematical bridge}
\;\neq\;
\text{physical identification}
\;\neq\;
\text{empirical validation}.
}
\]

\section{The half axis as an identity coordinate}
The binary complement
\[
C(\sigma)=1-\sigma
\]
has the unique fixed point
\[
\sigma=\frac12.
\]
Using the centered coordinate
\[
x=\sigma-\frac12,
\]
complement becomes the orientation reversal $x\mapsto-x$.  Hence negative, zero, and positive centered values can be read geometrically as the two orientations and the fixed axis of the same one-dimensional chart.  The same centering idea is used by the critical-axis submonograph for $u=s-1/2$.

In the qubit representation,
\[
n_z=2\sigma-1,
\]
so $\sigma=1/2$ is the Bloch equator.  Phase evolution can remain nontrivial while the identity coordinate remains balanced.  This is the sense in which the half value is an \emph{axis} rather than a frozen dynamical state.

\section{Normalized cycles before radians}
The critical-axis solver introduces a turn coordinate
\[
q\in\mathbb R/\mathbb Z
\]
with intrinsic winding frequency
\[
f_q=\frac{\Delta N}{\Delta t}.
\]
An angular representation is introduced only afterwards by
\[
\phi=Cq.
\]
For radians, $C=2\pi$.  The sequence
\[
\text{recurrence}\to\text{winding}\to\text{angle}
\]
is therefore conceptually prior to the conventional equation $\omega=2\pi f$.

This cross-references the phase and Poincar\'e structures of Chapter~\ref{ch:metatime_framework} and Appendix~\ref{app:poincare}, while retaining the publication claim boundary of Appendix~\ref{app:kappa}.

\section{The $24$ cross-factorization}
The exact integer identity is
\begin{equation}
24=8\cdot3=12\cdot2=6\cdot4.
\label{eq:appP-24}
\end{equation}
The equality is arithmetic.  Interpretations of these factors are model-level:
\begin{itemize}
\item $8\cdot3$: eight declared mixing sectors times three flavours;
\item $12\cdot2$: twelve projective cycles times two half-turn units;
\item $6\cdot4$: six spinor cycles times four half-turn units.
\end{itemize}
The three-flavour layer connects structurally to Chapter~\ref{ch:pmns_mixing}.  The projective/spinor interpretation connects to the Berry and holonomy structures reviewed in Chapter~\ref{ch:metatime_framework}.  None of these labels changes the status of $24$ from a discrete model normalization into an independently established physical constant.

\section{Arithmetic--geometry factorization of $\kappa$}
If the Metatime model assigns one binary information quantum to twelve normalized projective cycles,
\begin{equation}
\frac{\dd I}{\dd q}=\frac{\ln2}{12},
\label{eq:appP-info-turn}
\end{equation}
and if the chosen angular representation has closure period $C$, then
\begin{equation}
\frac{\dd I}{\dd\phi}=\frac{\ln2}{12C}.
\end{equation}
For $C=2\pi$ this reproduces
\begin{equation}
\kappa=\frac{\ln2}{24\pi}.
\label{eq:appP-kappa}
\end{equation}
Equation~\eqref{eq:appP-kappa} is exact conditional arithmetic once the cycle assignment and angular closure have been declared.  The independent physical derivation of the cycle assignment remains a model debt, consistently with Appendix~\ref{app:kappa}.

\section{Paired structures and implication directions}
\begin{longtable}{p{0.26\textwidth}p{0.29\textwidth}p{0.33\textwidth}}
\toprule
Source structure & Paired consequence & Status / direction \\
\midrule
$\sigma=1/2$ & $C(\sigma)=\sigma$ & Exact both-way characterization in the binary interval.\\
$\sigma=1/2$ & $H_{\rm bin}=\ln2$ & Exact; entropy maximum uniquely selects the half point.\\
$\sigma=1/2$ + qubit representation & Bloch equator $n_z=0$ & Exact inside the stated two-level representation.\\
Bloch equator + one azimuthal loop & Berry holonomy $-1$ & Standard geometric-phase result in the chosen state family.\\
Symmetric complementary readout + half-turn phase & exact destructive cancellation & Exact; the converse also selects $\sigma=1/2$.\\
$u=s-1/2$ + $J(s)=1-\bar s$ & fixed axis $\Re s=1/2$ & Exact/classical.\\
Centered zeta chart + $u\mapsto1/u$ & critical-axis set invariance & Exact; does not imply zero-set invariance.\\
Projective cycle + spin $1/2$ & double-cover period & Standard representation theory.\\
$\ln2/12$ per projective turn + $C=2\pi$ & $\kappa=\ln2/(24\pi)$ & Conditional arithmetic; cycle assignment is model-level.\\
TIR sector-holonomy operator & mass-spectrum candidate & Retrospective model layer; see Appendix~\ref{app:sector_holonomy_release}.\\
Critical-axis cancellation representation of every zeta zero & RH & Open bridge; not derived by Metatime holonomy alone.\\
\bottomrule
\end{longtable}

\section{Holonomy policy}
A relation edge is not automatically a holonomy edge.  Holonomy is meaningful only when a connection and a closed or composable path are defined.  This distinction matters for the parent programme because ``paired'' structures may be algebraic, geometric, semantic, or empirical without carrying a phase.

The CIEL/PhaseNav runtime therefore stores a normalized holonomy coordinate only on eligible path edges.  Generated PhaseNav geometry remains a candidate until independent TIR provenance supplies the semantic relation.  Conversely, a TIR relation without an executable geometric path does not acquire a synthetic phase merely to complete a diagram.

\section{Runtime dependency direction}
The companion CIEL/PhaseNav package is organized so that higher-level operators depend on primitive geometry rather than the reverse:
\[
\boxed{
\begin{array}{c}
\text{identity axis / centered coordinate}\\
\text{normalized cycle }\mathbb R/\mathbb Z\\
\text{phase crystal }(\mathbb R/\mathbb Z)^{36}\\
\text{typed relation edges}
\end{array}
\longrightarrow
\text{tangent/log-map vectorization}
\longrightarrow
\text{semantic and holonomic solvers}.
}
\]
On the phase crystal, the shortest oriented tangent displacement naturally lies in $[-1/2,1/2)$ per normalized coordinate.  Its sign is therefore generated by orientation; an identical coordinate produces zero.  This is the concrete computational meaning of the statement that signed vector components emerge from the geometry rather than being attached afterwards as semantic labels.

\section{Claim firewall}
The dedicated solver uses four statuses: \texttt{EXACT}, \texttt{STANDARD}, \texttt{MODEL}, and \texttt{OPEN}.  Exact closure permits only the first two classes.  Model closure may additionally use model assignments.  Open rules are never automatically promoted.

In particular, the solver can reproduce the mutually consistent half-axis routes and the conditional normalization \eqref{eq:appP-kappa}, while still reporting the canonical zeta zero-state representation as a missing premise.  This preserves the scientific-status policy of this publication candidate: technical consistency and structural unification do not by themselves constitute physical or number-theoretic closure.


% Expanded bibliography and source map for the Metatime monograph.
% Version v10.8, 2026-07-29.
% The grouping distinguishes external scientific context from original TIR claims.

\renewcommand{\bibname}{Bibliography and Source Map}
\begin{thebibliography}{999}
\small
\sloppy

\item[]\textit{Editorial note.} The references below document the established
mathematical, physical, experimental, and methodological context used throughout
this monograph. Inclusion of a source does not imply that the source endorses the
Metatime/TIR construction, nor that the original claims of this monograph follow
from the cited literature. Numerical particle data should be read against the
current Particle Data Group listings and their published updates.

\item[]\textbf{A. Reference data, precision measurements, and major experiments}

\bibitem{PDG2024}
S.~Navas \emph{et al.} (Particle Data Group),
``Review of Particle Physics,''
Phys.\ Rev.\ D \textbf{110}, 030001 (2024), with the 2025 online update.

\bibitem{CODATA2022}
E.~Tiesinga, P.~J.~Mohr, D.~B.~Newell, and B.~N.~Taylor,
``The 2022 CODATA Recommended Values of the Fundamental Physical Constants,''
NIST Standard Reference Database 121, Web Version 9.0 (2024),
\url{https://physics.nist.gov/constants}.

\bibitem{Planck2018}
N.~Aghanim \emph{et al.} (Planck Collaboration),
``Planck 2018 results. VI. Cosmological parameters,''
Astron.\ Astrophys.\ \textbf{641}, A6 (2020).

\bibitem{DESI2024}
A.~G.~Adame \emph{et al.} (DESI Collaboration),
``DESI 2024 VI: Cosmological Constraints from the Measurements of Baryon
Acoustic Oscillations,'' arXiv:2404.03002 (2024).

\bibitem{ATLAS2012Higgs}
G.~Aad \emph{et al.} (ATLAS Collaboration),
``Observation of a new particle in the search for the Standard Model Higgs
boson with the ATLAS detector at the LHC,''
Phys.\ Lett.\ B \textbf{716}, 1--29 (2012).

\bibitem{CMS2012Higgs}
S.~Chatrchyan \emph{et al.} (CMS Collaboration),
``Observation of a new boson at a mass of 125 GeV with the CMS experiment at
the LHC,'' Phys.\ Lett.\ B \textbf{716}, 30--61 (2012).

\bibitem{ATLAS2022Charm}
G.~Aad \emph{et al.} (ATLAS Collaboration),
``Direct constraint on the Higgs--charm coupling from a search for Higgs boson
decays into charm quarks with the ATLAS detector,''
Eur.\ Phys.\ J.\ C \textbf{82}, 717 (2022).

\bibitem{SuperK1998}
Y.~Fukuda \emph{et al.} (Super-Kamiokande Collaboration),
``Evidence for oscillation of atmospheric neutrinos,''
Phys.\ Rev.\ Lett.\ \textbf{81}, 1562--1567 (1998).

\bibitem{SNO2002}
Q.~R.~Ahmad \emph{et al.} (SNO Collaboration),
``Direct evidence for neutrino flavor transformation from neutral-current
interactions in the Sudbury Neutrino Observatory,''
Phys.\ Rev.\ Lett.\ \textbf{89}, 011301 (2002).

\bibitem{KamLAND2003}
K.~Eguchi \emph{et al.} (KamLAND Collaboration),
``First results from KamLAND: Evidence for reactor antineutrino disappearance,''
Phys.\ Rev.\ Lett.\ \textbf{90}, 021802 (2003).

\bibitem{DayaBay2012}
F.~P.~An \emph{et al.} (Daya Bay Collaboration),
``Observation of electron-antineutrino disappearance at Daya Bay,''
Phys.\ Rev.\ Lett.\ \textbf{108}, 171803 (2012).

\bibitem{Abel2020nEDM}
C.~Abel \emph{et al.},
``Measurement of the permanent electric dipole moment of the neutron,''
Phys.\ Rev.\ Lett.\ \textbf{124}, 081803 (2020).

\bibitem{MuonG2Final2025}
D.~P.~Aguillard \emph{et al.} (Muon $g-2$ Collaboration),
``Measurement of the positive muon anomalous magnetic moment to 127 ppb,''
Phys.\ Rev.\ Lett.\ \textbf{135} (2025), DOI: 10.1103/7clf-sm2v.

\bibitem{Riess1998}
A.~G.~Riess \emph{et al.},
``Observational evidence from supernovae for an accelerating universe and a
cosmological constant,'' Astron.\ J.\ \textbf{116}, 1009--1038 (1998).

\bibitem{Perlmutter1999}
S.~Perlmutter \emph{et al.},
``Measurements of $\Omega$ and $\Lambda$ from 42 high-redshift supernovae,''
Astrophys.\ J.\ \textbf{517}, 565--586 (1999).

\item[]\textbf{B. Quantum mechanics, quantum field theory, and the Standard Model}

\bibitem{Noether1918}
E.~Noether,
``Invariante Variationsprobleme,''
Nachr.\ Ges.\ Wiss.\ G\"ottingen, Math.-Phys.\ Kl. 235--257 (1918).

\bibitem{Dirac1928}
P.~A.~M.~Dirac,
``The quantum theory of the electron,''
Proc.\ R.\ Soc.\ Lond.\ A \textbf{117}, 610--624 (1928); \textbf{118},
351--361 (1928).

\bibitem{Weyl1929}
H.~Weyl,
``Elektron und Gravitation. I,''
Z.\ Phys.\ \textbf{56}, 330--352 (1929).

\bibitem{Tomonaga1946}
S.~Tomonaga,
``On a relativistically invariant formulation of the quantum theory of wave
fields,'' Prog.\ Theor.\ Phys.\ \textbf{1}, 27--42 (1946).

\bibitem{Schwinger1948}
J.~Schwinger,
``Quantum electrodynamics. I. A covariant formulation,''
Phys.\ Rev.\ \textbf{74}, 1439--1461 (1948).

\bibitem{Feynman1949}
R.~P.~Feynman,
``Space-time approach to quantum electrodynamics,''
Phys.\ Rev.\ \textbf{76}, 769--789 (1949).

\bibitem{Dyson1949}
F.~J.~Dyson,
``The radiation theories of Tomonaga, Schwinger, and Feynman,''
Phys.\ Rev.\ \textbf{75}, 486--502 (1949).

\bibitem{YangMills1954}
C.~N.~Yang and R.~L.~Mills,
``Conservation of isotopic spin and isotopic gauge invariance,''
Phys.\ Rev.\ \textbf{96}, 191--195 (1954).

\bibitem{Glashow1961}
S.~L.~Glashow,
``Partial-symmetries of weak interactions,''
Nucl.\ Phys.\ \textbf{22}, 579--588 (1961).

\bibitem{Weinberg1967}
S.~Weinberg,
``A model of leptons,'' Phys.\ Rev.\ Lett.\ \textbf{19}, 1264--1266 (1967).

\bibitem{Salam1968}
A.~Salam,
``Weak and electromagnetic interactions,'' in
\emph{Elementary Particle Theory}, edited by N.~Svartholm
(Almqvist and Wiksell, Stockholm, 1968), pp.~367--377.

\bibitem{EnglertBrout1964}
F.~Englert and R.~Brout,
``Broken symmetry and the mass of gauge vector mesons,''
Phys.\ Rev.\ Lett.\ \textbf{13}, 321--323 (1964).

\bibitem{Higgs1964}
P.~W.~Higgs,
``Broken symmetries and the masses of gauge bosons,''
Phys.\ Rev.\ Lett.\ \textbf{13}, 508--509 (1964).

\bibitem{GHK1964}
G.~S.~Guralnik, C.~R.~Hagen, and T.~W.~B.~Kibble,
``Global conservation laws and massless particles,''
Phys.\ Rev.\ Lett.\ \textbf{13}, 585--587 (1964).

\bibitem{tHooftVeltman1972}
G.~'t~Hooft and M.~Veltman,
``Regularization and renormalization of gauge fields,''
Nucl.\ Phys.\ B \textbf{44}, 189--213 (1972).

\bibitem{GrossWilczek1973}
D.~J.~Gross and F.~Wilczek,
``Ultraviolet behavior of non-Abelian gauge theories,''
Phys.\ Rev.\ Lett.\ \textbf{30}, 1343--1346 (1973).

\bibitem{Politzer1973}
H.~D.~Politzer,
``Reliable perturbative results for strong interactions?''
Phys.\ Rev.\ Lett.\ \textbf{30}, 1346--1349 (1973).

\bibitem{Wilson1974}
K.~G.~Wilson,
``Confinement of quarks,'' Phys.\ Rev.\ D \textbf{10}, 2445--2459 (1974).

\bibitem{PeskinSchroeder1995}
M.~E.~Peskin and D.~V.~Schroeder,
\emph{An Introduction to Quantum Field Theory}
(Addison-Wesley, Reading, 1995).

\bibitem{WeinbergQFT1995}
S.~Weinberg,
\emph{The Quantum Theory of Fields, Vols. I--III}
(Cambridge University Press, Cambridge, 1995--2000).

\item[]\textbf{C. Flavor, hadrons, neutrinos, and symmetry classification}

\bibitem{gellmann1961}
M.~Gell-Mann,
``The eightfold way: A theory of strong interaction symmetry,''
Caltech Synchrotron Laboratory Report CTSL-20 (1961).

\bibitem{Neeman1961}
Y.~Ne'eman,
``Derivation of strong interactions from a gauge invariance,''
Nucl.\ Phys.\ \textbf{26}, 222--229 (1961).

\bibitem{GellMann1962}
M.~Gell-Mann,
``Symmetries of baryons and mesons,''
Phys.\ Rev.\ \textbf{125}, 1067--1084 (1962).

\bibitem{okubo1962}
S.~Okubo,
``Note on unitary symmetry in strong interactions,''
Prog.\ Theor.\ Phys.\ \textbf{27}, 949--966 (1962).

\bibitem{Cabibbo1963}
N.~Cabibbo,
``Unitary symmetry and leptonic decays,''
Phys.\ Rev.\ Lett.\ \textbf{10}, 531--533 (1963).

\bibitem{GellMann1964}
M.~Gell-Mann,
``A schematic model of baryons and mesons,''
Phys.\ Lett.\ \textbf{8}, 214--215 (1964).

\bibitem{Zweig1964}
G.~Zweig,
``An SU(3) model for strong interaction symmetry and its breaking,''
CERN Reports TH-401 and TH-412 (1964).

\bibitem{kobayashi1973}
M.~Kobayashi and T.~Maskawa,
``CP-violation in the renormalizable theory of weak interaction,''
Prog.\ Theor.\ Phys.\ \textbf{49}, 652--657 (1973).

\bibitem{Wolfenstein1983}
L.~Wolfenstein,
``Parametrization of the Kobayashi--Maskawa matrix,''
Phys.\ Rev.\ Lett.\ \textbf{51}, 1945--1947 (1983).

\bibitem{Jarlskog1985}
C.~Jarlskog,
``Commutator of the quark mass matrices in the Standard Electroweak Model and
a measure of maximal CP nonconservation,''
Phys.\ Rev.\ Lett.\ \textbf{55}, 1039--1042 (1985).

\bibitem{Maki1962}
Z.~Maki, M.~Nakagawa, and S.~Sakata,
``Remarks on the unified model of elementary particles,''
Prog.\ Theor.\ Phys.\ \textbf{28}, 870--880 (1962).

\bibitem{Pontecorvo1957}
B.~Pontecorvo,
``Mesonium and antimesonium,''
Zh.\ Eksp.\ Teor.\ Fiz.\ \textbf{33}, 549--551 (1957)
[Sov.\ Phys.\ JETP \textbf{6}, 429 (1958)].

\bibitem{Wolfenstein1978}
L.~Wolfenstein,
``Neutrino oscillations in matter,''
Phys.\ Rev.\ D \textbf{17}, 2369--2374 (1978).

\bibitem{MikheyevSmirnov1985}
S.~P.~Mikheyev and A.~Y.~Smirnov,
``Resonance amplification of oscillations in matter and spectroscopy of solar
neutrinos,'' Sov.\ J.\ Nucl.\ Phys.\ \textbf{42}, 913--917 (1985).

\item[]\textbf{D. Geometric phase, differential geometry, topology, and holonomy}

\bibitem{Pancharatnam1956}
S.~Pancharatnam,
``Generalized theory of interference and its applications,''
Proc.\ Indian Acad.\ Sci.\ A \textbf{44}, 247--262 (1956).

\bibitem{Simon1983}
B.~Simon,
``Holonomy, the quantum adiabatic theorem, and Berry's phase,''
Phys.\ Rev.\ Lett.\ \textbf{51}, 2167--2170 (1983).

\bibitem{berry1984}
M.~V.~Berry,
``Quantal phase factors accompanying adiabatic changes,''
Proc.\ R.\ Soc.\ Lond.\ A \textbf{392}, 45--57 (1984).

\bibitem{WilczekZee1984}
F.~Wilczek and A.~Zee,
``Appearance of gauge structure in simple dynamical systems,''
Phys.\ Rev.\ Lett.\ \textbf{52}, 2111--2114 (1984).

\bibitem{AharonovAnandan1987}
Y.~Aharonov and J.~Anandan,
``Phase change during a cyclic quantum evolution,''
Phys.\ Rev.\ Lett.\ \textbf{58}, 1593--1596 (1987).

\bibitem{AnandanAharonov1990}
J.~Anandan and Y.~Aharonov,
``Geometry of quantum evolution,''
Phys.\ Rev.\ Lett.\ \textbf{65}, 1697--1700 (1990).

\bibitem{ProvostVallee1980}
J.~P.~Provost and G.~Vall\'ee,
``Riemannian structure on manifolds of quantum states,''
Commun.\ Math.\ Phys.\ \textbf{76}, 289--301 (1980).

\bibitem{Kahler1933}
E.~K\"ahler,
``\"Uber eine bemerkenswerte Hermitesche Metrik,''
Abh.\ Math.\ Sem.\ Univ.\ Hamburg \textbf{9}, 173--186 (1933).

\bibitem{Hopf1931}
H.~Hopf,
``\"Uber die Abbildungen der dreidimensionalen Sph\"are auf die Kugelfl\"ache,''
Math.\ Ann.\ \textbf{104}, 637--665 (1931).

\bibitem{Chern1946}
S.-S.~Chern,
``Characteristic classes of Hermitian manifolds,''
Ann.\ Math.\ \textbf{47}, 85--121 (1946).

\bibitem{AmbroseSinger1953}
W.~Ambrose and I.~M.~Singer,
``A theorem on holonomy,''
Trans.\ Am.\ Math.\ Soc.\ \textbf{75}, 428--443 (1953).

\bibitem{pontryagin1955}
L.~S.~Pontryagin,
``Smooth manifolds and their applications in homotopy theory,''
Trudy Mat.\ Inst.\ Steklov \textbf{45}, 1--139 (1955).

\bibitem{KobayashiNomizu1963}
S.~Kobayashi and K.~Nomizu,
\emph{Foundations of Differential Geometry}, Vols.~I--II
(Wiley-Interscience, New York, 1963--1969).

\bibitem{BottTu1982}
R.~Bott and L.~W.~Tu,
\emph{Differential Forms in Algebraic Topology}
(Springer, New York, 1982).

\bibitem{Nakahara2003}
M.~Nakahara,
\emph{Geometry, Topology and Physics}, 2nd ed.
(Institute of Physics Publishing, Bristol, 2003).

\bibitem{BengtssonZyczkowski2006}
I.~Bengtsson and K.~\.{Z}yczkowski,
\emph{Geometry of Quantum States}
(Cambridge University Press, Cambridge, 2006).

\item[]\textbf{E. Information theory, statistical geometry, and computation}

\bibitem{shannon1948}
C.~E.~Shannon,
``A mathematical theory of communication,''
Bell Syst.\ Tech.\ J.\ \textbf{27}, 379--423 and 623--656 (1948).

\bibitem{vonNeumann1932}
J.~von Neumann,
\emph{Mathematische Grundlagen der Quantenmechanik}
(Springer, Berlin, 1932).

\bibitem{Fisher1925}
R.~A.~Fisher,
``Theory of statistical estimation,''
Proc.\ Camb.\ Philos.\ Soc.\ \textbf{22}, 700--725 (1925).

\bibitem{Rao1945}
C.~R.~Rao,
``Information and the accuracy attainable in the estimation of statistical
parameters,'' Bull.\ Calcutta Math.\ Soc.\ \textbf{37}, 81--91 (1945).

\bibitem{Jaynes1957}
E.~T.~Jaynes,
``Information theory and statistical mechanics,''
Phys.\ Rev.\ \textbf{106}, 620--630 (1957); \textbf{108}, 171--190 (1957).

\bibitem{Landauer1961}
R.~Landauer,
``Irreversibility and heat generation in the computing process,''
IBM J.\ Res.\ Dev.\ \textbf{5}, 183--191 (1961).

\bibitem{Bures1969}
D.~Bures,
``An extension of Kakutani's theorem on infinite product measures to the
tensor product of semifinite $W^*$-algebras,''
Trans.\ Am.\ Math.\ Soc.\ \textbf{135}, 199--212 (1969).

\bibitem{Wootters1981}
W.~K.~Wootters,
``Statistical distance and Hilbert space,''
Phys.\ Rev.\ D \textbf{23}, 357--362 (1981).

\bibitem{Bennett1982}
C.~H.~Bennett,
``The thermodynamics of computation---a review,''
Int.\ J.\ Theor.\ Phys.\ \textbf{21}, 905--940 (1982).

\bibitem{BraunsteinCaves1994}
S.~L.~Braunstein and C.~M.~Caves,
``Statistical distance and the geometry of quantum states,''
Phys.\ Rev.\ Lett.\ \textbf{72}, 3439--3443 (1994).

\bibitem{AmariNagaoka2000}
S.-I.~Amari and H.~Nagaoka,
\emph{Methods of Information Geometry}
(American Mathematical Society, Providence, 2000).

\bibitem{NielsenChuang2000}
M.~A.~Nielsen and I.~L.~Chuang,
\emph{Quantum Computation and Quantum Information}
(Cambridge University Press, Cambridge, 2000).

\bibitem{Lloyd2000}
S.~Lloyd,
``Ultimate physical limits to computation,''
Nature \textbf{406}, 1047--1054 (2000).

\bibitem{CoverThomas2006}
T.~M.~Cover and J.~A.~Thomas,
\emph{Elements of Information Theory}, 2nd ed.
(Wiley, Hoboken, 2006).

\item[]\textbf{F. Synchronization, nonlinear dynamics, and collective phase models}

\bibitem{Kuramoto1975}
Y.~Kuramoto,
``Self-entrainment of a population of coupled nonlinear oscillators,'' in
\emph{International Symposium on Mathematical Problems in Theoretical Physics},
Lecture Notes in Physics Vol.~39 (Springer, Berlin, 1975), pp.~420--422.

\bibitem{Feigenbaum1978}
M.~J.~Feigenbaum,
``Quantitative universality for a class of nonlinear transformations,''
J.\ Stat.\ Phys.\ \textbf{19}, 25--52 (1978).

\bibitem{Strogatz2000}
S.~H.~Strogatz,
``From Kuramoto to Crawford: Exploring the onset of synchronization in
populations of coupled oscillators,''
Physica D \textbf{143}, 1--20 (2000).

\bibitem{Acebron2005}
J.~A.~Acebr\'on, L.~L.~Bonilla, C.~J.~P.~Vicente, F.~Ritort, and R.~Spigler,
``The Kuramoto model: A simple paradigm for synchronization phenomena,''
Rev.\ Mod.\ Phys.\ \textbf{77}, 137--185 (2005).

\bibitem{OttAntonsen2008}
E.~Ott and T.~M.~Antonsen,
``Low dimensional behavior of large systems of globally coupled oscillators,''
Chaos \textbf{18}, 037113 (2008).

\bibitem{Strogatz2015}
S.~H.~Strogatz,
\emph{Nonlinear Dynamics and Chaos}, 2nd ed.
(Westview Press, Boulder, 2015).

\item[]\textbf{G. Collatz dynamics, Ramanujan analysis, primes, and zeta structure}

\bibitem{Terras1976}
R.~Terras,
``A stopping time problem on the positive integers,''
Acta Arith.\ \textbf{30}, 241--252 (1976).

\bibitem{Lagarias1985}
J.~C.~Lagarias,
``The $3x+1$ problem and its generalizations,''
Am.\ Math.\ Mon.\ \textbf{92}, 3--23 (1985).

\bibitem{Wirsching1998}
G.~J.~Wirsching,
\emph{The Dynamical System Generated by the $3n+1$ Function}
(Springer, Berlin, 1998).

\bibitem{Tao2022Collatz}
T.~Tao,
``Almost all orbits of the Collatz map attain almost bounded values,''
Forum Math.\ Pi \textbf{10}, e12 (2022).

\bibitem{Ramanujan1918}
S.~Ramanujan,
``On certain trigonometrical sums and their applications in the theory of
numbers,'' Trans.\ Cambridge Philos.\ Soc.\ \textbf{22}, 259--276 (1918).

\bibitem{HardyRamanujan1918}
G.~H.~Hardy and S.~Ramanujan,
``Asymptotic formulae in combinatory analysis,''
Proc.\ Lond.\ Math.\ Soc.\ \textbf{17}, 75--115 (1918).

\bibitem{Rademacher1937}
H.~Rademacher,
``On the partition function $p(n)$,''
Proc.\ Natl.\ Acad.\ Sci.\ USA \textbf{23}, 78--84 (1937).

\bibitem{HardyLittlewood1923}
G.~H.~Hardy and J.~E.~Littlewood,
``Some problems of `Partitio Numerorum'; III: On the expression of a number as
a sum of primes,'' Acta Math.\ \textbf{44}, 1--70 (1923).

\bibitem{Riemann1859}
B.~Riemann,
``\"Uber die Anzahl der Primzahlen unter einer gegebenen Gr\"osse,''
Monatsber.\ K\"onigl.\ Preuss.\ Akad.\ Wiss.\ Berlin 671--680 (1859).

\bibitem{Selberg1956}
A.~Selberg,
``Harmonic analysis and discontinuous groups in weakly symmetric Riemannian
spaces with applications to Dirichlet series,''
J.\ Indian Math.\ Soc.\ \textbf{20}, 47--87 (1956).

\bibitem{BerryKeating1999}
M.~V.~Berry and J.~P.~Keating,
``The Riemann zeros and eigenvalue asymptotics,''
SIAM Rev.\ \textbf{41}, 236--266 (1999).

\bibitem{Connes1999}
A.~Connes,
``Trace formula in noncommutative geometry and the zeros of the Riemann zeta
function,'' Selecta Math.\ \textbf{5}, 29--106 (1999).

\bibitem{Zhang2014}
Y.~Zhang,
``Bounded gaps between primes,''
Ann.\ Math.\ \textbf{179}, 1121--1174 (2014).

\bibitem{Polymath2014}
D.~H.~J.~Polymath,
``New equidistribution estimates of Zhang type,''
Algebra Number Theory \textbf{8}, 2067--2199 (2014).

\bibitem{Maynard2015}
J.~Maynard,
``Small gaps between primes,''
Ann.\ Math.\ \textbf{181}, 383--413 (2015).

\bibitem{Apostol1976}
T.~M.~Apostol,
\emph{Introduction to Analytic Number Theory}
(Springer, New York, 1976).

\bibitem{HardyWright2008}
G.~H.~Hardy and E.~M.~Wright,
\emph{An Introduction to the Theory of Numbers}, 6th ed., revised by
D.~R.~Heath-Brown and J.~H.~Silverman
(Oxford University Press, Oxford, 2008).

\item[]\textbf{H. Anomalies, CP violation, strong CP, and electric dipole moments}

\bibitem{Adler1969}
S.~L.~Adler,
``Axial-vector vertex in spinor electrodynamics,''
Phys.\ Rev.\ \textbf{177}, 2426--2438 (1969).

\bibitem{BellJackiw1969}
J.~S.~Bell and R.~Jackiw,
``A PCAC puzzle: $\pi^0\to\gamma\gamma$ in the sigma model,''
Nuovo Cim.\ A \textbf{60}, 47--61 (1969).

\bibitem{tHooft1976}
G.~'t~Hooft,
``Symmetry breaking through Bell--Jackiw anomalies,''
Phys.\ Rev.\ Lett.\ \textbf{37}, 8--11 (1976).

\bibitem{PecceiQuinn1977}
R.~D.~Peccei and H.~R.~Quinn,
``CP conservation in the presence of pseudoparticles,''
Phys.\ Rev.\ Lett.\ \textbf{38}, 1440--1443 (1977).

\bibitem{Weinberg1978Axion}
S.~Weinberg,
``A new light boson?'' Phys.\ Rev.\ Lett.\ \textbf{40}, 223--226 (1978).

\bibitem{Wilczek1978Axion}
F.~Wilczek,
``Problem of strong $P$ and $T$ invariance in the presence of instantons,''
Phys.\ Rev.\ Lett.\ \textbf{40}, 279--282 (1978).

\bibitem{crewther1979}
R.~J.~Crewther, P.~Di~Vecchia, G.~Veneziano, and E.~Witten,
``Chiral estimate of the electric dipole moment of the neutron in quantum
chromodynamics,'' Phys.\ Lett.\ B \textbf{88}, 123--127 (1979).

\bibitem{witten1982}
E.~Witten,
``An SU(2) anomaly,'' Phys.\ Lett.\ B \textbf{117}, 324--328 (1982).

\bibitem{pospelov2005}
M.~Pospelov and A.~Ritz,
``Electric dipole moments as probes of new physics,''
Ann.\ Phys.\ (N.Y.) \textbf{318}, 119--169 (2005).

\item[]\textbf{I. Gravitation, cosmology, horizons, and information}

\bibitem{Einstein1915}
A.~Einstein,
``Die Feldgleichungen der Gravitation,''
Sitzungsber.\ Preuss.\ Akad.\ Wiss.\ Berlin 844--847 (1915).

\bibitem{Friedmann1922}
A.~Friedmann,
``\"Uber die Kr\"ummung des Raumes,''
Z.\ Phys.\ \textbf{10}, 377--386 (1922).

\bibitem{Lemaitre1927}
G.~Lema\^itre,
``Un univers homog\`ene de masse constante et de rayon croissant, rendant
compte de la vitesse radiale des n\'ebuleuses extra-galactiques,''
Ann.\ Soc.\ Sci.\ Bruxelles A \textbf{47}, 49--59 (1927).

\bibitem{Hubble1929}
E.~Hubble,
``A relation between distance and radial velocity among extra-galactic
nebulae,'' Proc.\ Natl.\ Acad.\ Sci.\ USA \textbf{15}, 168--173 (1929).

\bibitem{Penrose1965}
R.~Penrose,
``Gravitational collapse and space-time singularities,''
Phys.\ Rev.\ Lett.\ \textbf{14}, 57--59 (1965).

\bibitem{Bekenstein1973}
J.~D.~Bekenstein,
``Black holes and entropy,'' Phys.\ Rev.\ D \textbf{7}, 2333--2346 (1973).

\bibitem{Hawking1975}
S.~W.~Hawking,
``Particle creation by black holes,''
Commun.\ Math.\ Phys.\ \textbf{43}, 199--220 (1975).

\bibitem{HawkingEllis1973}
S.~W.~Hawking and G.~F.~R.~Ellis,
\emph{The Large Scale Structure of Space-Time}
(Cambridge University Press, Cambridge, 1973).

\bibitem{Wald1984}
R.~M.~Wald,
\emph{General Relativity}
(University of Chicago Press, Chicago, 1984).

\bibitem{Weinberg1989CC}
S.~Weinberg,
``The cosmological constant problem,''
Rev.\ Mod.\ Phys.\ \textbf{61}, 1--23 (1989).

\bibitem{Jacobson1995}
T.~Jacobson,
``Thermodynamics of spacetime: The Einstein equation of state,''
Phys.\ Rev.\ Lett.\ \textbf{75}, 1260--1263 (1995).

\bibitem{Carroll2001}
S.~M.~Carroll,
``The cosmological constant,''
Living Rev.\ Relativ.\ \textbf{4}, 1 (2001).

\bibitem{Verlinde2011}
E.~P.~Verlinde,
``On the origin of gravity and the laws of Newton,''
J.\ High Energy Phys.\ \textbf{04}, 029 (2011).

\item[]\textbf{J. Scientific methodology, falsifiability, preregistration, and reproducibility}

\bibitem{Popper1959}
K.~R.~Popper,
\emph{The Logic of Scientific Discovery}
(Hutchinson, London, 1959).

\bibitem{Lakatos1978}
I.~Lakatos,
\emph{The Methodology of Scientific Research Programmes}
(Cambridge University Press, Cambridge, 1978).

\bibitem{Ioannidis2005}
J.~P.~A.~Ioannidis,
``Why most published research findings are false,''
PLoS Med.\ \textbf{2}, e124 (2005).

\bibitem{Sandve2013}
G.~K.~Sandve, A.~Nekrutenko, J.~Taylor, and E.~Hovig,
``Ten simple rules for reproducible computational research,''
PLoS Comput.\ Biol.\ \textbf{9}, e1003285 (2013).

\bibitem{Wilson2014}
G.~Wilson \emph{et al.},
``Best practices for scientific computing,''
PLoS Biol.\ \textbf{12}, e1001745 (2014).

\bibitem{Munafo2017}
M.~R.~Munaf\`o \emph{et al.},
``A manifesto for reproducible science,''
Nat.\ Hum.\ Behav.\ \textbf{1}, 0021 (2017).

\bibitem{Nosek2018}
B.~A.~Nosek, C.~R.~Ebersole, A.~C.~DeHaven, and D.~T.~Mellor,
``The preregistration revolution,''
Proc.\ Natl.\ Acad.\ Sci.\ USA \textbf{115}, 2600--2606 (2018).

\bibitem{Stark2018}
P.~B.~Stark,
``Before reproducibility must come preproducibility,''
Nature \textbf{557}, 613 (2018).

\end{thebibliography}

\vfill
\noindent\small\textit{Source code and version history:}
\url{https://github.com/AdrianLipa90/The-Fundamental-Theory-of-Informational-Relations}
\par\vspace{2cm}
\backmatter
\printindex
\end{document}
